\pdfoutput=1
\documentclass[11pt]{amsart}
\usepackage{amsmath,amssymb,amsthm}
\usepackage[margin=1.2in]{geometry}
\usepackage[colorlinks=true,linkcolor=blue,citecolor=blue,urlcolor=blue]{hyperref}
\hypersetup{pdfauthor={Bernd Johannes Wuebben},
 pdftitle={Convergence of the minimizing flow on Kronecker quiver complexes},
 pdfsubject={MSC 2020: 16G20, 53D20; 14D20, 14L24, 15A45, 18G80, 47A63},
 pdfkeywords={quiver representation, complex of representations, stability condition, harmonic metric, polar decomposition, Kempf-Ness function, moment map, Schur multiplier, Lyapunov function, K\"ahler quotient}}
\setlength{\emergencystretch}{2em}
\newtheorem{theorem}{Theorem}[section]
\newtheorem{proposition}[theorem]{Proposition}
\newtheorem{lemma}[theorem]{Lemma}
\newtheorem{corollary}[theorem]{Corollary}
\theoremstyle{definition}
\newtheorem{definition}[theorem]{Definition}
\theoremstyle{remark}
\newtheorem{remark}[theorem]{Remark}
\newcommand{\CC}{\mathbb C}
\newcommand{\RR}{\mathbb R}
\newcommand{\id}{I}
\newcommand{\HS}{\mathrm{HS}}
\newcommand{\ii}{\mathrm i}
\newcommand{\coh}{\mathrm{coh}}
\newcommand{\cK}{\mathcal K}
\newcommand{\cD}{\mathcal D}
\newcommand{\cS}{\mathcal S}
\DeclareMathOperator{\tr}{tr}
\DeclareMathOperator{\Ein}{Ein}
\DeclareMathOperator{\Rea}{Re}
\DeclareMathOperator{\Ima}{Im}
\DeclareMathOperator{\diag}{diag}
\DeclareMathOperator{\spec}{spec}
\newcommand{\norm}[1]{\left\lVert#1\right\rVert}
\newcommand{\inner}[2]{\left\langle#1,#2\right\rangle}

\title[The minimizing flow on Kronecker quiver complexes]{Convergence of the
minimizing flow\\ on Kronecker quiver complexes}
\author{Bernd Johannes Wuebben}
\email{wuebben@gmail.com}
\subjclass[2020]{Primary 16G20, 53D20; Secondary 14D20, 14L24, 15A45, 18G80, 47A63}
\keywords{Quiver representation, complex of representations, stability
condition, harmonic metric, polar decomposition, Kempf--Ness function, moment
map, Schur multiplier, Lyapunov function, K\"ahler quotient}
\date{November 18, 2026\\This version: November 19, 2026}

\begin{document}
\begin{abstract}
For a two-term complex $(V\to W)\to(U\to0)$ of generalized Kronecker
quiver representations, we prove exponential cancellation of its acyclic
part under the polar-phase interpretation of the minimizing flow of
Haiden, Katzarkov, Kontsevich and Pandit. Suppose the differential
$D\colon V\to U$ is surjective, $U\ne0$, and the cohomology
$H=(\ker D\to W)$ is nonzero at both vertices. For all vertex charges in
the upper half-plane and every initial Hermitian metric, the endomorphism
defining the flow approaches that of $H$ exponentially in trace norm,
with the induced metric on $H$ and extension by zero. If the source
charge has larger argument than the target charge and $H$ is polystable, the induced metric
multiplied by $e^{t\arg Z(H)}$ converges exponentially to a harmonic
metric, where $Z(H)$ is the central charge of the cohomology. In this chamber semistability suffices for the shifted moment
map to tend to zero. The arrow-energy estimate uses the kernel of second divided differences
of $e^{-x}$ at the nodes $\ell_i,0,\ell_j$ with $\ell_i>0$. Its Schur
product $J$ with any positive semidefinite matrix satisfies
$\operatorname{tr}J_-\leq\frac13\operatorname{tr}J$, where $J_-$ is
the negative part. We also prove that, at fixed positive definite
imaginary part, minus the polar phase is strictly monotone for the trace
inner product.
\end{abstract}
\maketitle

\section{Introduction}\label{sec:intro}

Haiden, Katzarkov, Kontsevich and Pandit \cite{hkkp26} propose a notion of
K\"ahler geometry on a triangulated category with a stability condition in
the sense of Bridgeland \cite{bri07}. In their framework an object carries a
space of metrics $h$ and a complex-valued function $\cS_{\CC}$ on it with
differential $\delta h\mapsto\tr(\Omega\,h^{-1}\delta h)$, for an endomorphism
$\Omega$ whose trace is the central charge of the object. A \emph{phase}
$\arg\Omega$ is defined by a positivity condition, the \emph{mass} is
$m=\tr|\Omega|$, and the metrics move by the flow
$h^{-1}\dot h=-\arg\Omega$, called the minimizing flow in
\cite[\S4.1]{hkkp26}. In \cite[\S\S1,~4.1]{hkkp26} the mass is expected to
decrease along the flow, the flow to converge to a \emph{harmonic} metric,
one of constant phase, when the object admits one, and harmonic metrics to
exist on polystable objects, a categorical form of the theorem of Donaldson,
Uhlenbeck and Yau \cite{don85,uy86}. For a representation of a quiver,
$\Omega$ is the central charge of each vertex shifted by the moment map, and
the harmonic metrics are the zeros of the shifted moment map. For complexes
of quiver representations $\Omega$ is defined through a superconnection in a
graded algebra \cite[\S5.4]{hkkp26}, and it is asked there whether the mass
decreases \cite[Question~5.18]{hkkp26}.

For a complex the endomorphism $\Omega$ is in general not normal
\cite[Example~2.3]{wue26}, whereas in \cite[\S5.4]{hkkp26} its phase is
written with the functional calculus. Throughout this paper $\arg\Omega$ is
taken in the sense of the positivity condition of \cite[\S4.2]{hkkp26}, that
is, as the argument of the unitary factor in the polar decomposition
(Section~\ref{sec:setting}). Every statement about the flow below, including
the comparison with \cite[Question~5.18]{hkkp26}, refers to this reading. We
call the resulting flow the \emph{phase flow}: it is the minimizing flow of
\cite[\S4.1]{hkkp26} with this reading of the phase.

The simplest complexes on which both a differential and the arrows of a
quiver act are
\begin{equation}\label{eq:introcomplex}
 (V\xrightarrow{\,(A_j)\,}W)\ \xrightarrow{\,(D,0)\,}\ (U\to0),
\end{equation}
two-term complexes of representations of the generalized Kronecker quiver
with $k$ arrows, with $D:V\to U$ surjective and $A_1,\dots,A_k:V\to W$. The
cohomology is the representation $\ker D\to W$ in degree zero; it is a
subcomplex, and the quotient $(V/\ker D\to0)\to(U\to0)$ is an acyclic complex
of rank $r=\dim U$, the \emph{acyclic part}. So \eqref{eq:introcomplex} is
quasi-isomorphic to a representation. For these complexes $\Omega$ was
computed in \cite{wue26}, where it was shown that the mass can increase along
the phase flow \cite[Theorem~1.1]{wue26}; that the contribution of the
acyclic part to $\Omega$ nevertheless decays if the least singular
value of $D$ in the initial metrics is large \cite[Theorem~1.3]{wue26}; and
that for one pair of charges and $\dim\ker D=\dim W$ the moment map of the
cohomology representation then tends to zero if its Kempf--Ness function is
bounded below \cite[Theorem~1.4]{wue26}.

We prove that, for every initial metric, the contribution of the acyclic
part decays, for all charges, and that, when the charge at the source vertex
has the larger argument, the metrics induced on the cohomology converge,
after a scalar normalization, to a harmonic metric of the cohomology
representation whenever that representation is polystable. This
removes the assumption on the initial metrics of
\cite[Theorems~1.3,~1.4]{wue26}. The main results are
Theorems~\ref{thm:introA} and~\ref{thm:introB}.

\subsection{Results}
Let $h=(h_V,h_U,h_W)$ be Hermitian metrics on the three spaces of
\eqref{eq:introcomplex}; adjoints, written ${}^\dagger$, are taken for the
current metrics. Let $z_0,z_1$ be central charges in the upper half-plane,
attached to the source vertex $0$, which carries $V$ and $U$, and to the
target vertex $1$, which carries $W$; let $n_0=\dim\ker D$, $n_1=\dim W$,
$Z=n_0z_0+n_1z_1$ and $\varphi=\arg Z\in(0,\pi)$. The pair $(z_0,z_1)$ is the
central charge of a stability condition on the bounded derived category of
representations whose heart is the category of representations; the phase of
a representation of dimension vector $(a,b)$ is $\arg(az_0+bz_1)$, which is
$\pi$ times the usual normalized phase. The endomorphism of
\cite[\S5.4]{hkkp26} is $\Omega=\Omega_V\oplus\Omega_U\oplus\Omega_W$; it is
computed in \eqref{eq:coeff} below. Put
\[
 L=DD^\dagger,\qquad\lambda=\lambda_{\min}(L),\qquad
 D^+=D^\dagger(DD^\dagger)^{-1},\qquad s=\sum_j\norm{A_jD^+}_{\HS}^2 .
\]
For the metrics induced on $\ker D$ and $W$ let $H_j$ be the matrix of
$A_j|_{\ker D}$ in orthonormal frames, $C_0=\sum_jH_j^\dagger H_j$ and
$C_1=\sum_jH_jH_j^\dagger$. For the cohomology representation with these
metrics, $\Omega$ is $(z_0+C_0)\oplus(z_1-C_1)$, a scalar standing for the
corresponding multiple of the identity. We write $\norm\cdot_1$ for
the trace norm, and $(z_0+C_0)\oplus0$ for the endomorphism of
$V=\ker D\oplus(\ker D)^\perp$ that is $z_0+C_0$ on $\ker D$ and zero on the
orthogonal complement.

\begin{theorem}[Cancellation of the acyclic part]\label{thm:introA}
For all charges $z_0,z_1$ in the upper half-plane and every complex
\eqref{eq:introcomplex}, along the phase flow $\lambda(t)\to\infty$,
$s(t)\to0$, and for $t\geq t_1$
\[
 \norm{\Omega_V-(z_0+C_0)\oplus0}_1+\norm{\Omega_U}_1+\norm{\Omega_W-(z_1-C_1)}_1
 \leq C_{\mathrm{can}}\,e^{-\frac14(\arg z_1)(t-t_1)},
\]
where $C_{\mathrm{can}}$ is an explicit function of the dimensions and the
charges, and $t_1$ of these and of three numbers computed from the initial
metrics. No hypothesis is made on the cohomology representation.
\end{theorem}

The charges are \emph{ordered} if $\arg z_0>\arg z_1$, that is, if the simple
representation at the source vertex has the larger phase, and
\emph{reversed} if $\arg z_0<\arg z_1$. For reversed charges no
representation with $n_0,n_1\geq1$ is semistable. For ordered charges there
are unique $\theta_0,\theta_1>0$ with
$\arg(z_0+\theta_0)=\arg(z_1-\theta_1)=\varphi$; then
$n_0\theta_0=n_1\theta_1$, and a representation of dimension vector
$(n_0,n_1)$ is semistable, in the sense that no subrepresentation has phase
larger than $\varphi$, exactly when it is semistable in the sense of King
\cite{kin94} for the character $(-\theta_0,\theta_1)$. We call
$(\theta_0,\theta_1)$ the \emph{stability parameter}. The \emph{moment map} of
the cohomology representation, shifted by the stability parameter, is
$\mu=(C_0-\theta_0,\,C_1-\theta_1)$. Its zeros are the pairs of metrics for
which the phases $\arg(z_0+C_0)$ and $\arg(z_1-C_1)$ at the two vertices both
equal $\varphi$, that is, the harmonic metrics; they exist exactly for polystable
representations \cite{kn79,kin94}.

\begin{theorem}[Convergence]\label{thm:introB}
Along the phase flow, for every initial metric:
\begin{enumerate}
\item[(i)] If the charges are ordered and the cohomology representation is
 semistable, then $\mu(t)\to0$ and $m(h(t))\to|Z|$.
\item[(ii)] If the charges are ordered and the cohomology representation is
 polystable, the induced metrics
 $(e^{\varphi t}h_V|_{\ker D},\,e^{\varphi t}h_W)$ converge exponentially, as
 Hermitian forms on the fixed vector spaces $\ker D$ and $W$, to a zero of
 the moment map, that is, to a harmonic metric of the cohomology
 representation; the moment map and $m-|Z|$ tend to zero exponentially.
\item[(iii)] If the charges are reversed, the arrows $H_j$ tend to zero
 exponentially and $m(h(t))\to n_0|z_0|+n_1|z_1|$.
\end{enumerate}
\end{theorem}

A harmonic metric of a representation is not stationary under the flow but
moves by the scalar $e^{-\varphi t}$, which is why the limit is stated for
the normalized metrics. The complex itself has no harmonic metric, since the
phase on $U$ is negative and the phases on $V,W$ are positive, and its
metrics do not converge even after normalization, because $\lambda\to\infty$;
the limit exists for the metrics induced on the cohomology. In
\cite[Remark~5.15]{hkkp26} the metrics of an object are the metrized
complexes isomorphic to it in the derived category. The flow preserves the
underlying complex, and Theorem~\ref{thm:introB} identifies its limit with a
harmonic metric on another representative of the same object, the
cohomology representation. The rate in (ii)
is governed by the least eigenvalue of the Hessian of the Kempf--Ness
function transversally to the zeros of the moment map and by the rate of
Theorem~\ref{thm:introA}. The precise statements are
Theorems~\ref{thm:cancellation}, \ref{thm:relaxation}, \ref{thm:convergence}
and~\ref{thm:reversed}.

The proof of Theorem~\ref{thm:introA} does not use a Lyapunov function: the
mass is not monotone, and the function $\cK$ introduced before
Theorem~\ref{thm:introE} is monotone for the charges of that theorem but not
for all charges. It rests on an a priori bound for the
following quantity, which is not monotone either, and on the growth of
$\lambda$. Let $R(b)=(1-e^{-b})/b$ and
\[
 N=\sum_j\tr\bigl(A_jR(D^\dagger D)A_j^\dagger\bigr),
\]
the \emph{arrow energy}; it is the term of $-\cS_{\CC}$, see
\eqref{eq:primitive}, that does not contain the charges, and it bounds $\tr C_0$ and $(1-e^{-\lambda})s$.

\begin{theorem}[A priori bound]\label{thm:introC}
For all charges in the upper half-plane, along the phase flow
\[
 \frac{dN}{dt}\leq-\rho N+n_1M_1,\qquad
 \rho=\pi-\max\bigl\{\arg z_0,\tfrac{2\pi}5\bigr\}>0,
\]
with $M_1=\max\{\pi\max(\Rea z_1,0),\,2\Ima z_1\}$. Hence
$N(t)\leq\max\{N(0),n_1M_1/\rho\}$ for all $t\geq0$, and
$N(t)\leq2n_1M_1/\rho$ after an explicit time.
\end{theorem}

The three terms of $dN/dt$ come from the three blocks of $\Omega$
(Lemma~\ref{lem:energy}). The block $\Omega_W$ is a normal matrix and its
term is elementary. The blocks $\Omega_V$ and $\Omega_U$ are not normal, and
their terms are controlled by two matrix inequalities, stated and proved
independently of the flow in Sections~\ref{sec:kernel} and~\ref{sec:phase}.

\begin{theorem}[Two matrix inequalities]\label{thm:introD}
The following hold.
\begin{enumerate}
\item[(i)] Let $L>0$ and $E\geq0$ be Hermitian matrices, and let $J_L(E)$ be
 the integral of $e^{-xL}Ee^{-yL}$ over the simplex $x,y\geq0$, $x+y\leq1$.
 Then the negative part of the Hermitian matrix $J_L(E)$ satisfies
 $\tr J_L(E)_-\leq\frac25\tr J_L(E)$.
\item[(ii)] Let $\Omega',\Omega''$ be matrices with the same positive
 definite imaginary part and $\Phi',\Phi''$ the arguments of the unitary
 factors in their polar decompositions, with spectrum in $(0,\pi)$. Then
 $\tr\bigl((\Omega''-\Omega')(\Phi''-\Phi')\bigr)\leq0$, with equality only
 if $\Omega'=\Omega''$.
\end{enumerate}
\end{theorem}

By the Hermite--Genocchi formula for divided differences, in an eigenbasis
of $L$ the matrix
$J_L(E)$ is the entrywise product of $E$ with the matrix of second divided
differences of $e^{-x}$ at the points $\ell_i,0,\ell_k$; this matrix is not
positive semidefinite, and (i) says that the negative part of the product
carries at most two fifths of the trace. The best constant is at least
$(\sqrt2-1)/2$. The proof reduces (i) to an inequality in one real variable,
which is proved by estimates near $0$ and $\infty$ and by interval
arithmetic on a compact interval (Section~\ref{sec:kernel}). Statement (ii)
says that minus the argument of the polar factor is a strictly monotone
operator, in the sense of Minty \cite{min62}, on the Hermitian matrices with
the trace inner product; it is not monotone for the L\"owner order, because
the
arctangent is not operator monotone on the real line
(Remark~\ref{rem:phasefails}). For a scalar imaginary part it is the fact
that the argument of a complex number decreases in its real part. The proof
uses an integral representation of anti-L\"owner type \cite{aud11}
(Section~\ref{sec:phase}).

The mass is not monotone along the phase flow, also for the charges and
complexes of Theorem~\ref{thm:introB}: we give a complex with $512$ arrows,
stable cohomology and $r=1$, with metrics at which the mass increases, for
$z_0=-1+\ii$, $z_1=1+\ii$ (Proposition~\ref{prop:mass}). The largest phase
can increase and the smallest phase can decrease
(Proposition~\ref{prop:extreme}), whereas for representations the mass and
the extreme phases are monotone, as expected in \cite[\S4.1]{hkkp26}
(Proposition~\ref{prop:representations}). What converges for the complexes
\eqref{eq:introcomplex}, for ordered charges and semistable cohomology, is
the mass measure of \cite[\S4.2]{hkkp26}, recalled in
Section~\ref{sec:expectations} (Proposition~\ref{prop:measure}).

For a range of charges there is a Lyapunov function. The \emph{phase
potential} is
\[
 \cK=\frac{\Ima(e^{-\ii\varphi}\cS_{\CC})}{\sin\varphi},
\]
with $\cS_{\CC}$ the function \eqref{eq:primitive}. We call
$z_0=-1+\ii$, $z_1=1+\ii$ the \emph{symmetric charges}; for them and
$n_0=n_1$,
\[
 \cK(h)=\log\det h_V-\log\det h_U-\log\det h_W+\tr\Ein(D^\dagger D)+N,
 \qquad\Ein(b)=\int_0^1(1-e^{-xb})\frac{dx}x ,
\]
with the determinants of the Gram matrices in fixed reference bases.

\begin{theorem}[Lyapunov function]\label{thm:introE}
Let the charges be ordered, and suppose that $\varphi\geq\frac{2\pi}5$ or
$r=1$. Then the phase potential is strictly decreasing along the phase flow.
For the symmetric charges and $n_0=n_1$,
\[
 -\frac{d\cK}{dt}\ \geq\ \frac\pi{10}\tr J_L(DCD^\dagger)+\frac\pi2\tr e^{-L}>0,
 \qquad C=\sum_jA_j^\dagger A_j .
\]
There are ordered charges and complexes with $r=2$, with metrics at which
the phase potential increases along the flow.
\end{theorem}

The number $\frac{2\pi}5$ is $\pi$ times the constant of
Theorem~\ref{thm:introD}(i). The rate of the potential is a sum of three
terms, one for each of $V$, $W$ (degree zero) and $U$ (degree one); the first
two are nonnegative for elementary reasons, and the third is controlled by
Theorem~\ref{thm:introD}(i). See Theorem~\ref{thm:monotone} and
Propositions~\ref{prop:reversed}, \ref{prop:orderedfailure}.

\begin{theorem}[Infimum and moduli]\label{thm:introF}
Let the charges be ordered, and let $\gamma$ be Euler's constant.
\begin{enumerate}
\item[(i)] The potential depends on reference bases of $V,U,W$ only through
 their volume elements. With volume elements on $\ker D$ and $W$ induced
 from those of $V,U,W$ through $\det V\cong\det\ker D\otimes\det U$, the
 infimum of the phase potential over all metrics of the complex is the
 infimum of the Kempf--Ness function \eqref{eq:KN} of its cohomology
 representation for the stability parameter $(\theta_0,\theta_1)$, plus
 $\gamma\,\theta_0r$; it is not attained. Consequently
 the infima for two quasi-isomorphic complexes of the class
 \eqref{eq:introcomplex} with semistable cohomology and corresponding volume
 elements differ by $\gamma\,\theta_0(r'-r)$.
\item[(ii)] On the stable locus of a holomorphic family of complexes,
 $\ii\partial\bar\partial$ of the infimum is the pullback of the K\"ahler form
 of the moduli space of stable representations.
\item[(iii)] For semistable cohomology the potential along the flow tends to
 its infimum, exponentially and locally uniformly on the stable locus.
\end{enumerate}
\end{theorem}

Part (ii) follows from (i) and the classical description of the K\"ahler
potential of a K\"ahler quotient (Lemma~\ref{lem:kappa}). See
Theorems~\ref{thm:infimum}, \ref{thm:moduli} and~\ref{thm:families}.

\subsection{Remarks on the setting and the method}
The phase flow is not the gradient flow of the phase potential
(Section~\ref{sec:potential}). For a representation
regarded as a complex in degree zero $\Omega$ is normal, the phase
potential is the Kempf--Ness function, and its monotonicity along the flow
follows from formulas displayed in \cite[\S\S4.2, 5.1]{hkkp26}. On a
representation the flow has the same stationary points modulo scaling as the
downward gradient flow of the Kempf--Ness function, and its velocity pairs
negatively with the gradient, but it is not that gradient flow
(Section~\ref{sec:convergence}).

The estimates for the polar decomposition that isolate one block of a
matrix, and the differential inequality for $s$, are taken from \cite{wue26}
and recalled in Proposition~\ref{prop:imported}. What is new in the proof of
Theorem~\ref{thm:introA} is the control of the arrows and of $\lambda$ before
$\lambda$ is large: Theorem~\ref{thm:introC} bounds the arrows, and a second
differential inequality, for the trace of the exponential integral of $L$,
forces $\lambda$ to grow (Theorem~\ref{thm:growth}). For the charges of
Theorem~\ref{thm:introE} and semistable cohomology the Lyapunov function
gives the same control with different constants
(Proposition~\ref{prop:apriori}). Theorem~\ref{thm:introB} then follows from
the convexity of the Kempf--Ness function; the relevant statements are
classical \cite{kn79,grs21} and are proved for the present linear action in
Section~\ref{sec:convergence}.

The following is not proved here. For semistable cohomology that is not
polystable the induced metrics diverge
(Proposition~\ref{prop:divergence}); the asymptotics are determined here
only for the smallest such representation, regarded as a complex in degree
zero, where the metrics have the form $t^{\mp1/2}$ on a subrepresentation
and the quotient (Proposition~\ref{prop:semistable}). In general the
normalized induced metrics stay at bounded distance from a trajectory of the
gradient flow of the Kempf--Ness function for the invariant Riemannian
metric weighted by the imaginary parts of the charges
\cite[Theorem~6.3 and Corollary~6.5]{wue26b}, whose asymptotics are
described in \cite{hkkp23}; whether the bounded factor converges is open
beyond this example. The case
$\arg z_0=\arg z_1$ is covered by Theorems~\ref{thm:introA}
and~\ref{thm:introC} but not by Theorem~\ref{thm:introB}. Complexes whose
term of degree one sits at the target vertex, with acyclic parts at both
vertices, or with terms in more than two degrees, are not treated. For
cohomology that is not semistable, Theorems~\ref{thm:introA}
and~\ref{thm:introC} apply, and the induced metrics follow the
Harder--Narasimhan filtration of the cohomology representation
\cite[Theorem~6.11]{wue26b} (Remark~\ref{rem:unstable}).

\subsection{Organization}
Section~\ref{sec:setting} recalls the endomorphism $\Omega$, the flow and the
estimates
used from \cite{wue26}. Sections~\ref{sec:kernel} and~\ref{sec:phase} prove
Theorem~\ref{thm:introD}. Section~\ref{sec:energy} proves
Theorems~\ref{thm:introC} and~\ref{thm:introA}, and
Section~\ref{sec:convergence} Theorem~\ref{thm:introB}.
Section~\ref{sec:potential} treats the phase potential and proves
Theorem~\ref{thm:introE}, and Section~\ref{sec:moduli} proves
Theorem~\ref{thm:introF}. Section~\ref{sec:expectations} treats the mass, the
extreme phases and the fixed points, and Section~\ref{sec:remarks} the charges
for which the phase potential is not monotone, and open questions.

\section{The endomorphism \texorpdfstring{$\Omega$}{Omega}, the flow and the function \texorpdfstring{$\mathcal S_{\mathbb C}$}{S}}\label{sec:setting}

This section fixes notation and recalls what is used from \cite{hkkp26} and
\cite{wue26}. Nothing in it is new. Throughout, $T^\dagger$ is the adjoint of
a linear map for the Hermitian metrics in force (the conjugate transpose in
orthonormal frames), $T^*$ is the conjugate transpose of a matrix in fixed
reference bases, and $D^+$ is the Moore--Penrose inverse. Inner products
$\inner\cdot\cdot$ and Hermitian forms are conjugate-linear in the first
argument, and a scalar in a matrix expression denotes the corresponding
multiple of the identity. For a square matrix
$T$ we write
$\Rea T=(T+T^\dagger)/2$ and $\Ima T=(T-T^\dagger)/2\ii$; $\norm T_1$ is the
trace norm, $\norm T_{\HS}$ the Hilbert--Schmidt norm and $\norm T$ the
operator norm; for a Hermitian matrix $M$, $M_\pm\geq0$ are its positive and
negative parts, $M=M_+-M_-$, $M_+M_-=0$. The letter $\varphi$ always denotes
the argument of the total central charge; $\phi$ with indices denotes phases
of blocks and their eigenvalues. Likewise $\vartheta$, $\varkappa$,
$\varrho$, $\varsigma$, $\varpi$ denote objects different from those denoted
by $\theta$, $\kappa$, $\rho$, $\sigma$, $\pi$.

\subsection{Complexes and the endomorphism \texorpdfstring{$\Omega$}{Omega}}
Let $V,U,W$ be finite-dimensional complex vector spaces, $D:V\to U$ a linear
map and $A_1,\dots,A_k:V\to W$ linear maps. This is a two-term complex of
representations of the generalized Kronecker quiver, with two vertices $0,1$
and $k$ arrows from $0$ to $1$:
\begin{equation}\label{eq:complex}
 \mathsf E:\qquad(V\xrightarrow{\,(A_j)\,}W)\ \xrightarrow{\,(D,0)\,}\ (U\to0),
\end{equation}
with $V$ and $W$ in degree $0$ and $U$ in degree $1$. We assume throughout
that
\[
 D\text{ is surjective},\qquad n_0:=\dim\ker D\geq1,\quad n_1:=\dim W\geq1,
 \quad r:=\dim U\geq1 .
\]
The \emph{cohomology representation} is $\ker D\to W$ with the arrows
$A_j|_{\ker D}$, of dimension vector $(n_0,n_1)$. It is a subcomplex of
$\mathsf E$, and the quotient $(V/\ker D\to0)\to(U\to0)$ is an acyclic
complex, which we call the \emph{acyclic part}; $r$ is its rank. Fix
\emph{central charges} $z_0,z_1$ in the open upper half-plane, attached to
the two vertices, and write
\[
 z_i=x_i+\ii y_i=|z_i|e^{\ii\alpha_i},\qquad y_i>0,\quad\alpha_i\in(0,\pi).
\]

Let $h=(h_V,h_U,h_W)$ be Hermitian metrics. Adjoints refer to the current
metrics. Put
\[
 B=D^\dagger D,\qquad L=DD^\dagger,\qquad C=\sum_jA_j^\dagger A_j,
\]
and, for Hermitian matrices $P\geq0$ and $M$,
\begin{equation}\label{eq:J}
 \begin{gathered}
 R(P)=\int_0^1e^{-xP}dx,\qquad
 K_P(M)=\int_0^1e^{-(1-x)P}Me^{-xP}dx,\\
 J_P(M)=\int_{\substack{x,y\geq0\\x+y\leq1}}e^{-xP}Me^{-yP}\,dx\,dy .
 \end{gathered}
\end{equation}
In \cite[\S5.4]{hkkp26} a complex of quiver representations with metrics
carries an endomorphism $\Omega$ of its underlying graded vector space,
defined through a superconnection in a graded algebra attached to the
quiver: $\Omega=\varepsilon\,[\exp(-\sqrt{-1}F)]_f$, where
$F=(A-A^*)^2+\sqrt{-1}\sum_if_i\otimes z_i$ is formed from the element
$A=\sum_ie_i\otimes d_i+\sum_aa\otimes\rho(a)$ of the algebra of
\cite[Definition~5.13]{hkkp26} tensored with the endomorphisms, with $d_i$
the differentials and $\rho(a)$ the maps of the arrows, $\varepsilon$ is the
grading, and $[\,\cdot\,]_f$ is the component along the basis elements
$f_i$. Apart from $\Omega$ and $z_i$, the symbols of this sentence are those
of \cite{hkkp26} and have this meaning only here; $A^*$ is formed from the
involution of that algebra and the adjoints for the metrics
\cite[(5.14)]{hkkp26}. The
differential of the function $\cS_{\CC}$ below is
$\xi\mapsto\tr(\Omega\xi)$; the trace of $\Omega$ is the central charge
\eqref{eq:central}, and for a representation $\Omega$ is the central charge
of each vertex shifted by the moment map. For \eqref{eq:complex} it is
$\Omega=\Omega_V\oplus\Omega_U\oplus\Omega_W$ with
\cite[Proposition~2.1]{wue26}
\begin{equation}\label{eq:coeff}
 \Omega_V=z_0e^{-B}+K_B(C),\qquad
 \Omega_U=J_L(DCD^\dagger)-z_0e^{-L},\qquad
 \Omega_W=z_1\id-\sum_jA_jR(B)A_j^\dagger .
\end{equation}
We abbreviate
\begin{equation}\label{eq:KJG}
 K=K_B(C),\qquad J=J_L(DCD^\dagger),\qquad G=\sum_jA_jR(B)A_j^\dagger\geq0 ,
\end{equation}
three Hermitian matrices that do not depend on the charges. Hence
\begin{equation}\label{eq:imaginary}
 \Ima\Omega_V=y_0e^{-B}>0,\qquad
 \Ima(-\Omega_U)=y_0e^{-L}>0,\qquad
 \Ima\Omega_W=y_1\id>0,
\end{equation}
and \cite[Lemma~5.17]{hkkp26}
\begin{equation}\label{eq:central}
 \tr\Omega_V+\tr\Omega_U+\tr\Omega_W=n_0z_0+n_1z_1=:Z .
\end{equation}
The matrix $\Omega_W$ is normal. The matrix $\Omega_V$ is in general not
normal \cite[Example~2.3]{wue26}, and neither is $\Omega_U$, whose Hermitian
and anti-Hermitian parts $J-x_0e^{-L}$ and $-y_0e^{-L}$ commute only if $J$
commutes with $e^{-L}$.

\subsection{The phase and the flow}
For a matrix $T$ with $\Ima T>0$ there is a unique Hermitian matrix $\Psi$
with spectrum in $(0,\pi)$ such that $Te^{-\ii\Psi}>0$, the \emph{polar phase}
of $T$; it is the argument of the unitary factor in the polar decomposition
$T=(TT^\dagger)^{1/2}e^{\ii\Psi}$ \cite[Section~7.3]{hj13}: if
$e^{\ii\Psi}v=\zeta v$, then
$\inner v{(\Ima T)v}=(\Ima\zeta)\inner v{|T|v}$, so the eigenvalues of the
unitary factor lie in the open upper half-circle (see also
\cite[Lemma~2.2]{wue26}). Here
$|T|:=(TT^\dagger)^{1/2}$ is the left modulus; the unitary factor is the same
in the left and right polar decompositions, and so is the trace norm (in
\cite{bha97} the modulus is the right one, $(T^\dagger T)^{1/2}$). The phases
$\Phi_V,\Phi_W$ are the polar phases of $\Omega_V,\Omega_W$, and
$\Phi_U+\pi$ is the polar phase of $-\Omega_U$; so
\begin{equation}\label{eq:phase}
 |\Omega_\bullet|:=\Omega_\bullet e^{-\ii\Phi_\bullet}
 =(\Omega_\bullet\Omega_\bullet^\dagger)^{1/2}>0,\qquad\bullet\in\{V,U,W\},
\end{equation}
with $\spec\Phi_V,\spec\Phi_W\subset(0,\pi)$ and
$\spec\Phi_U\subset(-\pi,0)$. This is the phase defined by positivity in
\cite[\S4.2]{hkkp26}, with the branches of \cite[(5.15)]{hkkp26}. In
\cite[\S5.4]{hkkp26} the phase of $\Omega$ for a complex is written with the
functional calculus. For normal $\Omega$, in particular for
representations, the two readings coincide; for complexes $\Omega$ is not
normal in general \cite[Example~2.3]{wue26}. The flow studied here is the
one defined with the polar phase:
\begin{equation}\label{eq:flow}
 h_\bullet^{-1}\dot h_\bullet=-\Phi_\bullet,\qquad\bullet\in\{V,U,W\}.
\end{equation}
We call it the \emph{phase flow}. It is the minimizing flow of
\cite[\S4.1]{hkkp26} with the phase read as above; we use a neutral name
because on complexes the mass can increase along it
\cite[Theorem~1.1]{wue26}. The \emph{mass} is $m(h)=\sum_\bullet\tr|\Omega_\bullet|\geq|Z|$. The
vector field is smooth and its solutions exist for all $t\geq0$
\cite[Lemma~2.4]{wue26}. In orthonormal frames moving with the metrics the
maps evolve by
\begin{equation}\label{eq:mapflow}
 \dot D=\tfrac12(D\Phi_V-\Phi_UD),\qquad
 \dot A_j=\tfrac12(A_j\Phi_V-\Phi_WA_j).
\end{equation}
A \emph{state} is a tuple $(D,A_j)$ of matrices in orthonormal frames, up to
unitary changes of the frames; equivalently, a triple of metrics on the
fixed complex up to isometry, or a point of the orbit of $(D,A_j)$ under
$\mathrm{GL}(V)\times\mathrm{GL}(U)\times\mathrm{GL}(W)$ modulo the unitary
groups. The phases, mass and arrow energy are invariant under these changes of frames. The primitive below also uses fixed reference volume elements.

\begin{lemma}\label{lem:eigenbasis}
Let $T$ be a matrix with $\Ima T>0$ and polar phase $\Psi$, and let $v$ be a
unit vector with $\Psi v=\psi v$. Then
$\inner v{Tv}=p\,e^{\ii\psi}$ with $p=\inner v{|T|v}>0$. Consequently, in an
orthonormal eigenbasis $(v_k)$ of $\Psi$ with eigenvalues $\psi_k$ and
$p_k=\inner{v_k}{|T|v_k}$, $\tr(|T|f(\Psi))=\sum_kp_kf(\psi_k)$ for every
real function $f$.
\end{lemma}

\begin{proof}
$T=|T|e^{\ii\Psi}$, so $\inner v{Tv}=e^{\ii\psi}\inner v{|T|v}$. The second
statement is the trace computed in the basis $(v_k)$.
\end{proof}

The lemma applies to $\Omega_V$, $\Omega_W$ and $-\Omega_U$; for the block
$U$ it gives $\inner v{\Omega_Uv}=p\,e^{\ii\phi}$ for a unit eigenvector $v$
of $\Phi_U$ with eigenvalue $\phi\in(-\pi,0)$.

\subsection{The primitive}
Let $\Ein(b)=\int_0^1(1-e^{-xb})\,dx/x$. With respect to fixed reference
bases, in which $\log\det h_\bullet$ denotes the logarithm of the Gram
determinant, put
\begin{equation}\label{eq:primitive}
 \cS_{\CC}(h)=z_0\bigl(\log\det h_V-\log\det h_U\bigr)+z_1\log\det h_W
 +z_0\tr\Ein(B)-\tr\bigl(CR(B)\bigr).
\end{equation}
Then \cite[Proposition~7.1]{wue26}, for every variation
$\xi_\bullet=h_\bullet^{-1}\delta h_\bullet$ of the metrics,
\begin{equation}\label{eq:primitivevariation}
 d\cS_{\CC}(\xi)=\sum_\bullet\tr(\Omega_\bullet\xi_\bullet).
\end{equation}
Both sides are affine in $(z_0,z_1)$, so the terms without charges agree:
\begin{equation}\label{eq:energyvariation}
 d\bigl(\tr(CR(B))\bigr)(\xi)=-\tr(K\xi_V)-\tr(J\xi_U)+\tr(G\xi_W).
\end{equation}

\subsection{Charges and the stability parameter}
Put
\[
 \varphi=\arg Z\in(0,\pi),\quad\sigma=\sin\varphi,\quad
 c_0=|z_0|\sin(\alpha_0-\varphi),\quad c_1=|z_1|\sin(\varphi-\alpha_1).
\]
Since $\Ima(e^{-\ii\varphi}Z)=0$ we have $n_0c_0=n_1c_1$. The charges are
\emph{ordered} if $\alpha_0>\alpha_1$ and \emph{reversed} if
$\alpha_0<\alpha_1$; they are ordered exactly when $c_0>0$, or when
$\alpha_1<\varphi<\alpha_0$. For ordered charges put
\begin{equation}\label{eq:theta}
 \theta_0=c_0/\sigma,\qquad\theta_1=c_1/\sigma,\qquad
 \vartheta=n_0\theta_0=n_1\theta_1>0 ;
\end{equation}
then $z_0+\theta_0$ and $z_1-\theta_1$ have argument $\varphi$. We call the
pair $(\theta_0,\theta_1)$ the \emph{stability parameter}. A
representation whose dimension vector $(a_0,a_1)$ satisfies
$a_0\theta_0=a_1\theta_1$, for instance $(n_0,n_1)$, has phase $\varphi$. It
is \emph{semistable} if $a\theta_0-b\theta_1\leq0$ for the dimension vector
$(a,b)$ of each of its subrepresentations, that is, if no subrepresentation
has phase $\arg(az_0+bz_1)$ larger than $\varphi$; it is \emph{stable} if the
inequality is strict for proper nonzero subrepresentations, and
\emph{polystable} if it is a direct sum of stable representations of phase
$\varphi$. In the convention of King \cite[Definition~1.1]{kin94} this is
stability for the character $(-\theta_0,\theta_1)$, a positive multiple of
the integral character $(-n_1,n_0)$; the Kempf--Ness function \eqref{eq:KN}
carries the weights $(\theta_0,-\theta_1)$ on
$(\log\det h_0,\log\det h_1)$. For reversed charges no
representation with $n_0,n_1\geq1$ is semistable, because the
subrepresentation supported at the vertex $1$ has phase
$\alpha_1>\varphi$.

\begin{definition}\label{def:admissible}
Let $(n_0,n_1,r)$ be given. The charges $(z_0,z_1)$ are
\emph{admissible} if they are ordered and either $\varphi\geq\frac\pi3$ or
$r=1$.
\end{definition}

This is the range of charges covered by Theorem~\ref{thm:monotone}; the
number $\frac\pi3$ comes from the constant of Theorem~\ref{thm:kernel}.

The \emph{symmetric charges} are $z_0=-1+\ii$, $z_1=1+\ii$. If moreover
$n_0=n_1$, then $\varphi=\frac\pi2$, $\sigma=1$, $\theta_0=\theta_1=1$, and they
are admissible for every $r$.

\subsection{The splitting adapted to \texorpdfstring{$\ker D$}{ker D}}
Choose orthonormal frames adapted to $V=\ker D\oplus(\ker D)^\perp$ and write
\[
 D=[0,F],\qquad A_j=[H_j,T_j],\qquad
 \Phi_V=\begin{pmatrix}\Phi_{HH}&\Phi_{HA}\\\Phi_{AH}&\Phi_{AA}\end{pmatrix}.
\]
Here the indices $H$ and $A$ refer to $\ker D$ and to its orthogonal
complement, which carries the acyclic part. So $F$ is invertible and the
$H_j$ are the matrices of the arrows of the cohomology representation for the
induced metrics. Let
$D^+=D^\dagger(DD^\dagger)^{-1}$ be the Moore--Penrose inverse of $D$ and
\[
 X_j=A_jD^+=T_jF^{-1}:U\to W .
\]
These maps vanish exactly when the arrows vanish on $(\ker D)^\perp$. Put
\begin{equation}\label{eq:qslambda}
 C_0=\sum_jH_j^\dagger H_j,\quad C_1=\sum_jH_jH_j^\dagger,\quad
 q=\tr C_0=\tr C_1,\quad s=\sum_j\norm{X_j}_{\HS}^2,\quad
 \lambda=\lambda_{\min}(L),
\end{equation}
let $\ell_1\leq\dots\leq\ell_r$ be the eigenvalues of $L=FF^\dagger$, so
$\lambda=\ell_1$, and
\begin{equation}\label{eq:Y}
 Y=\sum_jX_j(\id-e^{-L})X_j^\dagger\geq0,\qquad G=C_1+Y,\qquad
 \tr Y\geq(1-e^{-\lambda})\,s .
\end{equation}
With the choice of the moving frames made in \cite[\S6.1]{wue26} (the frames
of $V$ are rotated so that the splitting is preserved),
\begin{equation}\label{eq:adapted}
 \dot F=\tfrac12(F\Phi_{AA}-\Phi_UF),\qquad
 \dot H_j=\tfrac12(H_j\Phi_{HH}-\Phi_WH_j),
\end{equation}
and
\begin{equation}\label{eq:adaptedcoeff}
 \Omega_V=\begin{pmatrix}z_0\id+C_0&E_R\\E_R^\dagger&z_0e^{-F^\dagger F}+K_{F^\dagger F}(C_T)\end{pmatrix},
 \qquad\Omega_W=z_1\id-C_1-Y,
\end{equation}
with $C_T=\sum_jT_j^\dagger T_j$ and $E_R=\sum_jH_j^\dagger T_jR(F^\dagger F)$.
For the cohomology representation with its induced metrics the
endomorphism $\Omega$ is
\[
 \Omega_H=(z_0\id+C_0)\oplus(z_1\id-C_1);
\]
it is normal, with phases the Hermitian matrices
$\phi_0=\arg(z_0\id+C_0)$ and $\phi_1=\arg(z_1\id-C_1)$, and mass
$m_H=\norm{z_0\id+C_0}_1+\norm{z_1\id-C_1}_1$. We write
\[
 \delta_0=\Phi_{HH}-\phi_0,\qquad\delta_1=\Phi_W-\phi_1,\qquad
 \mathcal E_V=\Omega_V-\diag\bigl(z_0\id+C_0,\ z_0e^{-F^\dagger F}\bigr).
\]
We say that \emph{the acyclic part cancels} along a trajectory if
$\lambda\to\infty$ and
$\norm{\Omega_V-(z_0\id+C_0)\oplus0}_1+\norm{\Omega_U}_1+\norm{\Omega_W-(z_1\id-C_1)}_1\to0$.
The name refers to the central charge, to which the two terms of the acyclic
part contribute $z_0r$ and $-z_0r$; the statement is that their
contributions to $\Omega$ itself tend to zero.

For ordered charges we write
$\mu=(\mu_0,\mu_1)=(C_0-\theta_0\id,\ C_1-\theta_1\id)$ and call it the
\emph{moment map} of the cohomology representation. The pair
$(-\mu_0,\mu_1)$ is the gradient of the Kempf--Ness function
(Lemma~\ref{lem:knvariation}), that is, up to a constant factor the moment
map of $\mathrm U(n_0)\times\mathrm U(n_1)$ in the usual convention, shifted
by King's character; we change the sign of the first component so that both
components read $C_i-\theta_i$. A pair of metrics on
$\ker D$ and $W$ with $\mu=0$ is a \emph{zero of the moment map}; such pairs
exist exactly when the cohomology representation is polystable
\cite[Proposition~6.5]{kin94}, which extends the theorem of Kempf and Ness
\cite{kn79}, and they are the harmonic metrics of \cite{hkkp26} for a
representation, since $\mu=0$ if and only if both phases $\phi_0,\phi_1$ equal
$\varphi$.

\subsection{Estimates for \texorpdfstring{$A_jD^+$}{A\_jD+} and for the phases}
We use the following facts from \cite{wue26}, valid for all charges in the
upper half-plane and all dimensions. Put $\kappa_0=\pi/(2e)$ and
$d_0=\min(n_0,r)$.

\begin{proposition}[{\cite{wue26}}]\label{prop:imported}
Along the phase flow and at every metric:
\begin{enumerate}
\item[(i)] $\lambda$ is nondecreasing along the flow. By the formula for the
 one-sided derivative of the least eigenvalue of a differentiable family of
 Hermitian matrices, its right derivative is the minimum of
 $\lambda\bigl(\inner v{\Phi_{AA}v}-\inner u{\Phi_Uu}\bigr)$ over the unit
 vectors $v$ with $F^\dagger Fv=\lambda v$, where $u=Fv/\sqrt\lambda$
 \cite[Lemma~4.3 and \S6.1]{wue26}.
\item[(ii)] $\norm{E_R}_{\HS}\leq\sqrt{qs/\lambda}$,
 $\norm{K_{F^\dagger F}(C_T)}_1\leq r\kappa_0s$, $\norm{J}_1\leq rs$ and
 $\norm{Y}_1\leq s$; hence
 $\norm{\mathcal E_V}_1\leq2\sqrt{d_0qs/\lambda}+r\kappa_0s$
 \cite[Lemma~4.2, \S6.1 and \S6.3]{wue26}.
\item[(iii)] $\dot s\leq\bigl(-\alpha_1+\frac{4\pi q}{y_0\lambda}\bigr)s$
 \cite[\S6.2]{wue26}.
\item[(iv)] If $\sqrt{qs/\lambda}\leq y_0/8$, then
 $\norm{\delta_0}_{\HS}\leq\frac{4\pi}{y_0}\sqrt{qs/\lambda}$; and
 $\norm{\delta_1}_{\HS}\leq s/y_1$ always
 \cite[Lemma~3.1 and \S6.2]{wue26}.
\item[(v)] Let $T_\circ=|T_\circ|\,e^{\ii\Psi_\circ}$ be an invertible
 matrix in polar form, with $\Psi_\circ$ Hermitian, and $T=T_\circ+M$ a
 matrix with $\Ima T>0$ and polar phase $\Psi$. If $v$ is a unit vector with
 $|T_\circ|v=pv$ and $\Psi_\circ v=\psi v$, where $p>0$ and $0<\psi<\pi$,
 then $|\inner v{(\Psi-\psi)v}|\leq\pi\sqrt{\norm M_1/p}$
 \cite[Lemma~3.2]{wue26}.
\end{enumerate}
\end{proposition}

In \cite{wue26} these estimates were used when the initial value of
$\lambda$ is large: then the acyclic part cancels, for all charges
\cite[Theorem~1.3]{wue26}, and for the symmetric charges the moment map of
the cohomology representation tends to zero if its Kempf--Ness function is
bounded below \cite[Theorem~1.4]{wue26}. Whether the assumption on the
initial value of $\lambda$ can be removed was left open there.

\section{A trace inequality for a Schur multiplier}\label{sec:kernel}

For a Hermitian $r\times r$ matrix $L>0$ and a Hermitian matrix $E\geq0$ put
\[
 J_L(E)=\int_{\substack{x,y\geq0\\x+y\leq1}}e^{-xL}Ee^{-yL}\,dx\,dy ,
\]
as in \eqref{eq:J}. This Hermitian matrix is in general indefinite
(Remark~\ref{rem:indefinite}). Recall that $M_\pm\geq0$ denote the positive
and negative parts of a Hermitian matrix $M=M_+-M_-$, $M_+M_-=0$.

\begin{theorem}[Trace inequality]\label{thm:kernel}
For all Hermitian $L>0$ and $E\geq0$,
\[
 \tr J_L(E)_-\leq\frac25\,\tr J_L(E).
\]
In particular $\tr J_L(E)\geq0$, and
$\tr J_L(E)_-\leq\frac27\tr J_L(E)_+$.
\end{theorem}

The proof occupies this section. It has three steps: a reduction to a Schur
multiplier, an explicit way of writing that multiplier as a difference of two
positive semidefinite kernels, and an inequality in one real variable, part
of which is verified with interval arithmetic.

\subsection{Reduction to a Schur multiplier}
Diagonalize $L=\diag(\ell_1,\dots,\ell_r)$ with $\ell_i>0$. Then
$J_L(E)_{ik}=\mathsf j(\ell_i,\ell_k)E_{ik}$, where
\[
 \mathsf j(a,b)=\int_{\substack{x,y\geq0\\x+y\leq1}}e^{-xa-yb}\,dx\,dy .
\]
Put $S=L^{-1}EL^{-1}\geq0$ and $\mathsf m(a,b)=ab\,\mathsf j(a,b)$. Then
\begin{equation}\label{eq:schurform}
 J_L(E)=S\circ\mathsf M,\qquad\mathsf M=[\mathsf m(\ell_i,\ell_k)]_{i,k},\qquad
 \tr J_L(E)=\sum_i\mathsf m(\ell_i,\ell_i)S_{ii},
\end{equation}
where $\circ$ denotes the entrywise (Schur) product. By the Hermite--Genocchi
formula $\mathsf j(a,b)$ is the second divided difference of $e^{-x}$ at the
points $a,0,b$. By the recursion for divided differences it equals
$-(R(a)-R(b))/(a-b)$ for $a\neq b$, with $R(x)=(1-e^{-x})/x$, so that
$J_L(E)$ is minus the Schur product of $E$ with the L\"owner matrix of $R$ at
the eigenvalues of $L$. The function $-R$ is increasing but, being entire and
not affine, it is not operator monotone; so by L\"owner's theorem
\cite{bha97} these matrices are not all positive semidefinite. Theorem~\ref{thm:kernel} bounds
the negative part of the Schur product in trace.

\begin{lemma}\label{lem:mformula}
Let $\eta(a)=1-e^{-a}$. For $a,b>0$, $a\ne b$,
\[
 \mathsf m(a,b)=\frac{b\,\eta(a)-a\,\eta(b)}{b-a},\qquad
 \mathsf m(a,a)=1-(1+a)e^{-a}>0 .
\]
\end{lemma}

\begin{proof}
Integrating in $y$ gives
$\mathsf j(a,b)=b^{-1}\int_0^1e^{-xa}(1-e^{-(1-x)b})\,dx
=b^{-1}[\eta(a)/a-(e^{-a}-e^{-b})/(b-a)]$. Multiply by $ab$ and use
$e^{-a}-e^{-b}=\eta(b)-\eta(a)$. The diagonal value is the limit $b\to a$,
and it is positive because its derivative is $ae^{-a}$.
\end{proof}

\subsection{An expansion of the kernel}
On $L^2[0,1]$ let $e_a(x)=e^{-ax}$, and let $\mathsf T$ be the integral
operator with kernel $\mathbf 1_{\{x+y>1\}}$. The integral operator with
kernel $\mathbf 1_{\{x+y\leq1\}}$ is
$\mathbf 1\otimes\mathbf 1-\mathsf T$, where $\mathbf 1\otimes\mathbf 1$ maps
$g$ to $(\int_0^1g)\mathbf 1$, and
$\mathsf j(a,b)=\inner{e_a}{(\mathbf 1\otimes\mathbf 1-\mathsf T)e_b}$.

\begin{lemma}\label{lem:eigen}
Put $\omega_k=(k+\frac12)\pi$ and $\chi_k(x)=\sqrt2\sin(\omega_kx)$ for
$k\geq0$. The functions $\chi_k$ form an orthonormal basis of $L^2[0,1]$,
and $\mathsf T\chi_k=(-1)^k\omega_k^{-1}\chi_k$.
\end{lemma}

\begin{proof}
The $\chi_k$ are the eigenfunctions of $-d^2/dx^2$ on $[0,1]$ with a
Dirichlet condition at $0$ and a Neumann condition at $1$, a self-adjoint
problem with simple eigenvalues $\omega_k^2$. Further
$(\mathsf T\chi_k)(x)=\int_{1-x}^1\sqrt2\sin(\omega_ky)\,dy
=\sqrt2\,\omega_k^{-1}[\cos(\omega_k(1-x))-\cos\omega_k]$, and
$\cos\omega_k=0$, $\sin\omega_k=(-1)^k$ give
$\cos(\omega_k(1-x))=(-1)^k\sin(\omega_kx)$.
\end{proof}

\begin{lemma}\label{lem:expansion}
For $a,b>0$,
\[
 \mathsf m(a,b)=\eta(a)\eta(b)-\sum_{k\geq0}\frac{(-1)^k}{\omega_k}\,t_k(a)t_k(b),
 \qquad
 t_k(a)=\frac{\sqrt2\,a\,\bigl(\omega_k-(-1)^kae^{-a}\bigr)}{a^2+\omega_k^2},
\]
and the series converges absolutely.
\end{lemma}

\begin{proof}
We have
$\mathsf m(a,b)=a\inner{e_a}{\mathbf 1}\,b\inner{\mathbf 1}{e_b}-\inner{ae_a}{\mathsf T\,be_b}$
and $a\inner{e_a}{\mathbf 1}=\eta(a)$. The operator $\mathsf T$ is symmetric and
Hilbert--Schmidt, so Lemma~\ref{lem:eigen} gives
\begin{align*}
 \inner{ae_a}{\mathsf Tbe_b}&=\sum_k\frac{(-1)^k}{\omega_k}\,
 a\inner{e_a}{\chi_k}\,b\inner{\chi_k}{e_b},\\
 \int_0^1e^{-ax}\sin(\omega x)\,dx
 &=\frac{\omega-e^{-a}(a\sin\omega+\omega\cos\omega)}{a^2+\omega^2},
\end{align*}
and at $\omega=\omega_k$ the integral equals
$(\omega_k-(-1)^kae^{-a})/(a^2+\omega_k^2)$.
Absolute convergence follows from
$|t_k(a)|\leq\sqrt2a(\omega_k+1)/\omega_k^2$.
\end{proof}

\subsection{A difference of two positive semidefinite kernels}
Put $u=\eta$ and $v=t_0/\sqrt{\omega_0}$, two functions on $(0,\infty)$; the
letters $u,v$ are used for them in this section only. Fix $\tau\in(-1,0)$
and define
\[
 u_\tau=\frac{u+\tau v}{\sqrt{1-\tau^2}},\qquad v_\tau=\frac{\tau u+v}{\sqrt{1-\tau^2}} .
\]
Then $u(a)u(b)-v(a)v(b)$ equals $u_\tau(a)u_\tau(b)-v_\tau(a)v_\tau(b)$ for all
$a,b$, and Lemma~\ref{lem:expansion} gives
$\mathsf m=\mathsf m^{\mathrm I}-\mathsf m^{\mathrm{II}}$ with
\[
 \mathsf m^{\mathrm I}(a,b)=u_\tau(a)u_\tau(b)+\sum_{k\ \mathrm{odd}}\frac{t_k(a)t_k(b)}{\omega_k},
 \qquad
 \mathsf m^{\mathrm{II}}(a,b)=v_\tau(a)v_\tau(b)+\sum_{\substack{k\ \mathrm{even}\\k\geq2}}
 \frac{t_k(a)t_k(b)}{\omega_k}.
\]
Both are sums, with nonnegative weights, of kernels of the form
$f(a)f(b)$. Hence for any points $\ell_1,\dots,\ell_r$ the matrices
$\mathsf M^{\mathrm I}=[\mathsf m^{\mathrm I}(\ell_i,\ell_k)]$ and
$\mathsf M^{\mathrm{II}}=[\mathsf m^{\mathrm{II}}(\ell_i,\ell_k)]$ are positive
semidefinite. They are not the positive and negative parts of $\mathsf M$:
their product need not vanish.

\begin{proposition}\label{prop:schurbound}
Let $\kappa_\tau=\sup_{a>0}\mathsf m^{\mathrm{II}}(a,a)/\mathsf m(a,a)$. Then
$\tr J_L(E)_-\leq\kappa_\tau\tr J_L(E)$ for all $L>0$ and $E\geq0$.
\end{proposition}

\begin{proof}
By \eqref{eq:schurform},
$J_L(E)=S\circ\mathsf M^{\mathrm I}-S\circ\mathsf M^{\mathrm{II}}$, and both
terms are positive semidefinite by the Schur product theorem
\cite[Section~7.5]{hj13}. Hence $J_L(E)\geq-S\circ\mathsf M^{\mathrm{II}}$,
and since $\tr M_-=\max_{0\leq\Pi\leq\id}(-\tr(M\Pi))$ for Hermitian $M$,
\[
 \tr J_L(E)_-\leq\tr(S\circ\mathsf M^{\mathrm{II}})
 =\sum_i\mathsf m^{\mathrm{II}}(\ell_i,\ell_i)S_{ii}
 \leq\kappa_\tau\sum_i\mathsf m(\ell_i,\ell_i)S_{ii}=\kappa_\tau\tr J_L(E).\qedhere
\]
\end{proof}

\subsection{The inequality in one variable}
Fix $\tau=-\frac9{25}$ and write, for $a>0$,
\[
 \mathsf n(a)=\mathsf m^{\mathrm{II}}(a,a)=\frac{(\tau u(a)+v(a))^2}{1-\tau^2}
 +\sum_{j\geq1}\frac{t_{2j}(a)^2}{\omega_{2j}},
 \qquad\mathsf m(a)=\mathsf m(a,a)=1-(1+a)e^{-a}.
\]

\begin{lemma}\label{lem:onevariable}
$\mathsf n(a)\leq\frac25\,\mathsf m(a)$ for all $a>0$.
\end{lemma}

Theorem~\ref{thm:kernel} follows from Proposition~\ref{prop:schurbound} and
Lemma~\ref{lem:onevariable}; the second inequality of the theorem is the
first one rewritten with $\tr J=\tr J_+-\tr J_-$.

Two elementary bounds are used throughout the proof. For even $k$,
$0<ae^{-a}\leq e^{-1}<\omega_0\leq\omega_k$, so
$0<\omega_k-ae^{-a}\leq\omega_k$ and
\begin{equation}\label{eq:gbound}
 \frac{t_k(a)^2}{\omega_k}\leq g_a(\omega_k),\qquad
 g_a(\omega)=\frac{2a^2\omega}{(a^2+\omega^2)^2}.
\end{equation}
The function $g_a$ increases on $[0,a/\sqrt3]$ and decreases on
$[a/\sqrt3,\infty)$; its maximum is $9/(8\sqrt3\,a)$ and
$\int_0^\infty g_a=1$. The even frequencies
$\omega_{2j}=2\pi j+\frac\pi2$ have spacing $2\pi$. Secondly,
\begin{equation}\label{eq:vbound}
 0<v(a)\leq\sqrt{\frac2{\omega_0}}\,\frac{a\,\omega_0}{a^2+\omega_0^2}
 \leq\min\Bigl(c_*a,\ \frac{\sqrt{2\omega_0}}a\Bigr),\qquad
 c_*=\sqrt2\,\omega_0^{-3/2}=0.718348\ldots
\end{equation}

\begin{lemma}\label{lem:unimodal}
Let $g\geq0$ be integrable on $[0,\infty)$, nondecreasing on $[0,c]$ and
nonincreasing on $[c,\infty)$, and let $p_j=p_0+j\epsilon$ for $j\geq1$, with
$\epsilon>0$ and $p_1\geq0$. Then
$\sum_{j\geq1}g(p_j)\leq\epsilon^{-1}\int_{p_1}^\infty g+\max g$.
\end{lemma}

\begin{proof}
For $x_1\leq x\leq x_2$ one has $g(x)\geq\min(g(x_1),g(x_2))$. The values
$\gamma_j=g(p_j)$ are nonincreasing once $p_j\geq c$, so some index $j^*$
maximizes $\gamma_j$. If $j<j^*$, then $g\geq\gamma_j$ on
$[p_j,p_{j^*}]\supset[p_j,p_{j+1}]$, so
$\gamma_j\leq\epsilon^{-1}\int_{p_j}^{p_{j+1}}g$. If $j>j^*$, then $g\geq\gamma_j$ on
$[p_{j^*},p_j]\supset[p_{j-1},p_j]$, so
$\gamma_j\leq\epsilon^{-1}\int_{p_{j-1}}^{p_j}g$. These intervals have disjoint
interiors and lie in $[p_1,\infty)$. Add $\gamma_{j^*}\leq\max g$.
\end{proof}

\begin{proof}[Proof of Lemma~\ref{lem:onevariable}]
\emph{The range $a\geq10$.} Since $0<u<1$ and $\tau<0$ we have
$\tau u+v\geq\tau$. By \eqref{eq:vbound}, $v\leq\sqrt{2\omega_0}/10<0.178$,
while $\tau u\leq\tau(1-e^{-10})<-0.359$; hence $\tau\leq\tau u+v<0$ and the
first term of $\mathsf n$ is at most $\tau^2/(1-\tau^2)=81/544<0.148898$. By
\eqref{eq:gbound} and Lemma~\ref{lem:unimodal} with $g=g_a$,
$p_j=\omega_{2j}$, $\epsilon=2\pi$,
\[
 \sum_{j\geq1}\frac{t_{2j}(a)^2}{\omega_{2j}}
 \leq\frac1{2\pi}+\frac9{8\sqrt3\,a}\leq0.159155+0.064952 .
\]
Since $\mathsf m$ is increasing, $\mathsf m(a)\geq1-11e^{-10}$, and
$\mathsf n(a)/\mathsf m(a)\leq(0.148898+0.224107)/(1-11e^{-10})<0.3733$.

\emph{The range $0<a\leq\frac1{20}$.} Here $a-\frac{a^2}2\leq u\leq a$. With
$\tau<0$ and \eqref{eq:vbound},
$\tau u+v\leq a(c_*+\tau+\frac{|\tau|}2a)\leq0.367349\,a$. Also
$\tau u+v>0$: indeed
$v\geq c_*a(1-a/\omega_0)(1-a^2/\omega_0^2)>0.68\,a$ while
$\tau u\geq-0.36\,a$. Hence the first term of $\mathsf n$ is at most
$0.367349^2a^2/(1-\tau^2)<0.155039\,a^2$. By \eqref{eq:gbound},
$t_k^2/\omega_k\leq2a^2/\omega_k^3$, so the sum is at most $2a^2\Sigma_3$ with
$\Sigma_3=\sum_{j\geq1}\omega_{2j}^{-3}<0.0026764$. The series
$\mathsf m(a)=\sum_{n\geq2}(-1)^n\frac{n-1}{n!}a^n$ is alternating with
decreasing terms for $0<a\leq1$, so
$\mathsf m(a)\geq\frac{a^2}2-\frac{a^3}3\geq\frac{29}{60}a^2$. Therefore
$\mathsf n(a)/\mathsf m(a)\leq(0.155039+0.005353)/0.483333<0.3319$.

\emph{The range $\frac1{20}\leq a\leq10$.} For an integer $j_*\geq1$ with
$\omega_{2j_*}>a/\sqrt3$, \eqref{eq:gbound} and the monotonicity of $g_a$ on
$[\omega_{2j_*},\infty)$ give
\begin{equation}\label{eq:tail}
 \sum_{j>j_*}\frac{t_{2j}(a)^2}{\omega_{2j}}
 \leq\frac1{2\pi}\int_{\omega_{2j_*}}^\infty g_a
 =\frac{a^2}{2\pi(a^2+\omega_{2j_*}^2)} .
\end{equation}
With $j_*=80$ the quantity
$\frac25\mathsf m(a)-(\tau u+v)^2/(1-\tau^2)-\sum_{j\leq80}t_{2j}^2/\omega_{2j}
-a^2/(2\pi(a^2+\omega_{160}^2))$ is an explicit elementary function of $a$.
Its positivity on $[\frac1{20},10]$ was verified with interval arithmetic:
the interval is covered by $3535$ subintervals $[a_k,a_{k+1}]$ with
$a_{k+1}/a_k\leq1.0015$, and on each of them an enclosure of the range of
this function, computed with $120$ bits of precision, is positive. The
smallest lower bound obtained for $\frac25-\mathsf n(a)/\mathsf m(a)$ is
$0.0766$, at the left end of the interval.\footnote{The enclosures were
computed with the ball arithmetic of Arb \cite{joh17}. The program and its
output accompany this paper as supplementary files
(\texttt{lemma\_3\_6\_interval\_check.py}).}
\end{proof}

\begin{remark}\label{rem:constant}
The method gives a constant close to $0.31$: evaluating the series on a grid
of values of $a$ gives $\sup_a\mathsf n(a)/\mathsf m(a)\approx0.3081$, with
limits
$\tau^2/(1-\tau^2)+\frac1{2\pi}$ as $a\to\infty$ and
$2[(c_*+\tau)^2/(1-\tau^2)+2\Sigma_3]\approx0.3058$ as $a\to0$. The value $\frac25$ is
chosen for a short verification. The best constant $\kappa_*$ in
Theorem~\ref{thm:kernel} satisfies
\[
 \frac{\sqrt2-1}2\leq\kappa_*\leq\frac25 .
\]
For the lower bound take $L=\diag(a,b)$ and $S=ww^\dagger$ with
$w=(a^{-1},2^{-1/2})$: by Lemma~\ref{lem:mformula},
$J_L(E)\to\left(\begin{smallmatrix}1/2&2^{-1/2}\\2^{-1/2}&1/2\end{smallmatrix}\right)$
as $a\to0$ and $b\to\infty$, with eigenvalues $\frac12\pm2^{-1/2}$. A numerical
maximization of $\tr J_L(E)_-/\tr J_L(E)$ over $L$ and $E$ for $r\leq12$, which
is not used in any proof, found no value above the lower bound; this
suggests that the lower bound is the best constant and that it is not
attained. We do not know a proof.
\end{remark}

\begin{remark}\label{rem:indefinite}
The matrices $\mathsf M=[\mathsf m(\ell_i,\ell_k)]$ are not positive
semidefinite, and $\mathsf m$ has infinitely many negative squares: for every
$n$ there are points $\ell_1,\dots,\ell_N$ for which $\mathsf M$ has at least
$n$ negative eigenvalues. Indeed $\mathsf M$ is the Gram matrix of the
quadratic form
$f\mapsto\inner f{(\mathbf 1\otimes\mathbf 1-\mathsf T)f}$ on the functions
$\ell_ie_{\ell_i}$. The functions $\sqrt2\cos(\omega_kx)$ are eigenfunctions
of $\mathbf 1\otimes\mathbf 1-\mathsf T$ with eigenvalues $(-1)^k/\omega_k$,
so the
form is negative definite on a space of dimension $n$; and the span of the
$e_a$, $a>0$, is dense in $L^2[0,1]$ by the Stone--Weierstrass theorem. The
inequality of Theorem~\ref{thm:kernel} also fails for indefinite $L$.
\end{remark}

\section{Trace monotonicity of the polar phase}\label{sec:phase}

This section is independent of the rest of the paper. Let $\Omega$ be a
complex $n\times n$ matrix with positive definite imaginary part
$\Ima\Omega=(\Omega-\Omega^\dagger)/2\ii$, where ${}^\dagger$ is the conjugate
transpose. Then $\Omega$ is invertible, and the unitary factor $\mathsf W$ in
the polar decompositions
$\Omega=(\Omega\Omega^\dagger)^{1/2}\mathsf W=\mathsf W(\Omega^\dagger\Omega)^{1/2}$
has its spectrum in the open upper half of the unit circle. The \emph{polar
phase} of $\Omega$ is the Hermitian matrix $\Phi$ with spectrum in $(0,\pi)$
and $\mathsf W=e^{\ii\Phi}$, and $|\Omega|=(\Omega\Omega^\dagger)^{1/2}=\Omega e^{-\ii\Phi}$
is the left modulus. The set of matrices with a prescribed positive definite
imaginary part is an affine space over the Hermitian matrices, and on it the
polar phase is a real-analytic map: the unitary polar factor of an invertible
matrix depends analytically on it, and $\Phi$ is $-\ii$ times its principal
logarithm.

\begin{theorem}[Monotonicity of the polar phase]\label{thm:phase}
Let $\Omega',\Omega''$ be matrices with the same positive definite imaginary
part, and let $\Phi',\Phi''$ be their polar phases. Then
\[
 \tr\bigl((\Omega''-\Omega')(\Phi''-\Phi')\bigr)\leq0,
\]
with equality only if $\Omega'=\Omega''$. More precisely, for every smooth
family $\Omega(\varepsilon)$ with constant positive definite imaginary part,
$\tr(\dot\Omega\,\dot\Phi)<0$ wherever $\dot\Omega\neq0$.
\end{theorem}

In other words, $\Rea\Omega\mapsto-\Phi$ is a strictly monotone operator
\cite{min62} on the real Hilbert space of Hermitian matrices with the trace
inner product.
For $n=1$ this says that $\arg(x+\ii y)$ decreases in $x$ when $y>0$, and for
a scalar imaginary part it follows from that by the functional calculus. In
general $\Omega$ is not normal and $\Phi$ is not a function of its
eigenvalues. The derivative of the unitary polar factor solves a Sylvester
equation, which is used in the perturbation theory of the polar
decomposition \cite{bha97}; we are not aware of a statement on the sign of the derivative
of its argument. The phases of a sectorial matrix in the sense of
\cite{wckq20} are defined through a congruence and are a different object.

\begin{corollary}\label{cor:comparison}
Let $\Ima\Omega^{(0)}>0$, let $K$ be Hermitian, and let $\Phi^{(0)}$ and $\Phi$
be the polar phases of $\Omega^{(0)}$ and $\Omega^{(0)}+K$. Then
$\tr(K\Phi)\leq\tr(K\Phi^{(0)})$. In particular, if
$\Omega^{(0)}=zP$ with $\Ima z>0$ and $P>0$, then
$\tr(K\Phi)\leq(\arg z)\tr K$.
\end{corollary}

\begin{proof}
Theorem~\ref{thm:phase} with $\Omega'=\Omega^{(0)}$ and
$\Omega''=\Omega^{(0)}+K$. The polar phase of $zP=(|z|P)e^{\ii\arg z}$ is the
scalar $\arg z$.
\end{proof}

\begin{corollary}\label{cor:blocks}
Let $\CC^n=H\oplus H'$ be an orthogonal decomposition and
\[
 \Omega=\begin{pmatrix}\Omega_{11}&K_{12}\\K_{12}^\dagger&\Omega_{22}\end{pmatrix},
 \qquad\Ima\Omega>0\text{ block diagonal},
\]
with polar phase $\Phi=(\Phi_{ik})$. Then
$\Rea\tr(K_{12}\,\Phi_{21})\leq0$, with equality only if $K_{12}=0$.
\end{corollary}

\begin{proof}
Let $K$ be the Hermitian matrix with off-diagonal blocks $K_{12},K_{12}^\dagger$
and zero diagonal blocks, $\Omega(\varepsilon)=\Omega-(1-\varepsilon)K$ and
$\Sigma=\diag(\id,-\id)$. Then $\Ima\Omega(\varepsilon)=\Ima\Omega$,
$\Sigma\Omega(\varepsilon)\Sigma=\Omega(-\varepsilon)$, and the polar phase is
equivariant under unitary conjugation; so
$f(\varepsilon)=\tr(K\Phi(\varepsilon))$ is odd, $f(0)=0$, and $f'<0$ for
$K\neq0$ by Theorem~\ref{thm:phase}. Finally
$f(1)=2\Rea\tr(K_{12}\Phi_{21})$.
\end{proof}

The rest of this section proves Theorem~\ref{thm:phase}. The pairing
$\tr(\dot\Omega\dot\Phi)$ is expressed through the solution of a Sylvester
equation, and the resulting quadratic form is a sum of forms with matrices
$(\Ima\Omega)\circ\mathsf L^{(a)}$, where the matrices $\mathsf L^{(a)}$ are
positive semidefinite by an integral representation of the kind that
characterizes anti-L\"owner functions \cite{aud11}.

\subsection{The first variation of the phase}
Put
\[
 \tilde P=e^{-\ii\Phi/2}\,|\Omega|\,e^{\ii\Phi/2}>0,\qquad\text{so that}\qquad
 \Omega=e^{\ii\Phi/2}\,\tilde P\,e^{\ii\Phi/2}.
\]
We work in an orthonormal eigenbasis of $\Phi$, $\Phi=\diag(\phi_a)$ with
$\phi_a\in(0,\pi)$. Then
$\Omega_{ab}=e^{\ii(\phi_a+\phi_b)/2}\tilde P_{ab}$, and since $\tilde P$ is
Hermitian,
\begin{equation}\label{eq:reim}
 (\Rea\Omega)_{ab}=\tilde P_{ab}\cos\tfrac{\phi_a+\phi_b}2,\qquad
 (\Ima\Omega)_{ab}=\tilde P_{ab}\,\varsigma_{ab},\qquad
 \varsigma_{ab}:=\sin\tfrac{\phi_a+\phi_b}2>0 .
\end{equation}

\begin{lemma}\label{lem:variation}
Let $\Omega(\varepsilon)$ be a smooth family with $\Omega(0)=\Omega$,
$\Ima\Omega(\varepsilon)>0$, and derivative $\dot\Omega$ at $0$. There is a
Hermitian matrix $\Theta$ with
\begin{equation}\label{eq:sylvester}
 \tilde P\Theta+\Theta\tilde P=\frac{X-X^\dagger}{2\ii},\qquad
 X=e^{-\ii\Phi/2}\,\dot\Omega\,e^{-\ii\Phi/2},
\end{equation}
and, in the eigenbasis of $\Phi$,
\begin{equation}\label{eq:phidot}
 \dot\Phi_{ab}=\frac{\phi_a-\phi_b}{\sin\frac{\phi_a-\phi_b}2}\,\Theta_{ab},
\end{equation}
where the factor is read as $2$ if $\phi_a=\phi_b$. If $\dot\Omega$ is
Hermitian, the right side of \eqref{eq:sylvester} has entries
$-\dot\Omega_{ab}\,\varsigma_{ab}$.
\end{lemma}

\begin{proof}
Let $\mathsf U(\varepsilon)=e^{\ii\Phi(\varepsilon)/2}$, so that
$\Omega=\mathsf U\tilde P\mathsf U$ for every $\varepsilon$. Since
$\mathsf U$ is unitary,
$\dot{\mathsf U}=\ii\Upsilon\mathsf U=\mathsf U\,\ii\Upsilon'$ with $\Upsilon$ Hermitian and
$\Upsilon'=\mathsf U^{-1}\Upsilon\mathsf U$. Hence
$\dot\Omega=\ii\Upsilon\Omega+\mathsf U\dot{\tilde P}\mathsf U+\ii\Omega\Upsilon'$ and
\[
 \mathsf U^{-1}\dot\Omega\,\mathsf U^{-1}
 =\dot{\tilde P}+\ii\bigl(\Upsilon'\tilde P+\tilde P\Upsilon\bigr).
\]
Apply $M\mapsto(M-M^\dagger)/2\ii$. The Hermitian matrix $\dot{\tilde P}$
drops out and $\ii(\Upsilon'\tilde P+\tilde P\Upsilon)$ contributes
$\frac12[(\Upsilon+\Upsilon')\tilde P+\tilde P(\Upsilon+\Upsilon')]$; so \eqref{eq:sylvester}
holds with $\Theta=\frac12(\Upsilon+\Upsilon')$. For an entire function $f$ and
a smooth family of Hermitian matrices, differentiating the Cauchy integral
$f(\Phi)=\frac1{2\pi\ii}\oint f(\zeta)(\zeta-\Phi)^{-1}d\zeta$ gives, in an
eigenbasis of $\Phi$,
$(\frac d{d\varepsilon}f(\Phi))_{ab}=\dot\Phi_{ab}\,(f(\phi_a)-f(\phi_b))/(\phi_a-\phi_b)$,
with the value $f'(\phi_a)$ if $\phi_a=\phi_b$ (the formula of Daleckii and
Krein; see \cite{bha97}). For $f(x)=e^{\ii x/2}$ and
$\gamma=\phi_a-\phi_b$ this gives
$\Upsilon_{ab}=-\ii\dot{\mathsf U}_{ab}e^{-\ii\phi_b/2}
=\dot\Phi_{ab}\,2\gamma^{-1}\sin(\gamma/4)\,e^{\ii\gamma/4}$ and
$\Upsilon'_{ab}=e^{-\ii\gamma/2}\Upsilon_{ab}$, hence
$\Theta_{ab}=\dot\Phi_{ab}\sin(\gamma/2)/\gamma$. For Hermitian
$\dot\Omega$,
$(X-X^\dagger)_{ab}/2\ii=\dot\Omega_{ab}(e^{-\ii(\phi_a+\phi_b)/2}-e^{\ii(\phi_a+\phi_b)/2})/2\ii
=-\dot\Omega_{ab}\,\varsigma_{ab}$.
\end{proof}

The equation \eqref{eq:sylvester} has a unique solution because
$\tilde P>0$.

\begin{lemma}\label{lem:form}
Let
\[
 \varrho_{ab}=\frac{\phi_a-\phi_b}{2(\cos\phi_b-\cos\phi_a)}\quad(\phi_a\neq\phi_b),
 \qquad\varrho_{ab}=\frac1{2\sin\phi_a}\quad(\phi_a=\phi_b),
\]
a positive symmetric function of $(\phi_a,\phi_b)$. For a family with
constant imaginary part, $\tr(\dot\Omega\dot\Phi)=-4\,Q(\Theta)$, where
\[
 Q(\Theta)=\sum_{a,b}\varrho_{ab}\,(\tilde P\Theta+\Theta\tilde P)_{ba}\,\Theta_{ab}
 =\sum_a\inner{\theta^a}{\bigl((\Ima\Omega)\circ\mathsf L^{(a)}\bigr)\theta^a},\qquad
 \mathsf L^{(a)}_{bd}=\frac{\varrho_{ba}+\varrho_{da}}{\varsigma_{bd}},
\]
and $\theta^a$ is the $a$-th column of $\Theta$.
\end{lemma}

\begin{proof}
By Lemma~\ref{lem:variation},
$\dot\Omega_{ab}=-(\tilde P\Theta+\Theta\tilde P)_{ab}/\varsigma_{ab}$ and
$\dot\Phi_{ab}=\gamma\,\Theta_{ab}/\sin(\gamma/2)$ with
$\gamma=\phi_a-\phi_b$, and
$\gamma/(\sin(\gamma/2)\,\varsigma_{ab})=2\gamma/(\cos\phi_b-\cos\phi_a)=4\varrho_{ab}$.
So
$\tr(\dot\Omega\dot\Phi)=\sum_{a,b}\overline{\dot\Omega_{ab}}\dot\Phi_{ab}=-4Q(\Theta)$,
the conjugate of an entry of a Hermitian matrix being the transposed entry.
Expanding,
$Q=\sum_{a,b,c}\varrho_{ab}[\tilde P_{bc}\Theta_{ca}\Theta_{ab}
+\Theta_{bc}\tilde P_{ca}\Theta_{ab}]$; renaming the indices $(a,b,c)$ of the
second sum as $(c,a,b)$ turns it into
$\sum_{a,b,c}\varrho_{ca}\tilde P_{bc}\Theta_{ca}\Theta_{ab}$. Hence
$Q=\sum_{a,b,c}(\varrho_{ab}+\varrho_{ca})\tilde P_{bc}\,\Theta_{ab}\Theta_{ca}$,
and $\tilde P_{bc}=(\Ima\Omega)_{bc}/\varsigma_{bc}$ by \eqref{eq:reim}, while
$\Theta_{ab}\Theta_{ca}=\overline{\theta^a_b}\,\theta^a_c$.
\end{proof}

\subsection{Positivity of the kernels}

\begin{lemma}\label{lem:integral}
For $x,c>0$ let
$g_c(x)=(1+x^2)(\arctan x-\arctan c)/(x^2-c^2)$ for $x\neq c$ and
$g_c(c)=1/(2c)$. Then
\[
 g_c(x)=\frac2\pi\int_0^\infty\frac x{x^2+y^2}\,w_c(y)\,dy,\qquad
 w_c(y)=\frac{|1-y^2|\,k_c(y)}{y^2+c^2}\geq0,
\]
where $k_c(y)=\arctan c$ for $0<y<1$ and $k_c(y)=\frac\pi2-\arctan c$ for
$y>1$.
\end{lemma}

\begin{proof}
For $x\neq c$,
$(1-y^2)/((x^2+y^2)(c^2+y^2))=p/(x^2+y^2)+p'/(c^2+y^2)$ with
$p=(1+x^2)/(c^2-x^2)$ and $p'=(1+c^2)/(x^2-c^2)$. Moreover
$\int_0^1x\,dy/(x^2+y^2)=\frac\pi2-\arctan x$,
$\int_1^\infty x\,dy/(x^2+y^2)=\arctan x$,
$\int_0^1x\,dy/(c^2+y^2)=\frac xc(\frac\pi2-\arctan c)$ and
$\int_1^\infty x\,dy/(c^2+y^2)=\frac xc\arctan c$. With $\beta=\arctan c$ and
$\beta'=\frac\pi2-\beta$,
\begin{align*}
 \int_0^\infty\frac{x\,w_c(y)}{x^2+y^2}\,dy
 &=\beta\bigl[p(\tfrac\pi2-\arctan x)+p'\tfrac xc\beta'\bigr]
 -\beta'\bigl[p\arctan x+p'\tfrac xc\beta\bigr]\\
 &=\tfrac\pi2\,p\,(\arctan c-\arctan x)=\tfrac\pi2\,g_c(x).
\end{align*}
Both sides are continuous in $x$, which gives the case $x=c$.
\end{proof}

\begin{lemma}\label{lem:antiloewner}
For every $c>0$ and all $x_1,\dots,x_n>0$ the matrix
$\bigl[(g_c(x_b)+g_c(x_d))/(x_b+x_d)\bigr]_{b,d}$ is positive semidefinite.
\end{lemma}

\begin{proof}
For $f_y(x)=x/(x^2+y^2)$,
\[
 \frac{f_y(x)+f_y(x')}{x+x'}=\frac{xx'+y^2}{(x^2+y^2)(x'^2+y^2)}
 =f_y(x)f_y(x')+y^2\,\frac1{x^2+y^2}\,\frac1{x'^2+y^2},
\]
a sum of two kernels of the form $\psi(x)\psi(x')$. Integrate against the
positive measure $\frac2\pi w_c(y)\,dy$ of Lemma~\ref{lem:integral}.
\end{proof}

Functions $g$ on $(0,\infty)$ for which all matrices
$[(g(x_b)+g(x_d))/(x_b+x_d)]$ are positive semidefinite are called
anti-L\"owner; Audenaert \cite{aud11} characterizes them by an integral
representation, and Lemma~\ref{lem:integral} puts $g_c$ in that form. The
sufficient condition of Kwong quoted in \cite{aud11}, operator monotonicity
of $g$, does not apply, since $g_c$ is not monotone: for $c=1$ the limits at
$0$ and at $\infty$ are $\frac\pi4$, while the value at $1$ is $\frac12$.

\begin{lemma}\label{lem:antiloewnermatrix}
For every index $a$ the matrix $\mathsf L^{(a)}$ of Lemma~\ref{lem:form} is
positive semidefinite, with positive diagonal.
\end{lemma}

\begin{proof}
Put $x_b=\tan(\phi_b/2)>0$ and $c=x_a$. Then
$\cos\phi_a-\cos\phi_b=2(x_b^2-c^2)/((1+c^2)(1+x_b^2))$ and
$\phi_b-\phi_a=2(\arctan x_b-\arctan c)$, so
$\varrho_{ba}=\frac{1+c^2}2g_c(x_b)$, also when $\phi_b=\phi_a$; and
$\varsigma_{bd}=(x_b+x_d)/(\sqrt{1+x_b^2}\sqrt{1+x_d^2})$. Hence
\[
 \mathsf L^{(a)}_{bd}=\frac{1+c^2}2\,\sqrt{1+x_b^2}\,\sqrt{1+x_d^2}\;
 \frac{g_c(x_b)+g_c(x_d)}{x_b+x_d},
\]
the matrix of Lemma~\ref{lem:antiloewner} conjugated by a positive diagonal
matrix. The diagonal entries $2\varrho_{ba}/\varsigma_{bb}$ are positive.
\end{proof}

\begin{proof}[Proof of Theorem~\ref{thm:phase}]
By Lemma~\ref{lem:form},
$\tr(\dot\Omega\dot\Phi)=-4\sum_a\inner{\theta^a}{((\Ima\Omega)\circ\mathsf L^{(a)})\theta^a}$.
Let $y_{\min}>0$ be the least eigenvalue of $\Ima\Omega$. Since
$\Ima\Omega-y_{\min}\id\geq0$ and $\mathsf L^{(a)}\geq0$, the Schur product
theorem \cite[Section~7.5]{hj13} gives
$(\Ima\Omega)\circ\mathsf L^{(a)}\geq y_{\min}\diag(\mathsf L^{(a)}_{bb})>0$ by
Lemma~\ref{lem:antiloewnermatrix}. So $Q(\Theta)>0$ unless $\Theta=0$, and
$\Theta=0$ forces $\dot\Omega=0$ by \eqref{eq:sylvester}. The integrated
statement follows along the segment
$\Omega'+\varepsilon(\Omega''-\Omega')$, on which the imaginary part is
constant.
\end{proof}

\begin{remark}\label{rem:phasefails}
The positivity of $\Ima\Omega$ is needed. (The numbers of this remark are
computed in floating-point arithmetic and are not used in any proof.) The
polar phase is defined and
analytic near any invertible matrix whose unitary polar factor has spectrum
in the open upper half-circle. For $\Phi=\diag(\frac15,3)$,
$\tilde P=\left(\begin{smallmatrix}1&9/10\\9/10&1\end{smallmatrix}\right)$ and
$\Omega=e^{\ii\Phi/2}\tilde Pe^{\ii\Phi/2}$, the matrix $\Ima\Omega$ has
eigenvalues $-0.730\ldots$ and $1.0699\ldots$, and the symmetric part of the
derivative of the phase on Hermitian variations has the positive eigenvalue
$0.2335\ldots$ For $\Ima\Omega>0$ the phase is in general not the gradient of
a function of $\Rea\Omega$ for the trace inner product: at
$\Omega=\left(\begin{smallmatrix}\ii&1\\1&3\ii\end{smallmatrix}\right)$ the
derivative of the entry $\Phi_{11}$ with respect to $(\Rea\Omega)_{22}$ is
$0.0488\ldots$, while that of $\Phi_{22}$ with respect to $(\Rea\Omega)_{11}$
is $0.0306\ldots$ Nor is $\Rea\Omega\mapsto-\Phi$ monotone for the L\"owner
order, even for a scalar imaginary part $y$: there
$-\Phi=\arctan(\Rea\Omega/y)-\frac\pi2$, and the arctangent is not operator
monotone on the real line, since an operator monotone function on the real
line is affine. The
trace norm $\tr|\Omega|$ is a convex function of $\Rea\Omega$ with gradient
$\cos\Phi$, so $\cos\Phi$ is a monotone map; Theorem~\ref{thm:phase} concerns
$\Phi$ itself and does not follow from this.
\end{remark}

\section{The arrow energy and the cancellation of the acyclic part}\label{sec:energy}

In this section the charges are arbitrary in the upper half-plane, and no
hypothesis on the cohomology representation is made. The subscripts of
$z_i$, $x_i$, $y_i$, $\alpha_i$, $n_i$, $C_i$, $\phi_i$, $\delta_i$ and $M_1$
refer to the vertices. In $\lambda_0$ and
$\Xi_0$ the subscript marks the value at time zero, and in $N_0$, $t_0$,
$t_1$, $t_2$, $\kappa_0$, $\kappa_1$, $d_0$, $s_1$, $\Lambda_1$ and
$\mathfrak c_1$, $\mathfrak c_2$, $\mathfrak c_3$ it is a label.

\subsection{The arrow energy}

\begin{definition}\label{def:energy}
The \emph{arrow energy} of a state is
\[
 N:=\sum_j\tr\bigl(A_jR(D^\dagger D)A_j^\dagger\bigr)=\tr\bigl(CR(B)\bigr)=\tr G .
\]
\end{definition}

It is the term of the primitive \eqref{eq:primitive} that does not contain
the charges. For $D=0$ it would be the sum of the squared norms of the
arrows.

\begin{lemma}\label{lem:energy}
The arrow energy has the following properties.
\begin{enumerate}
\item[(a)] $N=q+\tr Y$.
\item[(b)] $\tr K=\tr(Ce^{-B})\geq0$,
 $\tr J=\tr\bigl(C(R(B)-e^{-B})\bigr)\geq0$, and $\tr K+\tr J=N$.
\item[(c)] Along the phase flow, for all charges,
 \[
 \frac{dN}{dt}=\tr(K\Phi_V)+\tr(J\Phi_U)-\tr(G\Phi_W).
 \]
\end{enumerate}
\end{lemma}

\begin{proof}
(a) In the adapted splitting $R(B)=\id_{\ker D}\oplus R(F^\dagger F)$, so
$\tr(CR(B))=\tr C_0+\sum_j\tr(T_jR(F^\dagger F)T_j^\dagger)$, and
$T_jR(F^\dagger F)T_j^\dagger=X_jFR(F^\dagger F)F^\dagger X_j^\dagger
=X_j(\id-e^{-L})X_j^\dagger$.

(b) $\tr K_B(C)=\int_0^1\tr(e^{-(1-x)B}Ce^{-xB})\,dx=\tr(Ce^{-B})$. By
\eqref{eq:J}, $\tr J_L(M)=\tr(M\,g(L))$ with
$g(\ell)=\int_{x,y\geq0,\,x+y\leq1}e^{-(x+y)\ell}\,dx\,dy
=\int_0^1xe^{-x\ell}dx=(1-e^{-\ell}-\ell e^{-\ell})/\ell^2$; for
$M=DCD^\dagger$, and since $D^\dagger g(L)D=B\,g(B)=R(B)-e^{-B}$ (both sides
vanish on $\ker B$), this is $\tr(C(R(B)-e^{-B}))$. Both traces are
nonnegative because $C\geq0$, $e^{-B}\geq0$ and $R(b)\geq e^{-b}$ for
$b\geq0$.

(c) This is \eqref{eq:energyvariation} at $\xi=-\Phi$.
\end{proof}

\begin{lemma}\label{lem:three}
At every state and for all charges in the upper half-plane:
\begin{enumerate}
\item[(a)] $\tr(K\Phi_V)\leq\alpha_0\tr K$;
\item[(b)] $\tr(J\Phi_U)\leq\pi\tr J_-\leq\frac{2\pi}5\tr J$, and
 $\tr(J\Phi_U)\leq0$ if $r=1$;
\item[(c)] $\tr(G\Phi_W)\geq\pi N-n_1M_1$, where
 $M_1:=\max\{\pi\max(x_1,0),\,2y_1\}$.
\end{enumerate}
\end{lemma}

\begin{proof}
(a) Corollary~\ref{cor:comparison} with $\Omega^{(0)}=z_0e^{-B}$ and the
Hermitian matrix $K$: $\Omega^{(0)}+K=\Omega_V$, and $e^{-B}>0$.

(b) Let $(u_k)$ be an orthonormal eigenbasis of $\Phi_U$ with eigenvalues
$\phi_k\in(-\pi,0)$, and $j_k=\inner{u_k}{Ju_k}$. Then
$\tr(J\Phi_U)=\sum_k\phi_kj_k\leq\pi\sum_k(j_k)_-$, where
$(j_k)_-=\max(-j_k,0)$, and
$\sum_k(j_k)_-\leq\sum_k\inner{u_k}{J_-u_k}=\tr J_-$. By
Theorem~\ref{thm:kernel}, $\tr J_-\leq\frac25\tr J$, because
$DCD^\dagger\geq0$. For $r=1$, $J\geq0$ is a number and $\Phi_U<0$.

(c) $\Omega_W=z_1\id-G$ is normal and $\Phi_W=f_1(G)$ with
$f_1(x)=\arg(z_1-x)$. We claim that $x(\pi-f_1(x))\leq M_1$ for all
$x\geq0$. If $x\leq\max(x_1,0)$ the left side is at most $\pi\max(x_1,0)$.
If $x>\max(x_1,0)$, then $x-x_1>0$ and
$\pi-f_1(x)=\arctan\frac{y_1}{x-x_1}\leq\min\{\frac\pi2,\frac{y_1}{x-x_1}\}$;
for $x\leq2x_1$ the product is at most $2x_1\cdot\frac\pi2$, and for
$x>2\max(x_1,0)$ we have $x-x_1\geq x/2$ and the product is at most $2y_1$.
Summing over the eigenvalues of $G\geq0$ gives
$\tr(G\,f_1(G))\geq\pi\tr G-n_1M_1$.
\end{proof}

\begin{theorem}[A priori bound]\label{thm:energy}
Let
\[
 \rho:=\pi-\max\bigl\{\alpha_0,\tfrac{2\pi}5\bigr\}\quad
 (\rho:=\pi-\alpha_0\text{ if }r=1),\qquad
 \bar N:=\frac{2n_1M_1}\rho,\qquad
 t_0:=\frac1\rho\log^+\frac{\rho N(0)}{n_1M_1},
\]
where $\log^+x=\max(\log x,0)$.
For all charges in the upper half-plane and along every trajectory of the
phase flow,
\[
 \frac{dN}{dt}\leq-\rho N+n_1M_1,\qquad
 N(t)\leq N(0)e^{-\rho t}+\frac{n_1M_1}\rho\bigl(1-e^{-\rho t}\bigr).
\]
Hence $N(t)\leq N_0:=\max\{N(0),n_1M_1/\rho\}$ for all $t\geq0$, and
$N(t)\leq\bar N$ for $t\geq t_0$. At all times $q\leq N$, $\tr Y\leq N$,
$\tr J\leq N$ and $s\leq N/(1-e^{-\lambda})$.
\end{theorem}

\begin{proof}
By Lemmas~\ref{lem:energy}(c) and~\ref{lem:three},
$\dot N\leq\alpha_0\tr K+\frac{2\pi}5\tr J-\pi N+n_1M_1$. Since $\tr K$ and
$\tr J$ are nonnegative with sum $N$, the first two terms are at most
$\max\{\alpha_0,\frac{2\pi}5\}N$; for $r=1$ the second term may be omitted.
Integrate; for $t\geq t_0$ the first term of the bound is at most
$n_1M_1/\rho$. The remaining statements are Lemma~\ref{lem:energy}(a),(b) and
\eqref{eq:Y}.
\end{proof}

The bound $\bar N$ depends only on the charges and on $n_1$; the initial
state enters only through the time $t_0$. For the symmetric charges
$\rho=\frac\pi4$, $M_1=\pi$ and $\bar N=8n_1$. The two matrix inequalities of
Sections~\ref{sec:kernel} and~\ref{sec:phase} enter through
Lemma~\ref{lem:three}(a),(b). If Lemma~\ref{lem:three}(a) is replaced by the
elementary estimates $\tr(K\Phi_V)\leq\pi\tr K_+$,
$\tr K_+=\frac12(\norm K_1+\tr K)$ and $\norm K_1\leq\tr(CR(2B))\leq N$, one
obtains only $\dot N\leq-\frac\pi{10}\tr J+n_1M_1$, which gives no decay of
the arrows on $\ker D$.

\subsection{Growth of the least eigenvalue of
\texorpdfstring{$DD^\dagger$}{D D-dagger}}

Let $E_1(b)=\int_1^\infty e^{-xb}\,dx/x$ be the exponential integral. For
$b>0$,
\begin{equation}\label{eq:E1bounds}
 \frac{e^{-b}}{1+b}\leq E_1(b)\leq e^{-b}\log\bigl(1+\tfrac1b\bigr),\qquad
 E_1'(b)=-\frac{e^{-b}}b :
\end{equation}
the difference $E_1(b)-e^{-b}/(1+b)$ has derivative $-e^{-b}/(b(1+b)^2)<0$
and the difference $e^{-b}\log(1+\frac1b)-E_1(b)$ has derivative
$e^{-b}[\frac1{1+b}-\log(1+\frac1b)]<0$, and both tend to $0$ as $b\to\infty$.

\begin{lemma}\label{lem:growth}
For all charges, along the phase flow,
\[
 \frac d{dt}\tr E_1(L)=-\tr\bigl(e^{-F^\dagger F}\Phi_{AA}\bigr)-\tr\bigl(e^{-L}|\Phi_U|\bigr)
 \leq-\tr\bigl(e^{-L}|\Phi_U|\bigr),
\]
and at every state
\[
 \tr\bigl(e^{-L}|\Phi_U|\bigr)\geq
 \frac{y_0\,(\tr e^{-L})^2}{(y_0+|x_0|)\tr e^{-L}+\tr J_+} .
\]
\end{lemma}

\begin{proof}
By \eqref{eq:mapflow},
$\dot L=D\Phi_VD^\dagger-\frac12(\Phi_UL+L\Phi_U)$, so by \eqref{eq:E1bounds}
\[
 \frac d{dt}\tr E_1(L)=-\tr\bigl(L^{-1}e^{-L}D\Phi_VD^\dagger\bigr)+\tr\bigl(e^{-L}\Phi_U\bigr),
\]
and $D^\dagger L^{-1}e^{-L}D=0\oplus e^{-F^\dagger F}$ in the adapted
splitting. Both terms are nonpositive, because $\Phi_{AA}>0>\Phi_U$.

Let $(u_k)$ be an orthonormal eigenbasis of $\Phi_U$ with eigenvalues
$\phi_k\in(-\pi,0)$, and put $e_k=\inner{u_k}{e^{-L}u_k}>0$,
$j_k=\inner{u_k}{Ju_k}$. By Lemma~\ref{lem:eigenbasis} applied to
$-\Omega_U=z_0e^{-L}-J$, whose polar phase is $\Phi_U+\pi$,
$\inner{u_k}{(z_0e^{-L}-J)u_k}=p_ke^{\ii(\phi_k+\pi)}$ with $p_k>0$. Hence
$\cot(\phi_k+\pi)=(x_0e_k-j_k)/(y_0e_k)$ and
\[
 |\phi_k|=\pi-(\phi_k+\pi)=\operatorname{arccot}\frac{j_k-x_0e_k}{y_0e_k}
 \geq\operatorname{arccot}\frac{(j_k)_++|x_0|e_k}{y_0e_k}
 \geq\frac{y_0e_k}{(y_0+|x_0|)e_k+(j_k)_+},
\]
where $(j_k)_+=\max(j_k,0)$; we used that $\operatorname{arccot}$ is
decreasing and that $\operatorname{arccot}x\geq\frac1{1+x}$ for $x\geq0$ (the
difference is positive at $0$, tends to $0$ at infinity and has derivative
$-2x/((1+x^2)(1+x)^2)\leq0$). Therefore
\[
 \tr\bigl(e^{-L}|\Phi_U|\bigr)=\sum_k|\phi_k|e_k
 \geq\sum_k\frac{y_0e_k^2}{(y_0+|x_0|)e_k+(j_k)_+}
 \geq\frac{y_0\,(\sum_ke_k)^2}{(y_0+|x_0|)\sum_ke_k+\sum_k(j_k)_+}
\]
by the Cauchy--Schwarz inequality in the form
$\sum_ka_k^2/d_k\geq(\sum_ka_k)^2/\sum_kd_k$ for $d_k>0$. Finally
$\sum_ke_k=\tr e^{-L}$ and $\sum_k(j_k)_+\leq\tr J_+$.
\end{proof}

\begin{theorem}[Growth]\label{thm:growth}
Let $\bar N$ and $t_0$ be as in Theorem~\ref{thm:energy}, $y_0=\Ima z_0$,
$\lambda_0=\lambda(0)>0$, $c_\lambda=1/\log(1+1/\lambda_0)$,
$\Xi_0=\tr E_1(L)|_{t=0}$ and
\[
 \kappa_1:=\frac{y_0\,c_\lambda^2}{\sqrt2\,|z_0|\,c_\lambda\,\Xi_0+\frac75\bar N}.
\]
Then for all charges and all $t>t_0$
\[
 E_1(\lambda(t))\leq\tr E_1(L(t))\leq\frac1{\Xi_0^{-1}+\kappa_1(t-t_0)}<\frac1{\kappa_1(t-t_0)}.
\]
In particular $\lambda(t)\to\infty$, and $\lambda(t)\geq\Lambda$ for all
$t\geq t_0+(1+\Lambda)e^\Lambda/\kappa_1$, for every $\Lambda>0$.
\end{theorem}

\begin{proof}
Put $\Xi=\tr E_1(L)=\sum_iE_1(\ell_i)$. By
Proposition~\ref{prop:imported}(i) all $\ell_i\geq\lambda_0$ for $t\geq0$, so
\eqref{eq:E1bounds} gives $\Xi\leq\log(1+1/\lambda_0)\tr e^{-L}$, that is
$\tr e^{-L}\geq c_\lambda\Xi$. By Theorem~\ref{thm:kernel} and
Theorem~\ref{thm:energy}, for $t\geq t_0$,
$\tr J_+=\tr J+\tr J_-\leq\frac75\tr J\leq\frac75\bar N$, and
$y_0+|x_0|\leq\sqrt2|z_0|$. The function $x\mapsto x^2/(cx+d)$ is increasing
on $x>0$ for $c,d>0$. Hence by Lemma~\ref{lem:growth}, for $t\geq t_0$,
\[
 \dot\Xi\leq-\frac{y_0c_\lambda^2\,\Xi^2}{\sqrt2|z_0|c_\lambda\,\Xi+\frac75\bar N}
 \leq-\kappa_1\Xi^2,
\]
where we used $\Xi\leq\Xi_0$ in the denominator; $\Xi$ is nonincreasing by
Lemma~\ref{lem:growth}. Integrating from $t_0$,
$\Xi(t)^{-1}\geq\Xi(t_0)^{-1}+\kappa_1(t-t_0)\geq\Xi_0^{-1}+\kappa_1(t-t_0)$.
Since $E_1(\lambda)$ is one of the nonnegative summands of $\Xi$,
$E_1(\lambda)\leq\Xi$. For the last statement, if
$t\geq t_0+(1+\Lambda)e^\Lambda/\kappa_1$ then
$E_1(\lambda(t))<e^{-\Lambda}/(1+\Lambda)\leq E_1(\Lambda)$ by
\eqref{eq:E1bounds}.
\end{proof}

\subsection{Cancellation}
Put
\begin{equation}\label{eq:constants}
 \begin{gathered}
 \Lambda_1=\max\Bigl\{1,\frac{8\pi\bar N}{\alpha_1y_0}\Bigr\},\qquad
 t_1=t_0+\frac{(1+\Lambda_1)e^{\Lambda_1}}{\kappa_1},\qquad
 s_1=\frac{\bar N}{1-e^{-1}},\\
 \mathfrak c_1=2\sqrt{d_0\bar Ns_1}+r\kappa_0s_1,\qquad
 \mathfrak c_2=\frac\pi{\sqrt{|z_0|}}\bigl(\sqrt{\mathfrak c_1}+\sqrt{rs_1}\bigr),\qquad
 \mathfrak c_3=\max\Bigl\{0,2\log\frac{2\mathfrak c_2}\pi\Bigr\},\\
 C_{\mathrm{can}}=\mathfrak c_1+(r+1)s_1+2|z_0|re^{\mathfrak c_3},\qquad
 t_2=t_1+\frac2{\alpha_1}\log^+\frac{8\alpha_1s_1}{\pi y_0}.
 \end{gathered}
\end{equation}
The numbers $\Lambda_1$, $s_1$, $\mathfrak c_1$, $\mathfrak c_2$,
$\mathfrak c_3$, $C_{\mathrm{can}}$ depend only on the dimensions and the charges; the times $t_1$ and $t_2$
depend in addition on the three numbers $N(0)$, $\lambda(0)$ and
$\tr E_1(L)|_{t=0}$ of the initial state.

\begin{theorem}[Cancellation]\label{thm:cancellation}
For all charges in the upper half-plane and along every trajectory of the
phase flow:
\begin{enumerate}
\item[(i)] for $t\geq t_1$: $\lambda(t)\geq\Lambda_1$,
 $\dot s\leq-\frac{\alpha_1}2s$ and
 $s(t)\leq s_1e^{-\frac{\alpha_1}2(t-t_1)}$;
\item[(ii)] for $t\geq t_1$:
 $\lambda(t)\geq\max\{\Lambda_1,\ \frac{\alpha_1}4(t-t_1)-\mathfrak c_3\}$;
\item[(iii)] for $t\geq t_1$:
 \[
 \norm{\Omega_V-(z_0\id+C_0)\oplus0}_1+\norm{\Omega_U}_1+\norm{\Omega_W-(z_1\id-C_1)}_1
 \leq C_{\mathrm{can}}\,e^{-\frac{\alpha_1}4(t-t_1)},
 \]
 and $|m(h(t))-m_H(t)|$ satisfies the same bound;
\item[(iv)] for $t\geq t_2$:
 $\norm{\delta_0}_{\HS}^2\leq\frac{2\pi\alpha_1}{y_0}\,s$ and
 $\norm{\delta_1}_{\HS}\leq s/y_1$.
\end{enumerate}
So the acyclic part cancels, with $q\leq\bar N$ from the time $t_0$ on, and
after the time $t_2$ the metrics induced on the cohomology representation move by its own
phases $\phi_0,\phi_1$ up to an exponentially decaying error.
\end{theorem}

\begin{proof}
(i) By Theorem~\ref{thm:growth} with $\Lambda=\Lambda_1$,
$\lambda\geq\Lambda_1$ on $[t_1,\infty)$. By Theorem~\ref{thm:energy},
$q\leq\bar N$ for $t\geq t_0$, so $\frac{4\pi q}{y_0\lambda}\leq\frac{\alpha_1}2$ there, and
Proposition~\ref{prop:imported}(iii) gives $\dot s\leq-\frac{\alpha_1}2s$. At
$t_1$, $\lambda\geq1$ and $s\leq\tr Y/(1-e^{-\lambda})\leq s_1$ by
\eqref{eq:Y} and Theorem~\ref{thm:energy}.

(ii) By Proposition~\ref{prop:imported}(ii), Theorem~\ref{thm:energy} and
(i), for $t\geq t_1$,
\[
 \norm{\mathcal E_V}_1\leq2\sqrt{d_0qs/\lambda}+r\kappa_0s
 \leq\mathfrak c_1e^{-\frac{\alpha_1}4(t-t_1)},
 \qquad\norm J_1\leq rs_1e^{-\frac{\alpha_1}2(t-t_1)}.
\]
Let $v$ be a unit vector of $(\ker D)^\perp$ with $F^\dagger Fv=\lambda v$ and
$u=Fv/\sqrt\lambda$, so $Lu=\lambda u$. Apply
Proposition~\ref{prop:imported}(v) to $T=\Omega_V$ with
$T_\circ=\diag(z_0\id+C_0,\ z_0e^{-F^\dagger F})$, for which
$|T_\circ|v=|z_0|e^{-\lambda}v$ and $\Psi_\circ v=\alpha_0v$; and to
$T=-\Omega_U=z_0e^{-L}-J$ with $T_\circ=z_0e^{-L}$ and the vector $u$:
\[
 |\inner v{(\Phi_V-\alpha_0)v}|\leq\pi\sqrt{\frac{\norm{\mathcal E_V}_1e^\lambda}{|z_0|}},
 \qquad
 |\inner u{(\Phi_U+\pi-\alpha_0)u}|\leq\pi\sqrt{\frac{\norm J_1e^\lambda}{|z_0|}} .
\]
By Proposition~\ref{prop:imported}(i), the right derivative of $\lambda$
satisfies
\[
 \dot\lambda\geq\lambda\bigl[\pi-\mathfrak c_2e^{\lambda/2-\frac{\alpha_1}8(t-t_1)}\bigr]\qquad(t\geq t_1).
\]
Consider the line $t\mapsto\frac{\alpha_1}4(t-t_1)-\mathfrak c_3$. Where
$\lambda$ is at most its value we have
$\mathfrak c_2e^{\lambda/2-\frac{\alpha_1}8(t-t_1)}\leq\mathfrak c_2e^{-\mathfrak c_3/2}\leq\frac\pi2$,
hence $\dot\lambda\geq\frac\pi2\lambda\geq\frac\pi2>\frac{\alpha_1}4$, because
$\lambda\geq\Lambda_1\geq1$ and $\alpha_1<\pi$. Since
$\lambda(t_1)\geq1>-\mathfrak c_3$, the graph of $\lambda$ cannot cross below the line.

(iii) For $t\geq t_1$:
$\norm{\Omega_V-(z_0\id+C_0)\oplus0}_1\leq\norm{\mathcal E_V}_1+|z_0|re^{-\lambda}$;
$\norm{\Omega_U}_1\leq rs+|z_0|re^{-\lambda}$; and
$\norm{\Omega_W-(z_1\id-C_1)}_1=\tr Y\leq s$. By (ii),
$e^{-\lambda}\leq e^{\mathfrak c_3}e^{-\frac{\alpha_1}4(t-t_1)}$. Adding gives the
bound, and the statement on the masses follows from the triangle inequality
for the trace norm.

(iv) For $t\geq t_2$, (i) gives $s\leq\frac{\pi y_0}{8\alpha_1}$. Then
$\frac{qs}{\lambda y_0^2}\leq\frac{\bar N}{\Lambda_1}\frac s{y_0^2}
\leq\frac{\alpha_1s}{8\pi y_0}\leq\frac1{64}$, and
Proposition~\ref{prop:imported}(iv) gives
$\norm{\delta_0}_{\HS}^2\leq\frac{16\pi^2}{y_0^2}\frac{qs}\lambda
\leq\frac{2\pi\alpha_1}{y_0}s$.
\end{proof}

\begin{remark}\label{rem:rates}
The time $t_1$ is an a priori bound; it contains $e^{\Lambda_1}$ and
$\log^2(1+1/\lambda(0))$ and exceeds by many orders of magnitude the time at
which the decay is observed to begin in numerical integrations. For the
symmetric charges $\alpha_1=\frac\pi4$, so the theorem gives the rate
$\frac\pi8$ for $s$ and $\frac\pi{16}$ for the blocks of $\Omega$ on the
acyclic part. Numerical integration of \eqref{eq:mapflow}, which is not used
in any proof and is not reproduced here, suggests that
$\dot\lambda/\lambda$ approaches $\frac\pi2$, so that $\lambda$ would grow
exponentially.
\end{remark}

\section{Ordered and reversed charges: the limit on the cohomology}\label{sec:convergence}

By Theorem~\ref{thm:cancellation} the long-time behaviour is that of the
metrics induced on the cohomology representation, which move by its own
phases up to an exponentially small error. We treat ordered charges with
semistable and with polystable cohomology, and reversed charges.

For ordered charges, the \emph{Kempf--Ness function} of the cohomology
representation for the stability parameter $(\theta_0,\theta_1)$ of
\eqref{eq:theta} is the function of a pair $(h_0,h_1)$ of Hermitian metrics
on $\ker D$ and $W$
\begin{equation}\label{eq:KN}
 \cK_{\coh}(h_0,h_1)=q+\theta_0\log\det h_0-\theta_1\log\det h_1,\qquad
 q=\sum_j\tr\bigl(h_0^{-1}b_j^*h_1b_j\bigr),
\end{equation}
where $b_j$ is the matrix of $A_j|_{\ker D}$ in fixed reference bases,
$b_j^*$ its conjugate transpose, and $h_0,h_1$ stand for the Gram
matrices. Writing $h_i=g_i^*g_i$, one has
$\cK_{\coh}=\norm{g\cdot b}^2+\theta_0\log|\det g_0|^2-\theta_1\log|\det g_1|^2$
with $b=(b_j)$ and $(g\cdot b)_j=g_1b_jg_0^{-1}$; so $\cK_{\coh}$ is the
Kempf--Ness function of the point $b$ on the symmetric space
$\mathrm{GL}_{n_0}\times\mathrm{GL}_{n_1}/\mathrm U(n_0)\times\mathrm U(n_1)$
of pairs of metrics. It is invariant under a common rescaling of $h_0,h_1$ because
$n_0\theta_0=n_1\theta_1$, and its critical points are the zeros of the
moment map $\mu=(C_0-\theta_0\id,\,C_1-\theta_1\id)$
(Lemma~\ref{lem:knvariation}).

Put $\cK_*:=\inf\cK_{\coh}\in[-\infty,\infty)$. A zero of the moment map exists
if and only if the cohomology representation is polystable
\cite[Proposition~6.5]{kin94}, and $\cK_*>-\infty$ if and only if it is
semistable \cite[Theorem~5.1]{hkkp23}. King's results are stated for integral
characters. They apply here because $(\theta_0,\theta_1)$ is a positive
multiple $c\,(n_1,n_0)$ of an integral vector, and the Kempf--Ness function
for the stability parameter $c\,(n_1,n_0)$ at the arrows $b$ is $c$ times
the one for $(n_1,n_0)$ at the arrows $b/\sqrt c$, which have the same
subrepresentations.

For the second statement we recall the argument. If a subrepresentation
$(V',W')$ of dimension vector $(a_0,a_1)$ has $a_0\theta_0-a_1\theta_1>0$, let
$h_0=1\oplus e^\tau$ on $V'\oplus V'^\perp$ and $h_1=1\oplus e^\tau$ on
$W'\oplus W'^\perp$, for reference metrics; then $q$ stays bounded as
$\tau\to\infty$, and
$\cK_{\coh}=q-\tau(a_0\theta_0-a_1\theta_1)\to-\infty$, since
$n_0\theta_0=n_1\theta_1$. Conversely, for a semistable representation and
an integral character, the closure of the orbit of a lift of the
representation to the total space of the character does not meet the zero
section \cite[Proposition~3.1 and Lemma~2.2]{kin94}, the norm function on
the total space considered in \cite[\S6]{kin94} is proper on this closed set
\cite[Lemma~6.3]{kin94}, hence bounded below by a
positive constant there, and $\frac12\cK_{\coh}$ is its logarithm on the
orbit.

We write
$\norm\mu^2=\norm{\mu_0}_{\HS}^2+\norm{\mu_1}_{\HS}^2$ and
$\norm\delta^2=\norm{\delta_0}_{\HS}^2+\norm{\delta_1}_{\HS}^2$; for a pair of
matrices $\norm\cdot$ denotes this Hilbert--Schmidt norm, and for a single
matrix it remains the operator norm.

On a representation, regarded as a complex in degree zero, the normalized
metrics move with velocity $(G_0(C_0),-G_1(C_1))$ for increasing functions
$G_i$ with $G_i(\theta_i)=0$ (see \eqref{eq:velocities}), whereas the downward
gradient flow of the Kempf--Ness function has velocity
$(C_0-\theta_0,-(C_1-\theta_1))$. So the phase flow is not that gradient flow,
but it has the same stationary points and its velocity pairs negatively with
the gradient $(-\mu_0,\mu_1)$.

\subsection{The moment map tends to zero}
Let the charges be ordered. Besides \eqref{eq:constants} put
\[
 \Gamma=\max\Bigl\{\frac{(|z_0|+\bar N)^2}{y_0},\frac{(|z_1|+\bar N)^2}{y_1}\Bigr\},\qquad
 c_\delta^2=\frac{2\pi\alpha_1s_1}{y_0}+\frac{s_1^2}{y_1^2}
\]
(the formula for $c_\delta$ makes sense for all charges and is used again in
Theorem~\ref{thm:reversed}). We note that
\begin{equation}\label{eq:thetabound}
 \theta_0,\ \theta_1\ \leq\ \vartheta\ \leq\ \bar N .
\end{equation}
Indeed $\theta_1=x_1-y_1\cot\varphi$. If $\cot\varphi\geq0$, then
$0<\theta_1\leq x_1\leq M_1/\pi$. If $\cot\varphi<0$, then
$\frac\pi2<\varphi<\alpha_0$, so
$|\cot\varphi|=\cot(\pi-\varphi)\leq(\pi-\varphi)^{-1}<(\pi-\alpha_0)^{-1}\leq\rho^{-1}$
and $\theta_1\leq\max(x_1,0)+y_1/\rho\leq M_1/\pi+M_1/(2\rho)$. In both cases
$\theta_1\leq2M_1/\rho$, that is, $\vartheta=n_1\theta_1\leq\bar N$. Let
\[
 \mathfrak h(t)=\bigl(e^{\varphi t}h_V(t)|_{\ker D},\ e^{\varphi t}h_W(t)\bigr)
\]
be the pair of induced metrics on $\ker D$ and $W$, multiplied by
$e^{\varphi t}$. The function $\cK_{\coh}$ is invariant under common rescaling, so
$\mathcal E(t):=\cK_{\coh}(\mathfrak h(t))-\cK_*$ is the excess of the induced pair.

\begin{theorem}[Decay of the moment map]\label{thm:relaxation}
Let the charges be ordered and the cohomology representation semistable.
Then along every trajectory, for $t\geq t_2$,
\begin{equation}\label{eq:KNderivative}
 \frac{d}{dt}\cK_{\coh}(\mathfrak h(t))
 \leq-\frac1\Gamma\norm\mu^2+\norm\mu\norm\delta
 \leq-\frac1{2\Gamma}\norm\mu^2+\frac\Gamma2\norm\delta^2,
 \qquad\norm{\delta(t)}\leq c_\delta\,e^{-\frac{\alpha_1}4(t-t_1)},
\end{equation}
and
\[
 \mathcal E(t)+\frac1{2\Gamma}\int_{t_2}^t\norm\mu^2
 \leq\mathcal E^\sharp:=\mathcal E(0)+2\pi(N_0+\vartheta)\,t_2+\frac\Gamma{\alpha_1}c_\delta^2 ,
\]
with $N_0=\max\{N(0),n_1M_1/\rho\}$ as in Theorem~\ref{thm:energy} and
$\delta=(\delta_0,\delta_1)$, $\delta_0=\Phi_{HH}-\phi_0$,
$\delta_1=\Phi_W-\phi_1$, as in Section~\ref{sec:setting}. Moreover
$\mu(t)\to0$, $\cK_{\coh}(\mathfrak h(t))\to\cK_*$ and $m(h(t))\to|Z|$ as $t\to\infty$.
\end{theorem}

\begin{proof}
By \eqref{eq:flow} and \eqref{eq:adapted} the pair $\mathfrak h(t)$ has velocities
$\xi_0=\varphi-\Phi_{HH}$ and $\xi_1=\varphi-\Phi_W$ in the frames induced by
the adapted frames; that is,
\begin{equation}\label{eq:velocities}
 \begin{gathered}
 \xi_0=G_0(C_0)-\delta_0,\qquad\xi_1=-G_1(C_1)-\delta_1,\\
 G_0(x)=\varphi-\arg(z_0+x),\qquad G_1(x)=\arg(z_1-x)-\varphi .
 \end{gathered}
\end{equation}
The function $G_0$ is increasing with $G_0(\theta_0)=0$, because
$\Ima(e^{-\ii\varphi}(z_0+x))=c_0-\sigma x$, and
$G_0'(x)=y_0/|z_0+x|^2\in[y_0/(|z_0|+\bar N)^2,\ 1/y_0]$ on $[0,\bar N]$;
likewise $G_1$ is increasing with $G_1(\theta_1)=0$ and
$G_1'(x)=y_1/|z_1-x|^2\in[y_1/(|z_1|+\bar N)^2,\ 1/y_1]$. By
Theorem~\ref{thm:energy} the spectra of $C_0,C_1$ lie in $[0,\bar N]$ for
$t\geq t_0$, and $\theta_0,\theta_1\leq\bar N$ by \eqref{eq:thetabound}. The first variation of \eqref{eq:KN} is
$d\cK_{\coh}(\xi)=-\tr(\xi_0\mu_0)+\tr(\xi_1\mu_1)$
(Lemma~\ref{lem:knvariation}), so
\[
 \frac d{dt}\cK_{\coh}(\mathfrak h(t))
 =-\tr\bigl(\mu_0G_0(C_0)\bigr)-\tr\bigl(\mu_1G_1(C_1)\bigr)
 +\tr(\mu_0\delta_0)-\tr(\mu_1\delta_1),
\]
and $\tr(\mu_0G_0(C_0))=\sum(c-\theta_0)G_0(c)\geq\Gamma^{-1}\norm{\mu_0}_{\HS}^2$
over the eigenvalues $c$ of $C_0$, likewise for $\mu_1$. This is the first
inequality of \eqref{eq:KNderivative}, for all $t\geq t_0$; the bound on
$\norm\delta$ for $t\geq t_2$ is Theorem~\ref{thm:cancellation}(i),(iv).

Since $\norm{\xi_0},\norm{\xi_1}\leq\pi$ and
$\norm{\mu_0}_1\leq q+n_0\theta_0$, $\norm{\mu_1}_1\leq q+n_1\theta_1$, we have
$|\frac d{dt}\cK_{\coh}(\mathfrak h(t))|\leq2\pi(N_0+\vartheta)$ for all $t$, hence
$\mathcal E(t_2)\leq\mathcal E(0)+2\pi(N_0+\vartheta)t_2$. Integrating
\eqref{eq:KNderivative} from $t_2$ and using
$\frac\Gamma2\int_{t_2}^\infty\norm\delta^2\leq\frac\Gamma{\alpha_1}c_\delta^2$
and $\mathcal E\geq0$ gives the displayed bound. By \eqref{eq:adapted} and
$\norm{\Phi_\bullet}\leq\pi$,
$\norm{\dot\mu_0}_{\HS},\norm{\dot\mu_1}_{\HS}\leq2\pi q\leq2\pi N_0$, so
$\norm\mu^2$ is uniformly continuous and integrable and therefore tends to
zero. Then $C_0\to\theta_0\id$ and $C_1\to\theta_1\id$, so
$m_H\to n_0|z_0+\theta_0|+n_1|z_1-\theta_1|$. Both $z_0+\theta_0$ and
$z_1-\theta_1$ have argument $\varphi$, so this limit is
$\Rea\bigl(e^{-\ii\varphi}(Z+n_0\theta_0-n_1\theta_1)\bigr)=|Z|$, and
$m\to|Z|$ by Theorem~\ref{thm:cancellation}(iii).

It remains to show $\mathcal E(t)\to0$. By \eqref{eq:KNderivative} the
function $\mathcal E(t)+\frac\Gamma2\int_t^\infty\norm\delta^2$ is
nonincreasing on $[t_2,\infty)$, so $\mathcal E(t)$ has a limit
$\mathcal E_\infty\geq0$. Suppose $\mathcal E_\infty>0$ and choose a pair
$\mathfrak h'$ with $\cK_{\coh}(\mathfrak h')-\cK_*<\mathcal E_\infty/2$. By convexity of
$\cK_{\coh}$ along the geodesic from $\mathfrak h(t)$ to $\mathfrak h'$ in the
space of pairs of metrics, with its invariant distance $d$ (these are
recalled in the next subsection; Lemma~\ref{lem:knvariation}),
$\cK_{\coh}(\mathfrak h(t))-\cK_{\coh}(\mathfrak h')\leq\norm{\mu(t)}\,d(\mathfrak h(t),\mathfrak h')$, and the left
side is at least $\mathcal E_\infty/4$ for large $t$. By
\eqref{eq:velocities} the speed of $\mathfrak h(t)$ is at most
$\norm\mu/\min(y_0,y_1)+\norm\delta$, which is square integrable, so
$d(\mathfrak h(t),\mathfrak h')\leq c+c'\sqrt t$ by the Cauchy--Schwarz inequality. Then
$\norm{\mu(t)}^2\geq\mathcal E_\infty^2/(16(c+c'\sqrt t)^2)$ is not
integrable, a contradiction.
\end{proof}

Theorem~\ref{thm:relaxation} does not require a zero of the moment map to
exist: for semistable cohomology that is not polystable the infimum $\cK_*$
is not attained, and $\cK_{\coh}(\mathfrak h(t))\to\cK_*$ is the analogue for the phase
flow of the corresponding fact for gradient flow lines of a Kempf--Ness
function \cite[Theorem~4.3]{grs21}.

\subsection{The Kempf--Ness function on the space of metric pairs}
The convexity statements of this subsection are classical in the theory of
Kempf--Ness functions; see \cite[Theorem~4.3]{grs21}, where they are proved
for actions on closed K\"ahler manifolds. We include the short proofs for the
linear action at hand. Let $\mathcal X$ be the space of pairs $\mathfrak h=(\mathfrak h_0,\mathfrak h_1)$
of Hermitian metrics on $\ker D$ and $W$. A tangent vector at $\mathfrak h$ is a pair
$\xi=(\xi_0,\xi_1)$ of $\mathfrak h$-Hermitian endomorphisms, $\delta\mathfrak h=\mathfrak h\xi$, and
$\norm\xi^2=\tr\xi_0^2+\tr\xi_1^2$ defines an invariant Riemannian metric, for
which $\mathcal X$ is complete, simply connected and nonpositively curved.
Its geodesics are $u\mapsto \mathfrak he^{u\xi}$, the exponential map at any point is a
diffeomorphism, and closed balls are compact. Let $d$ be the distance. In a
$\mathfrak h$-orthonormal frame let $H_j$ be the matrix of $A_j|_{\ker D}$, and $C_0$,
$C_1$, $\mu_0$, $\mu_1$, $q$ as before.

\begin{lemma}\label{lem:knvariation}
$d\cK_{\coh}(\xi)=-\tr(\xi_0\mu_0)+\tr(\xi_1\mu_1)$, so the gradient of
$\cK_{\coh}$ is $(-\mu_0,\mu_1)$, of norm $\norm\mu$. Along a geodesic,
\[
 \frac{d^2}{du^2}\cK_{\coh}(\mathfrak he^{u\xi})
 =\sum_j\norm{\xi_1H_j(u)-H_j(u)\xi_0}_{\HS}^2\geq0,\qquad
 H_j(u)=e^{u\xi_1/2}H_je^{-u\xi_0/2}.
\]
\end{lemma}

\begin{proof}
In the $\mathfrak h$-orthonormal frame,
$q(\mathfrak he^{u\xi})=\sum_j\tr(e^{-u\xi_0}H_j^\dagger e^{u\xi_1}H_j)$, and
$\log\det$ is affine along geodesics with derivative $\tr\xi$. Differentiate
once and twice.
\end{proof}

So $\cK_{\coh}$ is geodesically convex and its critical points, the zeros of
the moment map, are its minimizers. Let $\mathcal Z\subset\mathcal X$ be the
set of zeros of the moment map. For $\mathfrak h\in\mathcal X$ let $\mathfrak a(\mathfrak h)$ be
the space of $\mathfrak h$-Hermitian pairs $\xi$ with $\xi_1A_j=A_j\xi_0$ on $\ker D$
for all $j$, the Hermitian infinitesimal automorphisms of the
representation.

\begin{lemma}\label{lem:minimumset}
For $p\in\mathcal Z$, $\mathcal Z=\{pe^\xi:\xi\in\mathfrak a(p)\}$. It is a
closed totally geodesic submanifold with $T_p\mathcal Z=\mathfrak a(p)$.
\end{lemma}

\begin{proof}
If $\xi\in\mathfrak a(p)$, the second derivative of $\cK_{\coh}$ along
$u\mapsto pe^{u\xi}$ vanishes identically by Lemma~\ref{lem:knvariation}, and
the first derivative vanishes at $u=0$; so $pe^\xi\in\mathcal Z$. Conversely,
if $pe^\xi\in\mathcal Z$, the convex function
$u\mapsto\cK_{\coh}(pe^{u\xi})$ takes its minimum at $u=0$ and $u=1$, hence is
constant on $[0,1]$, and its second derivative at $0$ gives
$\xi\in\mathfrak a(p)$. $\mathcal Z=\{\mu=0\}$ is closed, and it is the image
of a linear subspace under the exponential map at each of its points.
\end{proof}

\begin{lemma}\label{lem:gap}
For $p\in\mathcal Z$ put
\[
 \nu_*:=\min\Bigl\{\sum_j\norm{\xi_1H_j-H_j\xi_0}_{\HS}^2:\
 \xi\perp\mathfrak a(p),\ \norm\xi=1\Bigr\}>0,\qquad\varrho_*:=\frac{\log2}{\sqrt2}.
\]
Then $\nu_*$ does not depend on $p\in\mathcal Z$, and for every unit
$\xi\perp\mathfrak a(p)$ and $0\leq u\leq\varrho_*$,
$\frac{d^2}{du^2}\cK_{\coh}(pe^{u\xi})\geq\nu_*/2$.
\end{lemma}

\begin{proof}
The quadratic form is positive on the orthogonal complement of its kernel
$\mathfrak a(p)$. By Lemma~\ref{lem:knvariation},
$\xi_1H_j(u)-H_j(u)\xi_0=e^{u\xi_1/2}(\xi_1H_j-H_j\xi_0)e^{-u\xi_0/2}$ has
Hilbert--Schmidt norm at least
$e^{-u(\norm{\xi_1}+\norm{\xi_0})/2}\norm{\xi_1H_j-H_j\xi_0}_{\HS}$, and
$\norm{\xi_1}+\norm{\xi_0}\leq\sqrt2$ for a unit $\xi$. So the second
derivative is at least $e^{-\sqrt2u}\nu_*$. If $p'=pe^\zeta$ with
$\zeta\in\mathfrak a(p)$, then $e^{-\zeta/2}$ is a pair of automorphisms of
the representation that maps $p$-orthonormal frames to $p'$-orthonormal
frames and preserves the matrices $H_j$; conjugation by it carries
$\mathfrak a(p)$ and its orthogonal complement isometrically to those of
$p'$.
\end{proof}

\begin{lemma}\label{lem:kngrowth}
Let the representation be polystable, so that $\mathcal Z\neq\emptyset$.
For $\mathfrak h\in\mathcal X$ put $\mathsf d=d(\mathfrak h,\mathcal Z)$ and
$\mathcal E=\cK_{\coh}(\mathfrak h)-\cK_*$. Then
\begin{enumerate}
\item[(i)] $\mathcal E\leq\norm{\mu(\mathfrak h)}\,\mathsf d$;
\item[(ii)] $\mathcal E\geq\frac{\nu_*}4\min(\mathsf d^2,\varrho_*\mathsf d)$;
\item[(iii)] if $\mathsf d\leq\varrho_*$, then $\norm{\mu(\mathfrak h)}\geq\frac{\nu_*}2\mathsf d$
 and $\mathcal E\leq\frac2{\nu_*}\norm{\mu(\mathfrak h)}^2$.
\end{enumerate}
\end{lemma}

\begin{proof}
Let $p\in\mathcal Z$ be a nearest point and $\mathfrak h=pe^{\mathsf d\xi}$ with
$\norm\xi=1$; then $\xi\perp\mathfrak a(p)$. The function
$f(u)=\cK_{\coh}(pe^{u\xi})-\cK_*$ is convex with $f(0)=f'(0)=0$.
(i) $0=f(0)\geq f(\mathsf d)-\mathsf d f'(\mathsf d)$ and $f'(\mathsf d)\leq\norm{\mu(\mathfrak h)}$.
(ii) $f''\geq\nu_*/2$ on $[0,\varrho_*]$ by Lemma~\ref{lem:gap}, so
$f(u)\geq\frac{\nu_*}4u^2$ there, and $f(u)\geq f(\varrho_*)u/\varrho_*$ for
$u\geq\varrho_*$ by convexity. (iii) $\norm{\mu(\mathfrak h)}\geq f'(\mathsf d)
=\int_0^\mathsf d f''\geq\frac{\nu_*}2\mathsf d$; combine with (i).
\end{proof}

Part (iii) is a \L ojasiewicz gradient inequality with exponent $\frac12$
for $\cK_{\coh}$ near $\mathcal Z$.


\subsection{Convergence for polystable cohomology}

\begin{theorem}[Convergence to a harmonic metric]\label{thm:convergence}
Let the charges be ordered and the cohomology representation polystable.
Then along every trajectory there is a zero $\mathfrak h_\infty$ of the moment map,
that is, a harmonic metric of the cohomology representation in the sense of
\cite{hkkp26}, with
\[
 d\bigl(\mathfrak h(t),\mathfrak h_\infty\bigr)\leq\mathfrak C\,e^{-\beta t}\qquad(t\geq0)
\]
for every rate $\beta<\frac12\min\{\nu_*/(4\Gamma),\ \alpha_1/2\}$, where the
constant $\mathfrak C$ depends only on $\beta$, the dimensions, the charges, $N(0)$, $t_2$,
$\mathcal E(0)$ and $\nu_*$. Moreover, for $t\geq t_1$,
\[
 \norm{\mu(\mathfrak h(t))}\leq2\bar N\,d(\mathfrak h(t),\mathfrak h_\infty),\qquad
 0\leq m_H(t)-|Z|\leq\frac{\norm{\mu(\mathfrak h(t))}^2}{2\min(y_0,y_1)},
\]
so that the moment map of the cohomology representation and $m(h(t))-|Z|$
tend to zero exponentially.
\end{theorem}

The convergence is as Hermitian forms on the fixed vector spaces $\ker D$ and
$W$: no time-dependent automorphism of the representation is applied. The
limit depends on the initial metric through its position on $\mathcal Z$.

\begin{proof}
By Theorem~\ref{thm:relaxation}, \eqref{eq:KNderivative} holds for
$t\geq t_2$ and $\mathcal E(t)\leq\mathcal E^\sharp$ there.

\emph{Step 1.} By Lemma~\ref{lem:kngrowth}(ii),
$d(\mathfrak h(t),\mathcal Z)\leq\mathsf d_{\max}:=\max\{\varrho_*,4\mathcal E^\sharp/(\nu_*\varrho_*)\}$
for $t\geq t_2$, and by (i), $\mathcal E(t)\leq\norm{\mu(t)}\mathsf d_{\max}$.

\emph{Step 2.} By Theorem~\ref{thm:relaxation},
$\int_{t_2}^\infty\norm\mu^2\leq2\Gamma\mathcal E^\sharp$ and
$\mathcal E(t)\leq\mathcal E(t')+\frac\Gamma2\int_{t'}^\infty\norm\delta^2$ for
$t\geq t'\geq t_2$. For $T>0$ there is $t'\in[t_2+T/2,t_2+T]$ with
$\norm{\mu(t')}^2\leq4\Gamma\mathcal E^\sharp/T$, hence for $t\geq t_2+T$
\[
 \mathcal E(t)\leq\mathsf d_{\max}\sqrt{4\Gamma\mathcal E^\sharp/T}
 +\frac\Gamma{\alpha_1}c_\delta^2e^{-\frac{\alpha_1}4T}.
\]
Choose $T$, depending only on the dimensions, the charges,
$\mathcal E^\sharp$ and $\nu_*$, so that the right side is at most
$\frac{\nu_*}4\varrho_*^2$; then $d(\mathfrak h(t),\mathcal Z)\leq\varrho_*$ for
$t\geq t_3:=t_2+T$ by Lemma~\ref{lem:kngrowth}(ii). For such $t$,
Lemma~\ref{lem:kngrowth}(iii) and \eqref{eq:KNderivative} give
\[
 \dot{\mathcal E}\leq-\frac{\nu_*}{4\Gamma}\mathcal E
 +\frac\Gamma2c_\delta^2e^{-\frac{\alpha_1}2(t-t_1)},
\]
so $\mathcal E(t)\leq\mathfrak C_2e^{-2\beta t}$ for every
$2\beta<\min\{\nu_*/(4\Gamma),\alpha_1/2\}$.

\emph{Step 3.} Integrating \eqref{eq:KNderivative} over $[t,t+1]$,
$\int_t^{t+1}\norm\mu^2\leq2\Gamma\mathcal E(t)+\Gamma^2\int_t^{t+1}\norm\delta^2
\leq\mathfrak C_3e^{-2\beta t}$. By \eqref{eq:velocities} and $|G_0'|\leq1/y_0$,
$|G_1'|\leq1/y_1$, the speed of $\mathfrak h(t)$ is
$\norm\xi\leq\norm\mu/\min(y_0,y_1)+\norm\delta$. Hence
$\int_t^\infty\norm\xi\leq\mathfrak C_4e^{-\beta t}$ for $t\geq t_3$, and by
completeness $\mathfrak h(t)$ converges to some $\mathfrak h_\infty$ with
$d(\mathfrak h(t),\mathfrak h_\infty)\leq\mathfrak C_4e^{-\beta t}$; $\mu(\mathfrak h_\infty)=0$ by continuity. On
$[0,t_3]$ the speed is at most $\pi\sqrt{n_0+n_1}$, because the spectra of
$\Phi_{HH}$ and $\Phi_W$ lie in $(0,\pi)$; this gives the dependence of $\mathfrak C$.

\emph{The moment map and the mass.} Let $t\geq t_1$ and let
$u\mapsto \mathfrak h_\infty e^{u\zeta}$, $0\leq u\leq1$, be the geodesic from
$\mathfrak h_\infty$ to $\mathfrak h(t)$. The function $q$ is convex along geodesics, so
$q\leq\max(\vartheta,q(\mathfrak h(t)))\leq\bar N$ on it, by
\eqref{eq:thetabound}; here $q(\mathfrak h_\infty)=\tr C_0=\vartheta$. In the frames of
Lemma~\ref{lem:knvariation},
$\frac d{du}C_0=-\frac12(\zeta_0C_0+C_0\zeta_0)+\sum_jH_j^\dagger\zeta_1H_j$
and similarly for $C_1$, so $\norm{\frac d{du}\mu}\leq2q\norm\zeta$, and
$\mu(\mathfrak h_\infty)=0$. For the mass, the functions $x\mapsto|z_0+x|$ and
$x\mapsto|z_1-x|$ are convex with second derivatives at most $1/y_0$ and
$1/y_1$, and their first-order terms at $\theta_0,\theta_1$ cancel in the sum
over the eigenvalues of $C_0$ and $C_1$ because $\tr\mu_0=\tr\mu_1$. Combine
with Theorem~\ref{thm:cancellation}(iii).
\end{proof}

\begin{remark}\label{rem:rate}
The rate is limited by the gap $\nu_*$ of the Hessian of $\cK_{\coh}$
transversally to $\mathcal Z$ and by the decay of $A_jD^+$. When the
cohomology representation is stable, $\mathfrak a$ consists of the scalar
pairs $(c,c)$.
\end{remark}

\subsection{Semistable cohomology that is not polystable}
In this case Theorem~\ref{thm:relaxation} gives $\mu\to0$, and there is no
limit metric.

\begin{proposition}\label{prop:divergence}
Let the charges be ordered and let the cohomology representation be
semistable and not polystable. Then along every trajectory
$d(\mathfrak h(t),\mathfrak h(0))\to\infty$, and for $t\geq t_2$
\[
 d\bigl(\mathfrak h(t),\mathfrak h(t_2)\bigr)
 \leq\frac{\sqrt{2\Gamma\mathcal E^\sharp}}{\min(y_0,y_1)}\,\sqrt{t-t_2}
 +\frac{4c_\delta}{\alpha_1}.
\]
\end{proposition}

\begin{proof}
Suppose $d(\mathfrak h(t_k),\mathfrak h(0))$ is bounded along a sequence
$t_k\to\infty$. Closed balls of $\mathcal X$ are compact, so a subsequence
converges to a pair $\mathfrak h_\infty$. The moment map is continuous on
$\mathcal X$ and $\mu(\mathfrak h(t))\to0$ by
Theorem~\ref{thm:relaxation}; so $\mu(\mathfrak h_\infty)=0$ and the
representation is polystable, a contradiction. By \eqref{eq:velocities} and
$0<G_i'\leq1/y_i$ on $[0,\infty)$, the speed of $\mathfrak h(t)$ is at most
$\norm\mu/\min(y_0,y_1)+\norm\delta$. By Theorem~\ref{thm:relaxation} and the
Cauchy--Schwarz inequality
$\int_{t_2}^t\norm\mu\leq(2\Gamma\mathcal E^\sharp)^{1/2}\sqrt{t-t_2}$, and
$\int_{t_2}^\infty\norm\delta\leq4c_\delta/\alpha_1$ because $t_2\geq t_1$.
\end{proof}

The bound $\sqrt t$ is not attained in the following example, the smallest
one: there the distance grows like $\log t$, and the asymptotics can be
determined for every initial metric. Let $\mathsf R$ be the representation of
the Kronecker quiver with two arrows on $\CC^2\to\CC^2$ with the arrows
\[
 b_1=\id,\qquad b_2=E_{12},\qquad E_{12}e_2=e_1,\quad E_{12}e_1=0 .
\]
Here $n_0=n_1=2$, so $\theta_0=\theta_1=:\theta$. A subrepresentation
$(V',W')$ has $V'=b_1V'\subset W'$, so $\mathsf R$ is semistable; the
lines $\CC e_1\to\CC e_1$ form a subrepresentation of dimension vector
$(1,1)$, so it is not stable; and it is not polystable, since a direct sum
of two representations of dimension vector $(1,1)$ has arrows that are
diagonal in suitable bases of the two spaces, and then $b_1^{-1}b_2=E_{12}$
would be diagonalizable. Both $\CC e_1\to\CC e_1$ and the quotient are
isomorphic to the stable representation $\CC\to\CC$ with arrows $(1,0)$. Put
\[
 \mathsf g_0=G_0'(\theta)=\frac{y_0}{|z_0+\theta|^2},\qquad
 \mathsf g_1=G_1'(\theta)=\frac{y_1}{|z_1-\theta|^2},\qquad
 \mathsf k=\frac{2\mathsf g_0\mathsf g_1}{\mathsf g_0+\mathsf g_1},\qquad
 \Lambda_\infty=(\mathsf g_0+\mathsf g_1)\,\theta,
\]
with $G_0,G_1$ as in \eqref{eq:velocities}. For a pair of metrics
$(h_0,h_1)$ on $\mathsf R$, with Hermitian forms conjugate-linear in the
first argument, let
\[
 p_1=h_0(e_1,e_1),\qquad\chi_0=\frac{h_0(e_1,e_2)}{p_1},\qquad
 p_2=h_0(e_2-\chi_0e_1,\,e_2-\chi_0e_1),
\]
and define $p_1',\chi_1,p_2'$ in the same way from $h_1$; these are the
Gram--Schmidt data of the two metrics for the flag $\CC e_1\subset\CC^2$.
The metrics are diagonal in the basis $e_1,e_2$ exactly when
$\chi_0=\chi_1=0$, and the automorphisms $\id+xE_{12}$, $x\in\CC$, of
$\mathsf R$ shift $\chi_0$ and $\chi_1$ by $x$. Put
\[
 a=\frac{p_1'}{p_1},\qquad b=\frac{p_2'}{p_2},\qquad c=\frac{p_1'}{p_2},\qquad
 \chi=\chi_1-\chi_0 .
\]

\begin{proposition}[A semistable example]\label{prop:semistable}
Let the charges be ordered. For every initial pair of metrics on
$\mathsf R$, regarded as a complex in degree zero, along the phase flow:
\begin{enumerate}
\item[(i)] $\chi$ vanishes identically or has constant argument, it tends to
 zero with $|\chi(t)|=O(e^{-\Lambda't})$ for every $\Lambda'<\Lambda_\infty$,
 and $\chi_0$, $\chi_1$ converge to a common limit $\chi_\infty$;
\item[(ii)] $c(t)=\dfrac1{\mathsf kt}+O\Bigl(\dfrac{\log t}{t^2}\Bigr)$,
 $a-\theta=-\dfrac{\mathsf g_1}{\mathsf g_0+\mathsf g_1}c+O(c^2)$ and
 $b-\theta=-\dfrac{\mathsf g_0}{\mathsf g_0+\mathsf g_1}c+O(c^2)$;
\item[(iii)] $\displaystyle\lim_{t\to\infty}t\,\norm{\mu(t)}
 =\frac{\sqrt{2(\mathsf g_0^2+\mathsf g_1^2)}}{2\mathsf g_0\mathsf g_1}$ and
 $\cK_{\coh}(\mathfrak h(t))-\cK_*=\dfrac1{\mathsf kt}+O\Bigl(\dfrac{\log t}{t^2}\Bigr)$;
\item[(iv)] there is a constant $A>0$ with
 \[
 e^{\varphi t}p_1\sim A\,t^{-1/2},\quad e^{\varphi t}p_1'\sim\theta A\,t^{-1/2},\quad
 e^{\varphi t}p_2\sim\theta\mathsf kA\,t^{1/2},\quad
 e^{\varphi t}p_2'\sim\theta^2\mathsf kA\,t^{1/2},
 \]
 and $d(\mathfrak h(t),\mathfrak h(0))=\log t+O(1)$.
\end{enumerate}
\end{proposition}

So for every initial metric the normalized metrics decay like $t^{-1/2}$ on
the subrepresentation $\CC e_1\to\CC e_1$ and grow like $t^{1/2}$ on the
quotient. These rates hold in every basis $e_1,\ e_2-xe_1$; the value
$x=\chi_\infty$ is the one for which the metrics become diagonal, in the
sense that $h_i(e_1,e_2-\chi_\infty e_1)/h_i(e_1,e_1)\to0$ for $i=0,1$. The
exponents do not depend on the charges; the constants do, through $\theta$,
$\mathsf g_0$ and $\mathsf g_1$.

\begin{proof}
\emph{Frames.} The vectors $e_1/\sqrt{p_1}$, $(e_2-\chi_0e_1)/\sqrt{p_2}$ form
an $h_0$-orthonormal frame, and $e_1/\sqrt{p_1'}$,
$(e_2-\chi_1e_1)/\sqrt{p_2'}$ an $h_1$-orthonormal frame. Since
$e_2-\chi_0e_1=(e_2-\chi_1e_1)+\chi e_1$ and $E_{12}(e_2-\chi_0e_1)=e_1$, the
arrows have the matrices
$H_1=\bigl(\begin{smallmatrix}\sqrt a&\chi\sqrt c\\0&\sqrt b\end{smallmatrix}\bigr)$
and $H_2=\bigl(\begin{smallmatrix}0&\sqrt c\\0&0\end{smallmatrix}\bigr)$, so
\[
 C_0=\begin{pmatrix}a&\chi\sqrt{ac}\\\bar\chi\sqrt{ac}&b+c+|\chi|^2c\end{pmatrix},\qquad
 C_1=\begin{pmatrix}a+c+|\chi|^2c&\chi\sqrt{bc}\\\bar\chi\sqrt{bc}&b\end{pmatrix}.
\]
Let $F_0(x)=\arg(z_0+x)=\varphi-G_0(x)$ and
$F_1(x)=\arg(z_1-x)=\varphi+G_1(x)$; the matrices of the phases in these
frames are $F_0(C_0)$ and $F_1(C_1)$.

\emph{Evolution of the data.} For a metric $h$ with phase $\Phi$,
$\frac d{dt}h(u,w)=-h(u,\Phi w)$. Writing $e_1$, $e_2$ in the frame, one finds
for the data $(p_1,p_2,\chi_0)$ of $h_0$ and the matrix $\phi=F_0(C_0)$
\[
 \tfrac d{dt}\log p_1=-\phi_{11},\qquad\tfrac d{dt}\log p_2=-\phi_{22},\qquad
 \dot\chi_0=-\sqrt{p_2/p_1}\;\phi_{12},
\]
the second because $\det h_0=p_1p_2$ has logarithmic derivative $-\tr\phi$;
and likewise for $h_1$ with $\phi=F_1(C_1)$. For a $2\times2$ Hermitian matrix
$C$ with eigenvalues $c^\pm$ and a function $f$, $f(C)-f[c^+,c^-]\,C$ is a
scalar matrix, where $f[\,\cdot\,,\cdot\,]$ is the divided difference; so
$f(C)_{12}=f[c^+,c^-]\,C_{12}$. Let $\varrho_0=-F_0[c_0^+,c_0^-]>0$ and
$\varrho_1=F_1[c_1^+,c_1^-]>0$ for the eigenvalues of $C_0$, $C_1$. Since
$\sqrt{p_2/p_1}\sqrt{ac}=a$ and $\sqrt{p_2'/p_1'}\sqrt{bc}=b$,
\begin{equation}\label{eq:chi}
 \dot\chi_0=\varrho_0\,a\,\chi,\qquad\dot\chi_1=-\varrho_1\,b\,\chi,\qquad
 \dot\chi=-(a\varrho_0+b\varrho_1)\,\chi .
\end{equation}

\emph{Limits.} Theorem~\ref{thm:relaxation} applies to the flow of a
representation. Indeed, let $\mathsf E$ be the orthogonal direct sum of
$\mathsf R$ and a complex $(V_1\to0)\to(U\to0)$ with invertible differential.
It is of the form \eqref{eq:complex}, its endomorphism \eqref{eq:coeff} is
the direct sum of those of the two summands, so the flow preserves the
orthogonal splitting and restricts to the flow of each summand; the
cohomology of $\mathsf E$ is $\mathsf R$ with its own metrics, and
$\delta=0$. Hence $C_0\to\theta$ and $C_1\to\theta$, that is, $a\to\theta$,
$b\to\theta$, $c(1+|\chi|^2)\to0$, and $\varrho_i\to\mathsf g_i$. By
\eqref{eq:chi}, $\chi$ is a real positive multiple of $\chi(0)$, $|\chi|$ is
nonincreasing, and $a\varrho_0+b\varrho_1\to\Lambda_\infty$; this gives the
decay of $\chi$. Then
$\dot\chi_0$ and $\dot\chi_1$ are integrable, so $\chi_0$, $\chi_1$ converge,
to the same limit. This is (i).

\emph{The reduced system.} Put
\[
 w_1=-G_0(a)-G_1(a+c),\qquad w_2=-G_0(b+c)-G_1(b),\qquad
 w_3=-G_0(b+c)-G_1(a+c),
\]
functions of $(a,b,c)$. If $\chi=0$, the matrices $C_0$, $C_1$ are diagonal
and $\frac d{dt}\log a=w_1$, $\frac d{dt}\log b=w_2$,
$\frac d{dt}\log c=w_3$. In general, for a $2\times2$ Hermitian matrix $C$
with eigenvalues $c^\pm$ and a smooth function $f$, Newton interpolation of
$f$ at the nodes $C_{11},c^+,c^-$ gives
\[
 f(C)=f(C_{11})+f[C_{11},c^+]\,(C-C_{11})+f[C_{11},c^+,c^-]\,(C-C_{11})(C-c^+),
\]
and the upper left entries of $C-C_{11}$ and $(C-C_{11})(C-c^+)$ are $0$ and
$|C_{12}|^2$. So $|f(C)_{11}-f(C_{11})|\leq\frac12\sup|f''|\,|C_{12}|^2$, and
likewise for the other diagonal entry. Here $|C_{12}|^2\leq|\chi|^2c\,(a+b)$,
and the diagonal entries of $C_0$, $C_1$ differ from those for $\chi=0$ by
$|\chi|^2c$. Hence, with
$\varepsilon(t):=|\chi(t)|^2$,
\begin{equation}\label{eq:reduced}
 \tfrac d{dt}\log a=w_1+O(\varepsilon c),\qquad
 \tfrac d{dt}\log b=w_2+O(\varepsilon c),\qquad
 \tfrac d{dt}\log c=w_3+O(\varepsilon c),
\end{equation}
and $\frac d{dt}\log(e^{\varphi t}p_1)=G_0(a)+O(\varepsilon c)$, and similarly
for $p_2$, $p_1'$, $p_2'$ with $G_0(b+c)$, $-G_1(a+c)$, $-G_1(b)$.

\emph{The fast variables.} Differentiating along \eqref{eq:reduced},
\[
 \dot w_1=-\Lambda_aw_1-G_1'(a+c)\,c\,w_3+O(\varepsilon c),\qquad
 \dot w_2=-\Lambda_bw_2-G_0'(b+c)\,c\,w_3+O(\varepsilon c),
\]
with $\Lambda_a=a(G_0'(a)+G_1'(a+c))$ and $\Lambda_b=b(G_0'(b+c)+G_1'(b))$,
both tending to $\Lambda_\infty$. Near $(a,b,c)=(\theta,\theta,0)$ the
functions $(w_1,w_2,c)$ are real-analytic coordinates, with
$w_1=-(\mathsf g_0+\mathsf g_1)(a-\theta)-\mathsf g_1c$ and
$w_2=-(\mathsf g_0+\mathsf g_1)(b-\theta)-\mathsf g_0c$ to first order.
Therefore, with $|w|=(w_1^2+w_2^2)^{1/2}$,
\begin{equation}\label{eq:slow}
 a-\theta=-\frac{w_1+\mathsf g_1c}{\mathsf g_0+\mathsf g_1}+O(|w|^2+c^2),\qquad
 w_3=-\mathsf kc+\frac{\mathsf g_1w_1+\mathsf g_0w_2}{\mathsf g_0+\mathsf g_1}
 +O(|w|^2+c^2),
\end{equation}
and similarly for $b$; the factor of $c$ in $w_3$ is
$-(\mathsf g_0+\mathsf g_1)+(\mathsf g_0^2+\mathsf g_1^2)/(\mathsf g_0+\mathsf g_1)=-\mathsf k$.

We claim $|w|=O(c^2)$. Choose $T$ such that for $t\geq T$:
$\Lambda_a,\Lambda_b\geq\Lambda_\infty/2$,
$|\frac d{dt}\log c|\leq\Lambda_\infty/16$, $|w|\leq1$, $c\leq1$,
\eqref{eq:slow} holds, and $a\varrho_0+b\varrho_1\geq\Lambda_\infty/2$, so
that $\varepsilon(s)\leq\varepsilon(T)e^{-\Lambda_\infty(s-T)}$ for $s\geq T$
by \eqref{eq:chi}. Then, for constants $K',K''$ and after enlarging $T$,
\[
 \frac d{dt}|w|\leq-\frac{\Lambda_\infty}4|w|+K'\,c^2+K''\,\varepsilon\,c
\]
in the sense of upper right derivatives. Since
$c(s)\leq c(t)e^{\frac{\Lambda_\infty}{16}(t-s)}$ and
$e^{-\frac{\Lambda_\infty}{16}(t-T)}\leq c(t)/c(T)$ for $T\leq s\leq t$,
\[
 |w(t)|\leq\Bigl(\frac{c(t)}{c(T)}\Bigr)^4|w(T)|+\frac{8K'}{\Lambda_\infty}\,c(t)^2
 +K''c(t)\,\varepsilon(T)\int_T^te^{-\frac{3\Lambda_\infty}{16}(t-s)}e^{-\Lambda_\infty(s-T)}ds ,
\]
and the integral is at most
$\frac{16}{13\Lambda_\infty}e^{-\frac{3\Lambda_\infty}{16}(t-T)}\leq\frac{16}{13\Lambda_\infty}(c(t)/c(T))^3$.
This proves the claim. By \eqref{eq:reduced} and \eqref{eq:slow},
$\frac d{dt}\log c=-\mathsf kc+O(c^2)+O(\varepsilon c)$, so
$\frac d{dt}(1/c)=\mathsf k+O(c)+O(\varepsilon)$; then $c=O(1/t)$,
$\int^tc=O(\log t)$ and $1/c=\mathsf kt+O(\log t)$. With \eqref{eq:slow} this
is (ii).

(iii) By (ii), up to terms $O(c^2)$ and terms of order
$|\chi|\sqrt c$, which decay exponentially,
$C_0-\theta=\frac{\mathsf g_1c}{\mathsf g_0+\mathsf g_1}\diag(-1,1)$ and
$C_1-\theta=\frac{\mathsf g_0c}{\mathsf g_0+\mathsf g_1}\diag(1,-1)$. Moreover
$\cK_{\coh}=\tr C_0+\theta\log(\det h_0/\det h_1)=a+b+c+\varepsilon c-\theta\log(ab)$,
and $x-\theta\log x$ is stationary at $\theta$, so
$\cK_{\coh}=2\theta-2\theta\log\theta+c+O(c^2)+O(\varepsilon c)$; the constant
is $\cK_*$ by Theorem~\ref{thm:relaxation}.

(iv) By \eqref{eq:reduced} and (ii),
$\frac d{dt}\log(e^{\varphi t}p_1)=\mathsf g_0(a-\theta)+O(c^2)+O(\varepsilon c)
=-\frac{\mathsf k}2c+O(c^2)+O(\varepsilon c)=-\frac1{2t}+O(\frac{\log t}{t^2})$;
the error is integrable, so $e^{\varphi t}p_1\,t^{1/2}$ has a positive limit
$A$. Then $p_1'=ap_1$, $p_2=p_1'/c$ and $p_2'=bp_2$ give the other three. The
metrics are $g_i^*\diag g_i$ with the diagonal matrices of the data and
$g_i=\bigl(\begin{smallmatrix}1&\chi_i\\0&1\end{smallmatrix}\bigr)$, which
converge. For every pair $(g_0,g_1)$ of invertible matrices the map
$(h_0,h_1)\mapsto(g_0^*h_0g_0,\,g_1^*h_1g_1)$ is an isometry of the space of
pairs of metrics, and the $g_i(t)$ stay in a compact set; so the distance
from $\mathfrak h(0)$ is
$\bigl(\sum_i\log^2(e^{\varphi t}p_i)+\sum_i\log^2(e^{\varphi t}p_i')\bigr)^{1/2}+O(1)
=\log t+O(1)$.
\end{proof}

\begin{remark}\label{rem:semistable}
Apart from the limits $a,b\to\theta$ and $c\to0$, the proof uses of
$G_0,G_1$ only that they are smooth and increasing, with positive
derivatives at $\theta$, and vanish at $\theta$. For the downward gradient
flow of the Kempf--Ness function, for which these limits hold as well, the
same computation applies with $G_i(x)=x-\theta$, and gives the same exponents
with $\mathsf g_0=\mathsf g_1=\mathsf k=1$; asymptotics of such gradient flows
at semistable points are the subject of \cite{hkkp23}. For $\chi(0)\neq0$ the rate in (i) is
not attained for $\Lambda'=\Lambda_\infty$: the divided differences satisfy
$\varrho_i=\mathsf g_i+O(c^2)$, hence
$a\varrho_0+b\varrho_1=\Lambda_\infty-\mathsf kc+O(c^2)$ by (ii), and $|\chi(t)|$
is asymptotic to a constant multiple of $t\,e^{-\Lambda_\infty t}$. For
general semistable cohomology the induced metrics are determined up to a
bounded factor by \cite[Theorem~6.3 and Corollary~6.5]{wue26b}; whether the
factor converges is open.
\end{remark}

\subsection{Reversed charges}

\begin{theorem}[Reversed charges]\label{thm:reversed}
Let $\alpha_0<\alpha_1$. Along every trajectory, for
$t\geq t_{\mathrm{rev}}:=\max\bigl\{t_2,\ t_1+\frac4{\alpha_1}\log^+\frac{2\sqrt2\,c_\delta}{\alpha_1-\alpha_0}\bigr\}$,
\[
 q(t)\leq\bar N\,e^{-\frac12(\alpha_1-\alpha_0)(t-t_{\mathrm{rev}})} ,
\]
and $m(h(t))\to n_0|z_0|+n_1|z_1|$ as $t\to\infty$. The phases of the induced
metrics on $\ker D$ and $W$ converge to $\alpha_0$ and $\alpha_1$.
\end{theorem}

\begin{proof}
By \eqref{eq:adapted},
$\dot q=\tr(C_0\Phi_{HH})-\tr(C_1\Phi_W)
=\tr(C_0\phi_0)-\tr(C_1\phi_1)+\tr(C_0\delta_0)-\tr(C_1\delta_1)$. Since
$C_0,C_1\geq0$, $\phi_0=\arg(z_0\id+C_0)\leq\alpha_0$ and
$\phi_1=\arg(z_1\id-C_1)\geq\alpha_1$. Hence, for all charges,
$\dot q\leq-(\alpha_1-\alpha_0-\norm{\delta_0}-\norm{\delta_1})\,q$. By
Theorem~\ref{thm:cancellation}(i),(iv),
$\norm{\delta_0}+\norm{\delta_1}\leq\sqrt2\,c_\delta e^{-\frac{\alpha_1}4(t-t_1)}$
for $t\geq t_2$, which is at most $\frac12(\alpha_1-\alpha_0)$ for
$t\geq t_{\mathrm{rev}}$; and $q(t_{\mathrm{rev}})\leq\bar N$. Then $C_0,C_1\to0$, so
$\phi_0\to\alpha_0$, $\phi_1\to\alpha_1$, $\delta\to0$,
$m_H\to n_0|z_0|+n_1|z_1|$, and $|m-m_H|\to0$ by
Theorem~\ref{thm:cancellation}(iii).
\end{proof}

For reversed charges every representation with nonzero arrows is unstable;
its Harder--Narasimhan factors are the two semisimple representations
supported at one vertex each, of phases $\alpha_1>\alpha_0$, and Theorem~\ref{thm:reversed} says that the flow
converges to their direct sum. By the triangle inequality
$n_0|z_0|+n_1|z_1|>|Z|$ unless $z_0,z_1$ are parallel.

\begin{remark}\label{rem:unstable}
For ordered charges and cohomology that is not semistable,
Theorem~\ref{thm:cancellation} still gives the cancellation of the acyclic
part, with arrows bounded by $q\leq\bar N$ after the time $t_0$, and reduces
the long-time behaviour
to that of the phase flow of the cohomology representation, perturbed by the
exponentially decaying $\delta_0,\delta_1$. In \cite[\S4.1]{hkkp26} the
asymptotics of the flow on an object without harmonic metric are expected to
be controlled by a refinement of its Harder--Narasimhan filtration. For the
complexes of this paper this is proved in the following form
\cite[Theorem~6.11 and Corollary~6.5]{wue26b}: the induced metrics decay at
the exponential rates given by the phases of the Harder--Narasimhan
filtration of the cohomology representation, on the steps of that
filtration, and the metric induced on each factor is determined, up to a
bounded factor, by the theory for semistable representations. The simplest case is that in which all arrows of the
cohomology representation vanish: then $H_j=0$ for all time,
$\Omega_V$ is block diagonal with the block $z_0\id$ on $\ker D$, so
$\Phi_{HH}=\alpha_0$ exactly, and $\Phi_W=\arg(z_1\id-Y)\to\alpha_1$.
\end{remark}

\section{The phase potential}\label{sec:potential}

\subsection{The rate identity}

\begin{definition}\label{def:potential}
The \emph{phase potential} is
$\cK(h):=\sigma^{-1}\Ima\bigl(e^{-\ii\varphi}\cS_{\CC}(h)\bigr)$, and the
\emph{dissipation} is $\cD:=\cD_V+\cD_U+\cD_W$ with
\[
 \cD_\bullet:=\tr\Bigl(|\Omega_\bullet|\,(\Phi_\bullet-\varphi)\sin(\Phi_\bullet-\varphi)\Bigr).
\]
\end{definition}

For the symmetric charges and $n_0=n_1$, $\cK=-\Rea\cS_{\CC}$, that is
\begin{equation}\label{eq:Ksymmetric}
 \cK(h)=\log\det h_V-\log\det h_U-\log\det h_W+\tr\Ein(D^\dagger D)
 +\sum_j\tr\bigl(A_jR(D^\dagger D)A_j^\dagger\bigr).
\end{equation}
For a representation regarded as a complex in degree $0$ the potential is its
Kempf--Ness function for the stability parameter $(\theta_0,\theta_1)$; see
Proposition~\ref{prop:decomposition}. The function $\cK$ is a primitive of the
real one-form
$\xi\mapsto\sigma^{-1}\Ima\sum_\bullet\tr(e^{-\ii\varphi}\Omega_\bullet\xi_\bullet)$
on the space of metrics. The phase flow is not its gradient flow: the
gradient of $\cK$ for the invariant metric on the space of metrics has the
components $\sigma^{-1}\Ima(e^{-\ii\varphi}\Omega_\bullet)$, which do not differ
from the phases $\Phi_\bullet$ by a constant.

\begin{proposition}[Rate identity]\label{prop:rate}
Along the phase flow, for all charges in the upper half-plane,
\[
 \sigma\,\frac{d\cK}{dt}=-\cD .
\]
\end{proposition}

\begin{proof}
By \eqref{eq:primitivevariation} and \eqref{eq:flow},
$\frac d{dt}\cS_{\CC}=-\sum_\bullet\tr(\Omega_\bullet\Phi_\bullet)$. By
\eqref{eq:central}, $e^{-\ii\varphi}\sum_\bullet\tr\Omega_\bullet=|Z|$ is real,
so
\[
 \frac d{dt}\Ima(e^{-\ii\varphi}\cS_{\CC})
 =-\sum_\bullet\Ima\tr\bigl(e^{-\ii\varphi}\Omega_\bullet(\Phi_\bullet-\varphi)\bigr).
\]
By \eqref{eq:phase},
$e^{-\ii\varphi}\Omega_\bullet(\Phi_\bullet-\varphi)
=|\Omega_\bullet|\bigl[(\Phi_\bullet-\varphi)\cos(\Phi_\bullet-\varphi)
+\ii(\Phi_\bullet-\varphi)\sin(\Phi_\bullet-\varphi)\bigr]$, and the trace of a
product of two Hermitian matrices is real.
\end{proof}

The proof uses only \eqref{eq:primitivevariation}, \eqref{eq:flow} and
\eqref{eq:central}. For a representation in degree $0$ the endomorphism $\Omega$ is
normal and the statement follows from formulas displayed in
\cite[\S\S4.2, 5.1]{hkkp26}, although it is not stated there.

By Lemma~\ref{lem:eigenbasis}, in an orthonormal eigenbasis $(v_k)$ of
$\Phi_\bullet$ with eigenvalues $\phi_k$ and
$p_k=\inner{v_k}{|\Omega_\bullet|v_k}>0$,
\begin{equation}\label{eq:diagonal}
 \cD_\bullet=\sum_kp_k(\phi_k-\varphi)\sin(\phi_k-\varphi),\qquad
 p_k\sin(\phi_k-\varphi)=\inner{v_k}{\Ima\bigl(e^{-\ii\varphi}\Omega_\bullet\bigr)v_k}.
\end{equation}

\begin{proposition}\label{prop:degreezero}
$\cD_V\geq0$ and $\cD_W\geq0$, with equality on a block if and only if the
phase is $\varphi\id$ there.
\end{proposition}

\begin{proof}
For $\bullet=V,W$ the spectrum of $\Phi_\bullet$ lies in $(0,\pi)$ and
$\varphi\in(0,\pi)$, so $|\phi_k-\varphi|<\pi$ and
$(\phi_k-\varphi)\sin(\phi_k-\varphi)\geq0$, with equality only for
$\phi_k=\varphi$.
\end{proof}

The block $U$ of degree one is different: the spectrum of $\Phi_U-\varphi$
lies in $(-\pi-\varphi,-\varphi)$, and $x\sin x<0$ for
$x\in(-\pi-\varphi,-\pi)$.

\begin{proposition}\label{prop:degreeone}
Let $M=\sigma J+c_0e^{-L}$. For all charges in the upper half-plane,
\[
 \cD_U\geq\varphi\tr M-\pi\tr M_- .
\]
If $c_0\geq0$, then
\[
 \cD_U\geq\varphi c_0\tr e^{-L}+\sigma\bigl(\varphi\tr J-\pi\tr J_-\bigr).
\]
\end{proposition}

\begin{proof}
By \eqref{eq:coeff},
$\Ima(e^{-\ii\varphi}\Omega_U)=-\sigma J-|z_0|\sin(\alpha_0-\varphi)e^{-L}=-M$,
so by \eqref{eq:diagonal}, with $M_{kk}=\inner{v_k}{Mv_k}$,
$p_k\sin(\phi_k-\varphi)=-M_{kk}$ and
$\cD_U=\sum_k(\varphi-\phi_k)M_{kk}$. If $M_{kk}\geq0$, then
$\sin(\phi_k-\varphi)\leq0$, so $\varphi-\phi_k\in(\varphi,\pi]$ and the term
is at least $\varphi M_{kk}$. If $M_{kk}<0$, then
$\varphi-\phi_k\in(\pi,\pi+\varphi)$ and the term is at least
$(\pi+\varphi)M_{kk}$. In both cases it is at least
$\varphi M_{kk}-\pi(M_{kk})_-$. For any orthonormal basis
$M_{kk}\geq-(M_-)_{kk}$, hence $\sum_k(M_{kk})_-\leq\tr M_-$. This is the first
inequality. For the second, $\tr M_-=\max_{0\leq\Pi\leq\id}(-\tr(M\Pi))$ and
$c_0e^{-L}\geq0$, so $\tr M_-\leq\sigma\tr J_-$.
\end{proof}

\begin{remark}\label{rem:sharp}
As $e^{-L}\to0$ with $J$ fixed and invertible, $\Omega_U$ becomes Hermitian,
its phases tend to $0$ or $-\pi$, and the first inequality of
Proposition~\ref{prop:degreeone} becomes an equality. So the sign of $\cD_U$
is governed by the ratio $\tr J_-/\tr J$, which is the subject of
Theorem~\ref{thm:kernel}; this is the origin of the number
$\frac{2\pi}5=\pi\cdot\frac25$ in Definition~\ref{def:admissible}.
\end{remark}

\begin{theorem}[Lyapunov function]\label{thm:monotone}
Let the charges be admissible. Then at every state
\begin{equation}\label{eq:dissipationbound}
 \cD\geq\cD_U\geq\varphi c_0\tr e^{-L}
 +\sigma\bigl(\varphi-\tfrac{2\pi}5\bigr)_+\tr J>0,
\end{equation}
and the phase potential is strictly decreasing along the phase flow: it is a
strict Lyapunov function. For the symmetric charges and $n_0=n_1$,
\[
 \frac{d\cK}{dt}=-\cD,\qquad
 \cD\geq\frac\pi{10}\tr J+\frac\pi2\tr e^{-L}.
\]
\end{theorem}

\begin{proof}
Propositions~\ref{prop:rate}, \ref{prop:degreezero} and \ref{prop:degreeone},
together with Theorem~\ref{thm:kernel}: $\tr J_-\leq\frac25\tr J$ and
$\tr J\geq0$, since $DCD^\dagger\geq0$. If $r=1$ then $J$ is a nonnegative
number, by \eqref{eq:schurform} and Lemma~\ref{lem:mformula}, and $J_-=0$.
Here $L>0$ because $D$ is surjective, and $c_0>0$.
\end{proof}

No stability assumption on the cohomology representation and no assumption on
the size of $D$ or of the maps $A_jD^+$ enters Theorem~\ref{thm:monotone}.
Some condition on the charges is needed: see
Propositions~\ref{prop:reversed} and~\ref{prop:orderedfailure}.

\subsection{Decomposition and infimum}
In this subsection the charges are ordered. Let $H=\ker D$, with the induced
metric $h_0=h_V|_H$, and $h_1=h_W$. Fix reference bases adapted to
$H\subset V$: a basis of $H$ followed by vectors whose classes form a basis
of $V/H$. Let $\bar D:V/H\to U$ be the isomorphism induced by $D$ and
$c_D=\log|\det\bar D|^2$, the determinant taken in the reference bases. Let
$\cK_{\coh}$ be the Kempf--Ness function \eqref{eq:KN} of the cohomology
representation, in the reference basis of $H$ and that of $W$, and
$\cK_*$ its infimum.

Let $E_1(b)=\int_1^\infty e^{-xb}\,dx/x$ be the exponential integral and
$\gamma$ Euler's constant.

\begin{proposition}[Decomposition]\label{prop:decomposition}
\[
 \cK=\cK_{\coh}(h_0,h_1)+\theta_0\Bigl(c_D+r\gamma+\sum_{i=1}^rE_1(\ell_i)\Bigr)
 +\tr Y .
\]
\end{proposition}

\begin{proof}
$B$ vanishes on $H$ and equals $F^\dagger F$ on $H^\perp$, so
$\tr\Ein(B)=\sum_i\Ein(\ell_i)$, and $R(B)$ is the identity on $H$, so
$\tr(CR(B))=\sum_j\tr(A_jR(B)A_j^\dagger)=q+\tr Y$, using
$T_jR(F^\dagger F)T_j^\dagger=X_j(\id-e^{-L})X_j^\dagger$. For the Gram matrix
in the adapted reference basis, $\det h_V=\det h_0\cdot\det h_{V/H}$ with
$h_{V/H}$ the quotient metric, which is the metric of $H^\perp\cong V/H$. In
the reference bases of $V/H$ and $U$ the endomorphism $F^\dagger F$ of
$H^\perp$ has matrix $h_{V/H}^{-1}[\bar D]^*h_U[\bar D]$, so
$\prod_i\ell_i=|\det\bar D|^2\det h_U/\det h_{V/H}$. Hence
$\log\det h_V-\log\det h_U=\log\det h_0+c_D-\sum_i\log\ell_i$, and with
$\Ein(b)=\gamma+\log b+E_1(b)$ for $b>0$,
\[
 \cS_{\CC}=z_0\Bigl(\log\det h_0+c_D+r\gamma+\sum_iE_1(\ell_i)\Bigr)
 +z_1\log\det h_1-q-\tr Y .
\]
Finally $\Ima(e^{-\ii\varphi}z_0)=c_0$, $\Ima(e^{-\ii\varphi}z_1)=-c_1$ and
$\Ima(-e^{-\ii\varphi})=\sigma$.
\end{proof}

The identity
$\log\det h_V-\log\det h_U+\log\det(F^\dagger F)=\log\det h_0+c_D$ in this
proof is the finite-dimensional comparison between the metric induced on the
determinant of the cohomology and the metric of the cohomology, with torsion
$\log\det(F^\dagger F)$; compare \cite[Theorem~1.1]{bz92} for complexes of
real vector spaces. The constant $r\gamma$ arises because $\Ein$ is the heat
integral cut off at time $1$ rather than its zeta-regularization
(Remark~\ref{rem:normalization}).

\begin{corollary}\label{cor:infimum}
$\inf_h\cK=\cK_*+\theta_0(c_D+r\gamma)$, and the infimum is not attained.
\end{corollary}

\begin{proof}
The terms $\theta_0\sum_iE_1(\ell_i)$ and $\tr Y$ are nonnegative and the
first is positive, which gives ``$\geq$'' and non-attainment. For the
converse fix a complement $H'$ of $H$ in $V$ with a metric, a metric $h_U^0$,
and take $h_V=h_0\oplus h_{H'}$, $h_U=\epsilon^{-1}h_U^0$. As $\epsilon\to0$
all $\ell_i\to\infty$, so $E_1(\ell_i)\to0$ and
$\tr Y\leq\sum_j\norm{T_j}_{\HS}^2/\min_i\ell_i\to0$, while $\cK_{\coh}$
depends only on $(h_0,h_1)$.
\end{proof}

\begin{lemma}\label{lem:qbound}
For every pair $(h_0,h_1)$ with $q>0$,
\[
 q-\vartheta-\vartheta\log\frac q\vartheta\leq\cK_{\coh}(h_0,h_1)-\cK_* .
\]
\end{lemma}

\begin{proof}
Replacing $h_1$ by $e^xh_1$ changes $\cK_{\coh}$ to
$e^xq+(\cK_{\coh}-q)-\vartheta x$; at $e^x=\vartheta/q$ this is
$\vartheta+\cK_{\coh}-q+\vartheta\log(q/\vartheta)\geq\cK_*$.
\end{proof}

The function $Q\mapsto Q-\vartheta-\vartheta\log(Q/\vartheta)$ is convex,
vanishes only at $Q=\vartheta$ and tends to $+\infty$ as $Q\to0$ and as
$Q\to\infty$.

\subsection{Bounds through the potential}
Let the charges be admissible and the cohomology representation semistable.
Then the Lyapunov function gives bounds that do not involve the arrow energy.
Put
\begin{equation}\label{eq:gap}
 \Delta_0:=\cK(h(0))-\inf_h\cK
 =\cK_{\coh}(0)-\cK_*+\theta_0\sum_iE_1(\ell_i(0))+\tr Y(0)\in(0,\infty),
\end{equation}
and let $Q_\infty\geq\vartheta$ be the larger root of
$Q-\vartheta-\vartheta\log(Q/\vartheta)=\Delta_0$.

\begin{proposition}[Bounds through the potential]\label{prop:apriori}
Let the charges be admissible and $\cK_*>-\infty$. For all $t\geq0$:
\begin{enumerate}
\item[(i)] $0\leq\cK_{\coh}(t)-\cK_*\leq\Delta_0$ and $q(t)\leq Q_\infty$;
\item[(ii)] $\theta_0\sum_iE_1(\ell_i(t))\leq\Delta_0$ and
 $\tr Y(t)\leq\Delta_0$, hence $s(t)\leq\Delta_0/(1-e^{-\lambda(t)})$;
\item[(iii)] $\int_0^\infty\cD\,dt\leq\sigma\Delta_0$;
\item[(iv)] $\int_0^\infty e^{-\lambda(t)}dt\leq\dfrac{\Delta_0}{\varphi\theta_0}$
 and
 $\lambda(t)\geq\max\Bigl\{\lambda(0),\ \log\dfrac{\varphi\theta_0t}{\Delta_0}\Bigr\}$.
\end{enumerate}
In particular the conclusions of Theorem~\ref{thm:cancellation} hold with the
constants \eqref{eq:constants} computed from
$\bar N':=\max\{Q_\infty,\Delta_0\}$ in place of $\bar N$, and with $t_1$
replaced by $t_1':=\frac{\Delta_0}{\varphi\theta_0}e^{\Lambda_1'}$, where
$\Lambda_1'=\max\{1,8\pi\bar N'/(\alpha_1y_0)\}$.
\end{proposition}

\begin{proof}
By Proposition~\ref{prop:decomposition} and Corollary~\ref{cor:infimum},
$\cK-\inf\cK=(\cK_{\coh}-\cK_*)+\theta_0\sum_iE_1(\ell_i)+\tr Y$ is a sum of
three nonnegative terms, and it is at most $\Delta_0$ by
Theorem~\ref{thm:monotone}. This gives (i), with Lemma~\ref{lem:qbound}, and
(ii), by \eqref{eq:Y}. (iii) is the integral of $\sigma\dot\cK=-\cD$. For
(iv), \eqref{eq:dissipationbound} gives
$\cD\geq\varphi c_0\tr e^{-L}\geq\varphi\sigma\theta_0e^{-\lambda}$, and
$\lambda$ is nondecreasing (Proposition~\ref{prop:imported}(i)), so
$te^{-\lambda(t)}\leq\int_0^te^{-\lambda}$. For the last statement, by (i)
and (ii) $q\leq\bar N'$ and $\tr Y\leq\bar N'$ for all $t\geq0$, and by (iv)
$\lambda\geq\Lambda_1'$ for $t\geq t_1'$. The proof of
Theorem~\ref{thm:cancellation} uses Theorems~\ref{thm:energy}
and~\ref{thm:growth} only to obtain these three facts, with $\bar N$,
$\Lambda_1$, $t_1$ in place of $\bar N'$, $\Lambda_1'$, $t_1'$; the rest of
that proof applies unchanged.
\end{proof}

For the symmetric charges and $n_0=n_1$, (iv) reads
$\lambda(t)\geq\log(\pi t/2\Delta_0)$.

\section{The infimum: quasi-isomorphic complexes and moduli}\label{sec:moduli}

The charges are ordered. This section identifies the infimum of the phase
potential and shows that, up to an explicit constant, it is the same for
quasi-isomorphic complexes of the class \eqref{eq:complex}.

\subsection{Volume data}
A \emph{volume datum} for a complex $\mathsf E$ is a triple
$\varpi=(\varpi_V,\varpi_U,\varpi_W)$ of nonzero elements of the top exterior
powers $\det V,\det U,\det W$. If these are the wedges of the reference bases,
then by \eqref{eq:primitive}
\begin{equation}\label{eq:Kvol}
 \cK(h;\varpi)=\theta_0\bigl(\log\norm{\varpi_V}^2_{h_V}-\log\norm{\varpi_U}^2_{h_U}\bigr)
 -\theta_1\log\norm{\varpi_W}^2_{h_W}+\theta_0\tr\Ein(B)+\tr\bigl(CR(B)\bigr),
\end{equation}
and we use \eqref{eq:Kvol} as the definition of $\cK$ for an arbitrary volume
datum. Let $H=\ker D$. There is a canonical isomorphism
\[
 \iota_D:\det V\otimes(\det U)^{-1}\to\det H,\qquad
 (v_1\wedge\dots\wedge v_{n_0+r})\otimes(Dv_{n_0+1}\wedge\dots\wedge Dv_{n_0+r})^{-1}
 \mapsto v_1\wedge\dots\wedge v_{n_0},
\]
for any basis $(v_i)$ of $V$ whose first $n_0$ vectors span $H$. Put
$\bar\varpi=(\iota_D(\varpi_V\otimes\varpi_U^{-1}),\varpi_W)$, a volume datum
for the cohomology representation $\mathcal H=(H\to W)$, and
\begin{gather*}
 \cK_{\coh}(\mathcal H;h_0,h_1;\bar\varpi)=q+\theta_0\log\norm{\iota_D(\varpi_V\otimes\varpi_U^{-1})}^2_{h_0}
 -\theta_1\log\norm{\varpi_W}^2_{h_1},\\
 \cK_*(\mathcal H;\bar\varpi)=\inf_{h_0,h_1}\cK_{\coh}.
\end{gather*}

\begin{lemma}\label{lem:determinant}
For all metrics, with $h_0=h_V|_H$ and $\det{}'$ the product of the nonzero
eigenvalues,
\[
 \log\norm{\varpi_V}^2_{h_V}-\log\norm{\varpi_U}^2_{h_U}+\log\det{}'(D^\dagger D)
 =\log\norm{\iota_D(\varpi_V\otimes\varpi_U^{-1})}^2_{h_0}.
\]
\end{lemma}

\begin{proof}
Take an $h_V$-orthonormal basis $(v_i)$ with $v_1,\dots,v_{n_0}\in H$ and an
$h_U$-orthonormal basis $(u_i)$, let $F$ be the matrix of $D|_{H^\perp}$, and
write $\varpi_V=a_Vv_1\wedge\dots\wedge v_{n_0+r}$,
$\varpi_U=a_Uu_1\wedge\dots\wedge u_r$. Then
$\iota_D(\varpi_V\otimes\varpi_U^{-1})=(a_V/a_U)\det F\cdot v_1\wedge\dots\wedge v_{n_0}$,
and $|\det F|^2=\det{}'(D^\dagger D)$.
\end{proof}

This is the comparison already used in the proof of
Proposition~\ref{prop:decomposition}, where $c_D$ is the logarithm of the
squared modulus of $\iota_D$ in the reference bases.

\begin{theorem}[Infimum of the potential]\label{thm:infimum}
Let the charges be ordered. For every complex and volume datum,
\begin{gather*}
 \cK(h;\varpi)=\cK_{\coh}(\mathcal H;h_0,h_1;\bar\varpi)
 +\theta_0\bigl[r\gamma+\tr{}'E_1(D^\dagger D)\bigr]+\tr Y,\\
 \inf_h\cK(h;\varpi)=\cK_*(\mathcal H;\bar\varpi)+\theta_0\gamma\,r,
\end{gather*}
where $\tr{}'$ is the trace over the nonzero spectrum; the infimum is not
attained. Consequently, let $\mathsf E,\mathsf E'$ be two complexes of the
class \eqref{eq:complex} with semistable cohomology and
$f=(f_V,f_U,f_W):\mathsf E\to\mathsf E'$ a quasi-isomorphism of complexes of
representations, that is, a morphism inducing an isomorphism
$\mathcal H\to\mathcal H'$ of the cohomology representations. If the volume
data correspond, $\bar\varpi'=f_*\bar\varpi$, then
\[
 \inf_{h'}\cK_{\mathsf E'}(h';\varpi')-\inf_h\cK_{\mathsf E}(h;\varpi)=\theta_0\gamma\,(r'-r).
\]
\end{theorem}

\begin{proof}
The first identity is Proposition~\ref{prop:decomposition} with
Lemma~\ref{lem:determinant} and
\[
 \tr\Ein(D^\dagger D)=\log\det{}'(D^\dagger D)+r\gamma+\tr{}'E_1(D^\dagger D);
\]
the infimum is Corollary~\ref{cor:infimum}. The function $\cK_*$ is
transported by isomorphisms of representations with volume data.
\end{proof}

We note three special cases. For an invertible triple $g=(g_V,g_U,g_W)$ and fixed
$\varpi$,
$\inf\cK_{g\cdot\mathsf E}=\inf\cK_{\mathsf E}-\theta_0\log|\det g_V/\det g_U|^2+\theta_1\log|\det g_W|^2$.
A shear $g_V=\left(\begin{smallmatrix}\id&S\\0&\id\end{smallmatrix}\right)$ in
a splitting $V=H\oplus H'$, which leaves $D$ and the arrows on $H$ unchanged
and replaces the arrows $T_j$ on $H'$ by $T_j-A_j|_HS$, does not change the
infimum. Adding an acyclic summand $H''\xrightarrow{F''}U''$ with arbitrary
arrows $H''\to W$ adds
$\theta_0[\gamma\dim U''+\log|\det F''|^2]$, independently of those arrows.
In \cite[\S8.1]{wue26} the shear and the removal of orthogonal acyclic
summands were compared for fixed metrics on $V,U,W$ and $\lambda$ large.

\begin{remark}\label{rem:normalization}
The \emph{normalized potential} $\cK-\theta_0\gamma\,r$ has infimum
$\cK_*(\mathcal H;\bar\varpi)$, the infimum of the Kempf--Ness function of the
cohomology representation, for every complex with that cohomology. The
constant reflects the choice of regularization: for $b>0$,
$\log b=\int_0^1(1-e^{-xb})\frac{dx}x-\int_1^\infty e^{-xb}\frac{dx}x-\gamma$,
and $\Ein(b)$ is the first integral alone. The primitive of $\Omega$
is determined on the metrics of one complex only up to an additive constant,
and the normalized potential fixes these constants compatibly.
\end{remark}

\subsection{The potential on moduli}
Let $\mathcal R=\operatorname{Hom}(\CC^{n_0},\CC^{n_1})^k$ be the space of
representations of dimension vector $(n_0,n_1)$, with the action of
$\mathrm G=\mathrm{GL}_{n_0}\times\mathrm{GL}_{n_1}$,
$b_j\mapsto g_1b_jg_0^{-1}$, and $\mathrm K=\mathrm U(n_0)\times\mathrm U(n_1)$.
Put
\[
 \varkappa(b)=\inf_{g\in\mathrm G}\Bigl[\norm{g\cdot b}^2+\theta_0\log|\det g_0|^2
 -\theta_1\log|\det g_1|^2\Bigr],
\]
let $\mathcal R^{\mathrm s}$ be the set of stable points,
$\mathcal R_0=\{C_0=\theta_0\id,\ C_1=\theta_1\id\}\cap\mathcal R^{\mathrm s}$
the zeros of the moment map in it, and
$\mathcal M^{\mathrm s}=\mathcal R^{\mathrm s}/\mathrm G=\mathcal R_0/\mathrm K$
the moduli space of stable representations \cite{kin94}, with the K\"ahler
form $\omega_\theta$ obtained by reduction of
$\omega_{\mathrm{flat}}=\ii\partial\bar\partial\norm b^2$: its pullback to
$\mathcal R_0$ is the restriction of $\omega_{\mathrm{flat}}$. Let
$\mathrm{pr}:\mathcal R^{\mathrm s}\to\mathcal M^{\mathrm s}$ be the
projection. The stable locus may be empty, for instance for $k=2$ and
$n_0=n_1\geq2$; the statements about it are then vacuous.

\begin{lemma}\label{lem:kappa}
The function $\varkappa$ has the following properties.
\begin{enumerate}
\item[(a)] $\varkappa(g\cdot b)=\varkappa(b)-\theta_0\log|\det g_0|^2+\theta_1\log|\det g_1|^2$.
\item[(b)] $\varkappa$ is real-analytic on $\mathcal R^{\mathrm s}$ and
 $\ii\partial\bar\partial\varkappa=\mathrm{pr}^*\omega_\theta$ there.
\end{enumerate}
\end{lemma}

\begin{proof}
(a) Substitute $g'g$ in the infimum. (b) Let $b_0\in\mathcal R_0$ and, for a
pair $\zeta=(\zeta_0,\zeta_1)$ of Hermitian matrices, let
$f(\zeta,b)=\norm{e^\zeta\cdot b}^2+2\theta_0\tr\zeta_0-2\theta_1\tr\zeta_1$,
so $\varkappa(b)=\inf_\zeta f(\zeta,b)$. Then $\partial_\zeta f(0,b_0)=0$, and
$\partial_\zeta^2f(0,b_0)[\zeta,\zeta]=4\norm{\zeta\cdot b_0}^2$, with
$\zeta\cdot b=(\zeta_1b_j-b_j\zeta_0)_j$ the infinitesimal action, is positive
definite modulo the scalars $(c,c)$, on which $f$ is constant. By the
implicit function theorem there is a critical point $\zeta(b)$ of
$f(\cdot,b)$, real-analytic near $b_0$. It is the global minimizer:
$b'=e^{\zeta(b)}\cdot b$ lies in $\mathcal R_0$, the function
$x\mapsto f(x\zeta',b')$ is convex with zero derivative at $0$ for every
$\zeta'$, and every element of $\mathrm G$ is $ue^{\zeta'}$ with
$u\in\mathrm K$. For a holomorphic curve $b(w)$ with $b(0)=b_0$ and
$\dot b=b'(0)$, differentiating the minimum gives
$\partial_w\partial_{\bar w}\varkappa
=f_{w\bar w}-f_{w\zeta}f_{\zeta\zeta}^{-1}f_{\zeta\bar w}$ with
$f_{w\bar w}=\norm{\dot b}^2$,
$f_{w\zeta}[\zeta]=2\inner{\zeta\cdot b_0}{\dot b}$ and
$f_{\zeta\zeta}[\zeta,\zeta']=4\Rea\inner{\zeta\cdot b_0}{\zeta'\cdot b_0}$. At
a point of $\mathcal R_0$ the $\mathrm K$-orbit is isotropic, so the real
space $\{\zeta\cdot b_0\}$ is orthogonal to its multiple by $\ii$, and a real
orthonormal basis of it is a unitary basis of the tangent space
$\mathfrak g\cdot b_0$ of the complex orbit. Hence
\[
 \partial_w\partial_{\bar w}\varkappa(b(w))\big|_{w=0}
 =\norm{\dot b}^2-\norm{\Pi_{\mathfrak g\cdot b_0}\dot b}^2,
\]
the squared length of the component of $\dot b$ orthogonal to the complex
orbit, which is the reduced form evaluated on $d\,\mathrm{pr}(\dot b)$. For a
general stable point: every stable orbit meets $\mathcal R_0$
\cite{kn79,kin94}, and for fixed $g$ the translation $b\mapsto g\cdot b$ is a
biholomorphism of $\mathcal R^{\mathrm s}$ commuting with $\mathrm{pr}$ and
changing $\varkappa$ by a constant, by (a); so both sides of the identity are
invariant under it.
\end{proof}

Lemma~\ref{lem:kappa}(b) is the classical description of the K\"ahler
potential of a K\"ahler quotient through the shifted ambient potential; for
tori see \cite{bg04}. We include the short proof for the reader's
convenience.

Let now $\mathcal P$ be a complex manifold with holomorphic maps
$p\mapsto D(p),A_j(p)$ (matrices in fixed bases), $D(p)$ surjective, $\varpi$
the volume datum of the bases, $\mathsf E(p)$ the corresponding complexes,
$\cK_\infty(p)=\inf_h\cK_{\mathsf E(p)}(h;\varpi)$, and $\mathcal P^{\mathrm s}$
the open set where the cohomology representation is stable.

\begin{theorem}[The potential on moduli]\label{thm:moduli}
Let the charges be ordered. On $\mathcal P^{\mathrm s}$ the function
$\cK_\infty$ is real-analytic and
\[
 \ii\partial\bar\partial\cK_\infty=\mathrm{coh}^*\omega_\theta,
\]
where $\mathrm{coh}:\mathcal P^{\mathrm s}\to\mathcal M^{\mathrm s}$ sends $p$
to the class of the cohomology representation of $\mathsf E(p)$. In
particular the form is the same for all holomorphic families of complexes of
the class \eqref{eq:complex} with the same family of cohomology
representations, and it vanishes in every direction along which the
isomorphism class of the cohomology representation is constant.
\end{theorem}

\begin{proof}
$\ker D(p)$ is a holomorphic subbundle. In a local holomorphic frame
$(e_i(p))$ of it let $b(p)\in\mathcal R$ be the matrices of
$A_j(p)|_{\ker D(p)}$; then $b$ is holomorphic and
$\mathrm{coh}=\mathrm{pr}\circ b$. Write
$\iota_{D(p)}(\varpi_V\otimes\varpi_U^{-1})=a(p)\,e_1(p)\wedge\dots\wedge e_{n_0}(p)$
with $a$ holomorphic and nowhere zero. By Theorem~\ref{thm:infimum},
$\cK_\infty(p)=\varkappa(b(p))+\theta_0\log|a(p)|^2+\theta_0\gamma r$, and the
last two terms are pluriharmonic. Apply Lemma~\ref{lem:kappa}(b).
\end{proof}

\begin{remark}\label{rem:levi}
For fixed metrics on $V,U,W$ the function $(A,D)\mapsto\cK(h)$ is not
plurisubharmonic: for the symmetric charges it has a negative Levi direction
along a family of complexes with constant cohomology representation
\cite[Proposition~7.5]{wue26}. Theorem~\ref{thm:moduli} concerns the function
after minimizing over the metrics, where the infimum is not attained; along
such a family it is pluriharmonic.
\end{remark}

\subsection{Convergence along the flow in families}
For a continuous family of initial metrics $h^{\mathrm{in}}(p)$ let $h(t;p)$ be
the trajectory of the phase flow of $\mathsf E(p)$ starting at
$h^{\mathrm{in}}(p)$, and $\cK_t(p)=\cK_{\mathsf E(p)}(h(t;p);\varpi)$.

\begin{theorem}[Convergence in families]\label{thm:families}
Let the charges be ordered.
\begin{enumerate}
\item[(i)] If the cohomology representation of $\mathsf E(p)$ is semistable,
 then $\cK_t(p)\to\cK_\infty(p)$ as $t\to\infty$, although the infimum is not
 attained and no zero of the moment map need exist. For admissible charges
 $\cK_t(p)$ is strictly decreasing in $t$.
\item[(ii)] On compact subsets of $\mathcal P^{\mathrm s}$,
 $0<\cK_t-\cK_\infty\leq Ce^{-ct}$ with $C,c>0$ depending only on the subset.
\item[(iii)] On $\mathcal P^{\mathrm s}$ all parameter derivatives of $\cK_t$
 converge to those of $\cK_\infty$ in the sense of distributions; in
 particular $\ii\partial\bar\partial\cK_t\to\mathrm{coh}^*\omega_\theta$
 weakly.
\end{enumerate}
\end{theorem}

\begin{proof}
(i) By Theorem~\ref{thm:infimum},
$\cK_t-\cK_\infty=\mathcal E(t)+\theta_0\tr{}'E_1(D^\dagger D)+\tr Y$ with
$\mathcal E=\cK_{\coh}-\cK_*\geq0$ as in Section~\ref{sec:convergence}. The
second term is at most $\theta_0rE_1(\lambda(t))$ and the third at most
$s(t)$; both tend to zero by
Theorem~\ref{thm:cancellation}, and $\mathcal E(t)\to0$ by
Theorem~\ref{thm:relaxation}. Monotonicity for admissible charges is
Theorem~\ref{thm:monotone}; positivity of the difference is
Corollary~\ref{cor:infimum}.

(ii) For $t\geq t_1$, $\tr Y\leq s\leq s_1e^{-\frac{\alpha_1}2(t-t_1)}$ and
$\theta_0\tr{}'E_1\leq\theta_0re^{-\lambda}\leq\theta_0re^{\mathfrak c_3}e^{-\frac{\alpha_1}4(t-t_1)}$
by Theorem~\ref{thm:cancellation} and \eqref{eq:E1bounds}, and
$\mathcal E(t)\leq\mathfrak C_2e^{-2\beta t}$ by the proof of
Theorem~\ref{thm:convergence}. All constants there and in
\eqref{eq:constants} are continuous functions of $N(0)$, $\lambda(0)$,
$\tr E_1(L)|_{t=0}$, $\mathcal E(0)$ and $\nu_*$, which are continuous in $p$
for continuous initial metrics; and $\nu_*(p)$ is positive on
$\mathcal P^{\mathrm s}$, being the least eigenvalue of the Hessian form of
Lemma~\ref{lem:gap} at a zero of the moment map, which depends continuously
on $p$ modulo $\mathrm K$ by the proof of Lemma~\ref{lem:kappa}(b), on the
orthogonal complement of the scalars.

(iii) Locally uniform convergence of continuous functions implies convergence
of all derivatives in the sense of distributions.
\end{proof}

\begin{corollary}\label{cor:continuity}
In Theorem~\ref{thm:convergence} the limit $\mathfrak h_\infty$ depends continuously on
the initial metric. Along a continuous family of complexes whose cohomology
representations stay stable, with continuous initial metrics, it depends
continuously on the parameter, in continuous local frames of $\ker D$.
\end{corollary}

\begin{proof}
The constants of Theorem~\ref{thm:convergence} are locally uniform in such
families: $N(0)$, $t_2$ and $\mathcal E(0)$ are continuous, and $\nu_*$ is
continuous and positive on the stable locus, as shown in the proof of
Theorem~\ref{thm:families}(ii). So $d(\mathfrak h(t),\mathfrak h_\infty)\leq\mathfrak Ce^{-\beta t}$ uniformly,
and $\mathfrak h(T)$ depends continuously on the data for each fixed $T$.
\end{proof}


We do not know whether the parameter derivatives of $\cK_t$ converge
pointwise. Outside the stable locus the infimum need not be differentiable.
For example let $k=2$, $n_0=n_1=2$ (so $\theta_0=\theta_1$) and
$b=(\id_2,\left(\begin{smallmatrix}0&a\\a'&0\end{smallmatrix}\right))$ with
$a,a'\neq0$. This representation is isomorphic to the direct sum of the two
representations $(1,\pm c)$ of dimension vector $(1,1)$, $c^2=aa'$, by a
simultaneous conjugation, which does not change $\varkappa$ because
$\theta_0=\theta_1$; and on a direct sum of representations $(1,c_i)$ the
infimum is attained at a diagonal zero of the moment map. Hence
\[
 \varkappa(b)=2\theta_0+2\theta_0\log\frac{1+|a|\,|a'|}{\theta_0}\qquad(a,a'\neq0),
\]
which has no differentiable extension to $a=0$, where the representation is
semistable and not polystable.

\section{Mass, extreme phases and fixed points}\label{sec:expectations}

In \cite[\S4.1]{hkkp26} the structures of a categorical K\"ahler metric are
listed as axioms that ``should be read more as a collection of desiderata
rather than as a final definition''. For the flow the text reads: ``We expect
the following quantities to vary monotonically under this action: $m(X,h)$
decreases, $\phi^-(X,h)$ increases, and $\phi^+(X,h)$ decreases. Moreover,
the fixed points of the action on $\mathrm{Met}(X)/$rescaling are precisely
the harmonic metrics.'' Here $\phi^\pm$ are the ends of the support of the
mass measure, which is defined in \cite[\S4.2]{hkkp26} by
$\int f\,d\mu_{X,h}=\Omega\bigl(e^{-\sqrt{-1}\arg\Omega}f(\arg\Omega/\pi)\bigr)$;
there $\Omega$ is the functional $a\mapsto\tr(\Omega a)$, so that the right
side is $\tr\bigl(|\Omega|\,f(\arg\Omega/\pi)\bigr)$. We write $\nu_h$ for this
measure, since $\mu$ denotes the moment map. For
the endomorphism $\Omega$ of the complex \eqref{eq:complex}, the polar phase
$\Phi=\Phi_V\oplus\Phi_U\oplus\Phi_W$ and phases measured in radians, this is
the measure
\[
 \nu_h=\sum_kp_k\,\delta_{\phi_k},\qquad p_k=\inner{v_k}{|\Omega|v_k}>0,
\]
for an orthonormal eigenbasis $(v_k)$ of $\Phi$ with eigenvalues $\phi_k$;
equivalently $\int f\,d\nu_h=\tr(|\Omega|f(\Phi))$. Its total mass is $m(h)$,
and
\[
 \phi_+=\max\bigl(\spec\Phi_V\cup\spec\Phi_W\bigr),\qquad
 \phi_-=\min\spec\Phi_U .
\]
A metric is harmonic if $\Phi$ is a scalar. In \cite[Remark~5.15]{hkkp26} the
metrics of an object are the metrized complexes isomorphic to it in the
derived category. The statements of this section concern the metrics on a
given complex, which the flow preserves. For these we examine whether the
mass is nonincreasing, whether $\phi_+$ is nonincreasing and $\phi_-$
nondecreasing, and whether the fixed points modulo rescaling are the harmonic
metrics. For a representation all of this holds
(Proposition~\ref{prop:representations}). For the complexes
\eqref{eq:complex} the mass and both extreme phases can move in the opposite
direction (Propositions~\ref{prop:mass} and~\ref{prop:extreme}), and there
are neither fixed points nor harmonic metrics on the complex
(Proposition~\ref{prop:fixed}); what converges is the mass measure
(Proposition~\ref{prop:measure}).

\subsection{Representations}
Let a representation of the generalized Kronecker quiver with metrics be
given, regarded as a complex in degree zero, with dimension vector
$(n_0,n_1)$, $n_0,n_1\geq1$. In orthonormal frames let $H_j$ be the matrices
of the arrows, $C_0=\sum_jH_j^\dagger H_j$ and $C_1=\sum_jH_jH_j^\dagger$. Its
endomorphism $\Omega=(z_0+C_0)\oplus(z_1-C_1)$ is normal, with phases
$\phi_0=\arg(z_0+C_0)$ and $\phi_1=\arg(z_1-C_1)$, and the flow is
$\dot H_j=\frac12(H_j\phi_0-\phi_1H_j)$.

\begin{proposition}\label{prop:representations}
For a representation and all charges in the upper half-plane, along the phase
flow:
\begin{enumerate}
\item[(i)] the mass is nonincreasing;
\item[(ii)] the largest phase is nonincreasing and the smallest phase is
 nondecreasing;
\item[(iii)] the fixed points on metrics modulo rescaling are the metrics
 with $\phi_0=\varphi\,\id$ and $\phi_1=\varphi\,\id$, that is, the harmonic
 metrics.
\end{enumerate}
\end{proposition}

\begin{proof}
From the flow,
$\dot C_0=\phi_0C_0-\sum_jH_j^\dagger\phi_1H_j$ and
$\dot C_1=\sum_jH_j\phi_0H_j^\dagger-\phi_1C_1$, because $\phi_0$ commutes with
$C_0$ and $\phi_1$ with $C_1$.

(i) This is \cite[\S6.2]{wue26}: the derivative of $x\mapsto|z_0+x|$ is
$\cos\arg(z_0+x)$ and that of $x\mapsto|z_1-x|$ is $-\cos\arg(z_1-x)$, so in
eigenbases of $\phi_0$ (eigenvalues $a_i$) and $\phi_1$ (eigenvalues $b_l$)
\[
 \dot m=\tr(\cos\phi_0\,\dot C_0)-\tr(\cos\phi_1\,\dot C_1)
 =\sum_{j,i,l}|(H_j)_{li}|^2(\cos a_i-\cos b_l)(a_i-b_l)\leq0,
\]
because the cosine decreases on $(0,\pi)$.

(ii) We have $\phi_0=f_0(C_0)$ and $\phi_1=f_1(C_1)$ with
$f_0(x)=\arg(z_0+x)$ decreasing and $f_1(x)=\arg(z_1-x)$ increasing on
$[0,\infty)$. Let $\bar\phi_0=f_0(\lambda_{\min}(C_0))$ and
$\bar\phi_1=f_1(\lambda_{\max}(C_1))$ be the largest eigenvalues of $\phi_0$
and $\phi_1$. For a unit eigenvector $v$ of $C_0$ for its least eigenvalue
$c$,
$\inner v{\dot C_0v}=\bar\phi_0c-\sum_j\inner{H_jv}{\phi_1H_jv}
\geq(\bar\phi_0-\bar\phi_1)c$. The right lower derivative of
$\lambda_{\min}(C_0)$ is the minimum of $\inner v{\dot C_0v}$ over such $v$;
so whenever $\bar\phi_0\geq\bar\phi_1$, the upper right derivative of
$\bar\phi_0$ is nonpositive. For a unit eigenvector $w$ of $C_1$ for its
largest eigenvalue $c'$,
$\inner w{\dot C_1w}=\sum_j\inner{H_j^\dagger w}{\phi_0H_j^\dagger w}-\bar\phi_1c'
\leq(\bar\phi_0-\bar\phi_1)c'$; so whenever $\bar\phi_1\geq\bar\phi_0$, the
upper right derivative of $\bar\phi_1$ is nonpositive. Hence
$\max(\bar\phi_0,\bar\phi_1)$ is nonincreasing. The statement for the
smallest phase is the same argument with $\lambda_{\max}(C_0)$ and
$\lambda_{\min}(C_1)$.

(iii) At a fixed point modulo rescaling the velocities are a common scalar,
$\phi_0=c\,\id$ and $\phi_1=c\,\id$ with $c\in(0,\pi)$, and
$\tr\Omega=Z$ forces $c=\varphi$.
\end{proof}

\subsection{The mass}
The example of \cite[Theorem~1.1]{wue26} with increasing mass has charges
close to the real axis and an invertible differential, and the question
whether the mass decreases is asked in \cite[Question~5.18]{hkkp26}. The
mass can also increase for the symmetric charges and for complexes with
stable cohomology and $r=1$, to which Theorems~\ref{thm:monotone}
and~\ref{thm:convergence} apply.

\begin{proposition}[A state with increasing mass]\label{prop:mass}
Let $z_0=-1+\ii$, $z_1=1+\ii$, $k=512$, $V=\CC^2$, $U=\CC$, $W=\CC^{512}$ with
the standard metrics, $D=(0,f)$ with $f^2=\frac{383}{100}$, and let $A_j$ be
the matrix whose only nonzero row is the $j$-th, equal to
$(-0.1978,\,-0.6072)$ for $480$ values of $j$ and to $(-0.5038,\,0.8392)$ for
the other $32$. The cohomology representation $\CC\to\CC^{512}$ is stable,
$r=1$, and at this state
\[
 m=749.87188\ldots,\qquad\frac{dm}{dt}=1.77041\ldots>0,\qquad\cD=48.7558\ldots>0 .
\]
\end{proposition}

\begin{proof}
Since $r=1$ and the arrows have orthogonal images, $\Omega$ is in
closed form: with $v=D^\dagger$, $\ell=|v|^2$, $a_j$ the nonzero row of $A_j$
and $C=\sum_ja_j^\dagger a_j$,
\[
 \Omega_V=z_0\bigl(\id-R(\ell)vv^\dagger\bigr)+C-\tfrac{1-R(\ell)}\ell(vv^\dagger C+Cvv^\dagger)
 +\tfrac{1-2R(\ell)+e^{-\ell}}{\ell^2}(v^\dagger Cv)\,vv^\dagger,
\]
$\Omega_U=(v^\dagger Cv)(1-e^{-\ell}-\ell e^{-\ell})/\ell^2-z_0e^{-\ell}$, and
$\Omega_W$ is diagonal with entries
$z_1-|a_j|^2+\frac{1-R(\ell)}\ell|a_jv|^2$. The flow preserves this form, with
$\dot v=\frac12(\Phi_V-\Phi_U)v$ and
$\dot a_j=\frac12(a_j\Phi_V-(\Phi_W)_{jj}a_j)$, and
$\frac d{dt}\tr|T|=\Rea\tr(\dot T\,e^{-\ii\Psi})$ for a matrix $T$ with polar
phase $\Psi$. The $2\times2$ polar decomposition is explicit:
$|T|=(M+\sqrt{\det M})/\sqrt{\tr M+2\sqrt{\det M}}$ with $M=TT^\dagger$. The
three numbers were enclosed with interval arithmetic at $500$ bits; the
enclosures have radius below $10^{-20}$.\footnote{The program
\texttt{prop\_9\_2\_mass.py} and its output accompany this paper as
supplementary files. The program also checks
the trace identity \eqref{eq:central}, compares the derivative with a
difference quotient of the enclosed mass along the velocity, and gives
$\dot m<0$ for the same state with all rows divided by $20$.} Stability: every row $a_j$ has a nonzero
first entry, so the proper nonzero subrepresentations are subspaces of $W$,
of phase $\frac\pi4<\varphi$.
\end{proof}

With these two rows in the proportion $15:1$, the same formulas give
$\dot m>0$ for $k=400$ arrows and $\dot m<0$ for $k=368$; both signs are
enclosed by the same program. For $k\leq3$ arrows
and the symmetric charges, a numerical maximization of $\dot m/m$ over the
maps $D,A_j$ from random starting points, for small dimensions, found no
metric with increasing mass; this observation is not used in any
proof.

\subsection{The extreme phases}

\begin{proposition}[Extreme phases]\label{prop:extreme}
Let $z_0=-1+\ii$, $z_1=1+\ii$, and take the standard metrics.
\begin{enumerate}
\item[(i)] For $V=\CC^2$, $U=W=\CC$, $D=(0,f)$ with $f^2=\frac1{10}$,
 $A_1=(1,0)$ and $A_2=(0,10)$, the smallest phase satisfies
 $\frac d{dt}\phi_-=-0.38959518\ldots<0$.
\item[(ii)] For $V=\CC^3$, $U=\CC$, $W=\CC^2$, $D=(0,0,\frac92)$,
 \[
 A_1=\begin{pmatrix}5.7&-1.7&6\\6&0&-5.6\end{pmatrix},\qquad
 A_2=\begin{pmatrix}6&0.1&-6\\-6&-1.8&-6\end{pmatrix},
 \]
 the largest eigenvalue of $\Phi_V$ is simple and exceeds the eigenvalues of
 $\Phi_W$, and $\frac d{dt}\phi_+=0.0060800323\ldots>0$.
\end{enumerate}
In both cases $r=1$ and the cohomology representation is polystable; in (i)
it is stable.
\end{proposition}

\begin{proof}
(i) Here $C=\diag(1,100)$, all blocks of $\Omega$ are numbers, and the
flow preserves this form. With $\ell=f^2$, $x=100/\ell$ and
$\mathsf m(\ell)=1-(1+\ell)e^{-\ell}$, the phases are $\arg(z_0+1)$ and
$\Phi_A:=\arg(z_0+100)$ on the two lines of $V$,
$\Phi_W=\arg(z_1-1-x(1-e^{-\ell}))$, and
$\Phi_U=\arg(z_0e^{-\ell}-x\,\mathsf m(\ell))-\pi=-\arctan\frac1{1+y}$ with
$y=x\,\mathsf m(\ell)e^\ell$; moreover $\dot\ell=\ell(\Phi_A-\Phi_U)$ and
$\dot x=x(\Phi_U-\Phi_W)$ by \eqref{eq:mapflow} and \eqref{eq:adapted}. Hence
$\dot\Phi_U=\dot y/((1+y)^2+1)$ and
$\dot y/y=(\Phi_U-\Phi_W)+(\ell e^{-\ell}/\mathsf m(\ell)+1)\dot\ell$, an
explicit expression, which was enclosed with interval
arithmetic.\footnote{The program \texttt{prop\_9\_3\_extreme\_phases.py} and
its output accompany this paper as supplementary files.} The cohomology
representation is $1:\CC\to\CC$ with a second arrow equal to zero.

(ii) Since $r=1$, $\Omega$ is given by the formulas in the proof of
Proposition~\ref{prop:mass}, with $C=A_1^\dagger A_1+A_2^\dagger A_2$ and
$\Omega_W=z_1-\sum_jA_j\bigl(\id-\tfrac{1-R(\ell)}\ell vv^\dagger\bigr)A_j^\dagger$;
the derivative of $\Omega_V$ along the flow is obtained by the product rule
from $\dot v=\frac12(\Phi_V-\Phi_U)v$ and
$\dot A_j=\frac12(A_j\Phi_V-\Phi_WA_j)$. By Lemma~\ref{lem:variation}, the
derivative of a simple eigenvalue of $\Phi_V$ with spectral projection $\Pi$
is $2\tr(\Pi\Theta)$, where $\Theta$ solves \eqref{eq:sylvester} for the
variation $\dot\Omega_V$. The spectra of $\Phi_V$ and $\Phi_W$, the
projection $\Pi$ and the solution $\Theta$ were enclosed with interval
arithmetic at $400$ bits; the spectra are
$\{0.0071\ldots,\,0.1929\ldots,\,3.13307\ldots\}$ and
$\{3.12792\ldots,\,3.12883\ldots\}$. The eigenvalue $\phi_+$ carries the
weight $6.09\ldots\times10^{-6}$ of the mass measure, whose total mass is
$303.13\ldots$ The cohomology representation $\CC^2\to\CC^2$ has the arrows
$H_1,H_2$ given by the first two columns of $A_1,A_2$. The polynomial
$\det(H_1-\tau H_2)=-\frac{51}5\tau^2+\frac{1053}{50}\tau+\frac{51}5$ has two
distinct roots $\tau_1,\tau_2$; if $H_1v_i=\tau_iH_2v_i$, then the lines
$\CC v_i\to\CC H_2v_i$ are two non-isomorphic stable subrepresentations of
dimension vector $(1,1)$, and the representation is their direct sum.
\end{proof}

\subsection{Fixed points and harmonic metrics}

\begin{proposition}\label{prop:fixed}
Let $r\geq1$. For all charges in the upper half-plane the phase flow has no
fixed point and no periodic orbit on the metrics of the complex modulo
rescaling, and the complex carries no harmonic metric. If
$\alpha_0\neq\alpha_1$, the phase potential has no critical point.
\end{proposition}

\begin{proof}
The least eigenvalue $\lambda$ of $DD^\dagger$ is invariant under a common
rescaling of the metrics, and by Proposition~\ref{prop:imported}(i) its right
derivative is at least
$\lambda\min\bigl(\inner v{\Phi_{AA}v}-\inner u{\Phi_Uu}\bigr)>0$ over unit
vectors, since $\Phi_V>0>\Phi_U$. So $\lambda$ is strictly increasing along
every trajectory. A harmonic metric would have $\Phi_V$ and $\Phi_U$ equal to
the same scalar, which is impossible for the same reason. The differential
of the phase potential is
$\xi\mapsto\sigma^{-1}\sum_\bullet\tr(\Ima(e^{-\ii\varphi}\Omega_\bullet)\xi_\bullet)$
(Section~\ref{sec:potential}), so at a critical point
$\Ima(e^{-\ii\varphi}\Omega_\bullet)=0$ for all three blocks. But
$\tr\Ima(e^{-\ii\varphi}\Omega_U)=-\sigma\tr J-c_0\tr e^{-L}<0$ if $c_0>0$, and
$\tr\Ima(e^{-\ii\varphi}\Omega_V)=c_0\tr e^{-B}-\sigma\tr K<0$ if $c_0<0$; here
$\tr J\geq0$ by Theorem~\ref{thm:kernel}, $\tr K=\tr(Ce^{-B})\geq0$, and
$c_0\neq0$ exactly when $\alpha_0\neq\alpha_1$.
\end{proof}

So on the metrics of the given complex the set of fixed points and the set
of harmonic metrics coincide trivially, both being empty. In the sense of
\cite[Remark~5.15]{hkkp26} the harmonic metrics of the object live on other
representatives; for ordered charges and polystable cohomology
Theorem~\ref{thm:convergence} identifies the limit of the flow with a
harmonic metric of the cohomology representation.

\subsection{The mass measure}
The extreme phases need not converge to $\varphi$. If the arrows vanish on
$(\ker D)^\perp$, then $\Omega_V=(z_0+C_0)\oplus z_0e^{-F^\dagger F}$ and
$\Omega_U=-z_0e^{-L}$, the flow preserves this form, and the phases of the
acyclic part are $\alpha_0$ and $\alpha_0-\pi$ for all time, while their
weights $|z_0|e^{-\ell_i}$ in the mass measure tend to zero.

\begin{proposition}\label{prop:measure}
For every metric on the complex,
\[
 m(h)-|Z|=\int\bigl(1-\cos(\phi-\varphi)\bigr)\,d\nu_h(\phi)\geq0 .
\]
Along a trajectory, $\nu_{h(t)}\to|Z|\,\delta_\varphi$ weakly if and only if
$m(h(t))\to|Z|$. In particular $\nu_{h(t)}\to|Z|\,\delta_\varphi$ for ordered
charges and semistable cohomology.
\end{proposition}

\begin{proof}
For a unit eigenvector $v_k$ of $\Phi$,
$\inner{v_k}{\Omega v_k}=p_ke^{\ii\phi_k}$, because
$\Omega=|\Omega|e^{\ii\Phi}$. Hence $\int e^{\ii\phi}d\nu_h=\tr\Omega=Z$ and
$\int d\nu_h=m(h)$; take the real part of $e^{-\ii\varphi}$ times the first
identity. The phases lie in $(-\pi,\pi)$ and $\varphi\in(0,\pi)$, so on
$[-\pi,\pi]$ the function $1-\cos(\phi-\varphi)$ vanishes only at
$\phi=\varphi$. If $m\to|Z|$, the $\nu_{h(t)}$-measure of the complement of
any neighbourhood of $\varphi$ therefore tends to zero, and the total mass
tends to $|Z|$; the converse follows by integrating the constant $1$. The
last statement is Theorem~\ref{thm:relaxation}.
\end{proof}

\section{Charges for which the potential is not monotone, and questions}\label{sec:remarks}

The phase potential is monotone for the charges of
Theorem~\ref{thm:monotone}. This section gives charges for which it is not,
and states the questions left open.

\subsection{Reversed and ordered charges}
\begin{proposition}\label{prop:reversed}
If all $A_j=0$, then for all charges in the upper half-plane
\[
 \cD=c_0\bigl[\,n_0(\alpha_0-\alpha_1)+\pi\tr e^{-L}\,\bigr].
\]
In particular, for reversed charges one has $\cD<0$ at every such state with
$\pi\tr e^{-L}>n_0(\alpha_1-\alpha_0)$; such states exist if and only if
$\pi r>n_0(\alpha_1-\alpha_0)$.
\end{proposition}

\begin{proof}
The matrices $\Omega_V=z_0e^{-B}$, $-\Omega_U=z_0e^{-L}$, $\Omega_W=z_1\id$
are normal with phases $\alpha_0$, $\alpha_0-\pi$, $\alpha_1$. Hence
$\cD_V=|z_0|(n_0+\tr e^{-L})(\alpha_0-\varphi)\sin(\alpha_0-\varphi)$,
$\cD_U=|z_0|\tr e^{-L}(\pi-\alpha_0+\varphi)\sin(\alpha_0-\varphi)$ and
$\cD_W=n_1c_1(\varphi-\alpha_1)=n_0c_0(\varphi-\alpha_1)$. For reversed
charges $c_0<0$, and $\tr e^{-L}$ takes all values in $(0,r)$.
\end{proof}

For example, for $z_0=1+2\ii$, $z_1=-3+2\ii$, $n_0=n_1=r=1$, $D=(0,\frac12)$
and standard metrics, $\cD=-1.78927\ldots$

For ordered charges the potential can also fail to be monotone when
$r\geq2$.

\begin{proposition}\label{prop:orderedfailure}
Let $z_0=10^{-6}(-1+\ii)$ and $z_1=1+\frac\ii{500}$, so that
$\alpha_0=\frac{3\pi}4>\alpha_1$, and let $n_0=n_1=1$, $r=2$, $k=2$. With
standard metrics on $V=\CC^3$, $U=\CC^2$, $W=\CC$, let
\[
 D=\begin{pmatrix}0&\sqrt{\ell_1}&0\\0&0&\sqrt{\ell_2}\end{pmatrix},\quad
 A_1=(\tfrac15,0,0),\quad A_2=(0,a_1,a_2),\qquad\ell_1=\tfrac1{20},\ \ell_2=8,
\]
\[
 a_1=\sqrt{\epsilon/\ell_1},\quad a_2=\tfrac7{10}\sqrt{\epsilon\,\ell_2},\quad
 \epsilon=\tfrac3{10}\Bigl(\frac{1-e^{-\ell_1}}{\ell_1^2}+\frac{49}{100}(1-e^{-\ell_2})\Bigr)^{-1}.
\]
The cohomology representation is the nonzero map $\frac15:\CC\to\CC$, which
is stable, and
\[
 \cD=-1.6654354\ldots\times10^{-6}<0,\qquad
 (\cD_V,\cD_U,\cD_W)\approx(6.72,\,-9.09,\,0.70)\times10^{-6}.
\]
\end{proposition}

\begin{proof}
Since $C=\diag(\frac1{25},\,(a_1,a_2)^\dagger(a_1,a_2))$ has no block between
$\ker D$ and its complement, all blocks of $\Omega$ are of size $1$ or
$2$, and the polar decomposition of a $2\times2$ matrix is explicit, as in
the proof of Proposition~\ref{prop:mass}. The value of $\cD$ was enclosed
with interval arithmetic at $300$ bits; the enclosure has radius below
$10^{-30}$.\footnote{The program
\texttt{prop\_10\_1\_10\_2\_dissipation.py} and its output accompany this
paper as supplementary files.} At the same state and for the symmetric charges the same
computation gives $\cD=3.79\ldots>0$, in agreement with
Theorem~\ref{thm:monotone}.
\end{proof}

The mechanism is the following. As $z_0\to0$ with $z_1=e^{\ii\epsilon'}$ and
no arrows on the cohomology, $\Omega_V$ and $\Omega_U$ tend to the Hermitian
matrices $K_{F^\dagger F}(C_T)$ and $J$, whose phases have the eigenvalues
$0,\pi$ and $0,-\pi$ respectively, and
$\cD/\sigma\to\epsilon'\,\tr Y/(1-\tr Y)+\pi\bigl(\tr K_{F^\dagger F}(C_T)_--\tr J_-\bigr)$
up to terms of order $\epsilon'^2$ ($n_1=1$, $\tr Y<1$). Configurations with
$\tr J_->\tr K_{F^\dagger F}(C_T)_-$ exist, and then $\cD<0$ for small
$\epsilon'$. By Remark~\ref{rem:sharp} the block of degree one alone is
negative on a much larger set of charges; in numerical minimizations of
$\cD$ over the maps $D,A_j$, which are not used in any proof, the blocks of
degree zero compensated it except for charges very close to the real axis. By
Theorems~\ref{thm:cancellation} and~\ref{thm:convergence} the failure of
monotonicity of the potential does not affect the long-time behaviour.

\subsection{Questions}
We collect the questions left open.
\begin{enumerate}
\item Determine the best constant in Theorem~\ref{thm:kernel}. The value
 $(\sqrt2-1)/2$ would enlarge the admissible range to
 $\varphi\geq\frac\pi2(\sqrt2-1)$ and improve the constant $\rho$ of
 Theorem~\ref{thm:energy}.
\item Determine the set of charges for which $\cD\geq0$ at every state. It
 contains the admissible charges, and by
 Proposition~\ref{prop:orderedfailure} it does not contain all ordered
 charges when $r\geq2$.
\item For semistable cohomology that is not polystable, does the bounded
 factor between the induced metrics and the asymptotic solutions of
 \cite{hkkp23} converge, as in Proposition~\ref{prop:semistable}
 (Remark~\ref{rem:semistable}; \cite[Corollary~6.5]{wue26b})?
\item Is the mass nonincreasing along the flow for the symmetric charges and
 few arrows?
\item Do the parameter derivatives of $\cK_t$ converge pointwise on
 $\mathcal P^{\mathrm s}$ (Theorem~\ref{thm:families})?
\item Extend Theorems~\ref{thm:cancellation} and~\ref{thm:convergence} to
 complexes with acyclic parts at both vertices or in several degrees. The
 block of degree one then involves a Schur multiplier given by divided
 differences of the exponential function with one more interpolation point,
 to which Theorem~\ref{thm:kernel} does not apply.
\end{enumerate}

\begin{thebibliography}{WCKQ20}

\bibitem[Aud11]{aud11}
K.~M.~R.~Audenaert,
\textit{A characterisation of anti-L\"owner functions},
Proc. Amer. Math. Soc. \textbf{139} (2011), no.~12, 4217--4223.
\href{https://arxiv.org/abs/1008.2943v2}{arXiv:1008.2943v2}.

\bibitem[Bha97]{bha97}
R.~Bhatia,
\textit{Matrix analysis},
Graduate Texts in Mathematics \textbf{169}, Springer, New York, 1997.

\bibitem[BZ92]{bz92}
J.-M.~Bismut and W.~Zhang,
\textit{An extension of a theorem by Cheeger and M\"uller},
Ast\'erisque \textbf{205} (1992), 235~pp.

\bibitem[Bri07]{bri07}
T.~Bridgeland,
\textit{Stability conditions on triangulated categories},
Ann. of Math. (2) \textbf{166} (2007), no.~2, 317--345.

\bibitem[BG04]{bg04}
D.~Burns and V.~Guillemin,
\textit{Potential functions and actions of tori on K\"ahler manifolds},
Comm. Anal. Geom. \textbf{12} (2004), no.~1, 281--303.
\href{https://arxiv.org/abs/math/0302350}{arXiv:math/0302350}.

\bibitem[Don85]{don85}
S.~K.~Donaldson,
\textit{Anti self-dual Yang--Mills connections over complex algebraic
surfaces and stable vector bundles},
Proc. London Math. Soc. (3) \textbf{50} (1985), no.~1, 1--26.

\bibitem[GRS21]{grs21}
V.~Georgoulas, J.~W.~Robbin and D.~A.~Salamon,
\textit{The moment-weight inequality and the Hilbert--Mumford criterion},
Lecture Notes in Math. \textbf{2297}, Springer, Cham, 2021.
\href{https://arxiv.org/abs/1311.0410v2}{arXiv:1311.0410v2}.

\bibitem[HKKP23]{hkkp23}
F.~Haiden, L.~Katzarkov, M.~Kontsevich and P.~Pandit,
\textit{Semistability, modular lattices, and iterated logarithms},
J. Differential Geom. \textbf{123} (2023), no.~1, 21--66.
\href{https://doi.org/10.4310/jdg/1679503804}{doi:10.4310/jdg/1679503804};
\href{https://arxiv.org/abs/1706.01073v2}{arXiv:1706.01073v2}.

\bibitem[HKKP26]{hkkp26}
F.~Haiden, L.~Katzarkov, M.~Kontsevich and P.~Pandit,
\textit{Towards categorical K\"ahler geometry},
preprint (2026), \href{https://arxiv.org/abs/2609.00978v1}{arXiv:2609.00978v1}.

\bibitem[HJ13]{hj13}
R.~A.~Horn and C.~R.~Johnson,
\textit{Matrix analysis}, 2nd ed.,
Cambridge Univ. Press, Cambridge, 2013.

\bibitem[Joh17]{joh17}
F.~Johansson,
\textit{Arb: efficient arbitrary-precision midpoint-radius interval
arithmetic},
IEEE Trans. Comput. \textbf{66} (2017), no.~8, 1281--1292.
\href{https://doi.org/10.1109/TC.2017.2690633}{doi:10.1109/TC.2017.2690633}.

\bibitem[KN79]{kn79}
G.~Kempf and L.~Ness,
\textit{The length of vectors in representation spaces},
in \textit{Algebraic geometry (Copenhagen, 1978)},
Lecture Notes in Math. \textbf{732}, Springer, Berlin, 1979, 233--243.

\bibitem[Kin94]{kin94}
A.~D.~King,
\textit{Moduli of representations of finite dimensional algebras},
Quart. J. Math. Oxford Ser. (2) \textbf{45} (1994), no.~4, 515--530.
\href{https://doi.org/10.1093/qmath/45.4.515}{doi:10.1093/qmath/45.4.515}.

\bibitem[Min62]{min62}
G.~J.~Minty,
\textit{Monotone (nonlinear) operators in Hilbert space},
Duke Math. J. \textbf{29} (1962), no.~3, 341--346.
\href{https://doi.org/10.1215/S0012-7094-62-02933-2}{doi:10.1215/S0012-7094-62-02933-2}.

\bibitem[UY86]{uy86}
K.~Uhlenbeck and S.-T.~Yau,
\textit{On the existence of Hermitian--Yang--Mills connections in stable
vector bundles},
Comm. Pure Appl. Math. \textbf{39} (1986), no.~S1, S257--S293.

\bibitem[WCKQ20]{wckq20}
D.~Wang, W.~Chen, S.~Z.~Khong and L.~Qiu,
\textit{On the phases of a complex matrix},
Linear Algebra Appl. \textbf{593} (2020), 152--179.
\href{https://arxiv.org/abs/1904.07211v2}{arXiv:1904.07211v2}.

\bibitem[Wue26a]{wue26}
B.~J.~Wuebben,
\textit{Mass increase and acyclic cancellation for phase flows on quiver
complexes},
preprint (2026),
\href{https://doi.org/10.5281/zenodo.23921919}{doi:10.5281/zenodo.23921919}.

\bibitem[Wue26b]{wue26b}
B.~J.~Wuebben,
\textit{Convergence of the minimizing flow on acyclic quivers},
preprint (2026), \href{https://doi.org/10.5281/zenodo.24346923}{doi:10.5281/zenodo.24346923}.

\end{thebibliography}

\end{document}
